\documentclass[11pt, a4paper]{report}
%----------------------------------------------------------------------------------------
%	PACKAGES (宏包管理)
%----------------------------------------------------------------------------------------
% 1. 编码与语言
\usepackage[english]{babel}
\usepackage{fontspec}
\usepackage{xeCJK}
\setmainfont{Latin Modern Roman}
\setCJKmainfont{Microsoft YaHei}[AutoFakeSlant=0.2, AutoFakeBold=2]
% 2. 页面布局与排版
\usepackage{geometry}
\geometry{
    textwidth=6in, % 稍微加宽以便放入更好的box
    top=1in,
    bottom=1in,
    headheight=12pt,
    headsep=25pt,
    footskip=30pt
}
\usepackage{setspace}
\setstretch{1.2}
\usepackage{float}
\usepackage[shortlabels]{enumitem} % 更好的列表控制
\usepackage{tocloft}  % 加载目录控制宏包
\renewcommand{\cftsecleader}{\cftdotfill{\cftdotsep}} % 强制让 section 后面显示点号
% 3. 数学支持
\usepackage{amsmath, amsthm, amssymb, amsfonts, amscd}
\usepackage{mathtools}
\usepackage{thmtools} % 定理环境增强
\usepackage{extarrows}
% 4. 图形与颜色
\usepackage{graphicx}
\usepackage[dvipsnames, svgnames]{xcolor} % 加载更多颜色名
\usepackage{tikz}
\usepackage{tikz-cd}
\usepackage{ytableau} % 杨表
\usetikzlibrary{decorations.pathreplacing, calc, shadows}
% 5. 漂亮的文本框
\usepackage{tcolorbox}
\tcbuselibrary{breakable, skins, theorems}
% 6. 参考文献
%\usepackage[backend=biber,style=gb7714-2015,sorting=gb7714-2015]{biblatex}
%\addbibresource{references.bib}
% 7. 交互 (最后加载)
\usepackage[hidelinks]{hyperref}
\usepackage{array}
%----------------------------------------------------------------------------------------
%	CUSTOM COLORS (自定义配色方案)
%----------------------------------------------------------------------------------------
% 定理色系 - 蓝色
\definecolor{ThmColor}{RGB}{0, 71, 171}        % 深钴蓝
\definecolor{ThmBg}{RGB}{240, 244, 255}        % 极浅蓝
% 定义色系 - 绿色
\definecolor{DefColor}{RGB}{0, 100, 0}         % 深绿
\definecolor{DefBg}{RGB}{240, 255, 240}        % 蜜瓜绿
% 引理/命题色系 - 橙色
\definecolor{LemColor}{RGB}{204, 85, 0}        % 焦橙色
\definecolor{LemBg}{RGB}{255, 248, 240}        % 极浅橙
% 示例/注记色系 - 灰色/紫色
\definecolor{ExColor}{RGB}{105, 105, 105}      % 暗灰
\definecolor{ExBg}{RGB}{250, 250, 250}         % 极浅灰
% 专题/话题色系 - 靛紫色
\definecolor{TopicColor}{RGB}{75, 0, 130}      % 靛蓝色 (Indigo)
\definecolor{TopicBg}{RGB}{248, 245, 255}      % 极浅紫
%----------------------------------------------------------------------------------------
%	THEOREM ENVIRONMENTS (环境美化核心)
%----------------------------------------------------------------------------------------
% 1. Theorem (蓝色风格)
\declaretheorem[name=Theorem, numberwithin=section]{theorem}
\tcolorboxenvironment{theorem}{
    enhanced jigsaw,
    colframe=ThmColor,       % 边框颜色
    colback=ThmBg,           % 背景颜色
    coltitle=white,          
    boxrule=0.5pt,           % 边框粗细
    leftrule=2pt,            % 左边框加粗
    sharp corners=downhill,  % 样式
    breakable,               % 允许跨页
    drop shadow=black!5!white % 轻微阴影
}
% Corollary 使用与 Theorem 相同的风格，但通常紧随其后
\declaretheorem[name=Corollary, numberlike=theorem]{corollary}
\tcolorboxenvironment{corollary}{
    enhanced jigsaw,
    colframe=ThmColor!80!white, 
    colback=ThmBg,
    boxrule=0pt, leftrule=2pt, arc=0pt, outer arc=0pt,
    breakable
}
% 2. Definition (绿色风格)
\declaretheorem[name=Definition, numberlike=theorem]{definition}
\tcolorboxenvironment{definition}{
    enhanced jigsaw,
    colframe=DefColor,
    colback=DefBg,
    boxrule=0.5pt, leftrule=2pt,
    sharp corners=downhill,
    breakable,
    drop shadow=black!5!white
}
% 3. Lemma & Proposition (橙色风格)
\declaretheorem[name=Lemma, numberlike=theorem]{lemma}
\declaretheorem[name=Proposition, numberlike=theorem]{proposition}
\tcolorboxenvironment{lemma}{
    enhanced jigsaw,
    colframe=LemColor,
    colback=LemBg,
    boxrule=0.5pt, leftrule=2pt,
    breakable
}
\tcolorboxenvironment{proposition}{
    enhanced jigsaw,
    colframe=LemColor,
    colback=LemBg,
    boxrule=0.5pt, leftrule=2pt,
    breakable
}
% 4. Example & Remark (简约风格)
\declaretheorem[name=Example, numberlike=theorem]{example}
\declaretheorem[name=Remark, numberlike=theorem]{remark}
\declaretheorem[name=Exercise, numberlike=theorem]{exercise}
\tcolorboxenvironment{example}{
    enhanced jigsaw,
    colframe=ExColor,
    colback=ExBg,
    leftrule=2pt, boxrule=0pt,
    frame hidden, % 隐藏其他三边
    breakable
}
% 5. Topic (靛紫色风格，带标题支持)
\declaretheorem[name=Topic, numberlike=theorem]{topic}
\tcolorboxenvironment{topic}{
    enhanced jigsaw,
    breakable,
    colframe=TopicColor,        % 边框为靛紫色
    colback=TopicBg,           % 背景为极浅紫
    boxrule=0.5pt,             % 细边框
    leftrule=2pt,              % 左侧加粗线
    sharp corners=downhill,    % 与你其他环境一致的切角样式
    drop shadow=black!5!white, % 轻微阴影
    before skip=15pt,
    after skip=15pt
}
%----------------------------------------------------------------------------------------
%	MATH MACROS (保留你的常用定义)
%----------------------------------------------------------------------------------------
\newcommand{\g}{\mathfrak{g}}  
\newcommand{\der}{\mathrm{d}} 
\newcommand{\h}{\mathfrak{h}} 
\newcommand{\n}{\mathfrak{n}}  
\newcommand{\U}{\mathcal{U}}  
\newcommand{\C}{\mathbb{C}} 
\newcommand{\R}{\mathbb{R}} 
\newcommand{\borel}{\mathfrak{b}}
\newcommand{\Ind}{\mathrm{Ind}} 
\newcommand{\Hom}{\mathrm{Hom}}  
\newcommand{\Ext}{\mathrm{Ext}}
\newcommand{\HRule}[1]{\rule{\linewidth}{#1}}
%----------------------------------------------------------------------------------------
%	TITLE PAGE
%----------------------------------------------------------------------------------------
\begin{document}
\begin{titlepage}
    \centering
    \vspace*{2cm}
    
    \textsc{\Large Course Notes}\\[0.5cm]

    \HRule{1.5pt} \\[0.5cm]
    { \huge \bfseries Abstract Algebra III }\\[0.4cm] 
    \HRule{1.5pt} \\[1.5cm]
    
    \vspace{3cm}

    % 人员信息区域：垂直排列并区分字号
    \begin{spacing}{1.8}
        {\Large \textbf{Instructor:}} \\
        {\Large Prof. Zhicheng \textsc{Feng}} \\[0.8cm]
        
        {\normalsize \textit{Student:}} \\
        {\normalsize Leo}
    \end{spacing}

    \vfill
    {\large \today}
\end{titlepage}
\newpage
\tableofcontents
\newpage


%----------------------------------------------------------------------------------------
%    CONTENT STRUCTURE TEMPLATE
%----------------------------------------------------------------------------------------
\chapter{Lecture 1: Modules}

Message: QQ, BB. HW 30\%, Mid 20\%, Final 50\%

\section{Basic Definitions and Examples}

% --- Begin rewritten content from 抽代3.md ---
\begin{definition}
Let $R$ be a ring with $1$, and let $M$ be an abelian group. We say that $M$ is a left $R$-module if $R$ ``left acts'' on $M$:
\[
R \times M \rightarrow M, \quad (r, m) \mapsto rm
\]
such that
\[
\begin{aligned}
r(m+n) &= rm + rn, \\
(r+s)m &= rm + sm, \\
(rs)m &= r(sm), \\
1m &= m.
\end{aligned}
\quad \text{for any } r, s \in R, m,n \in M
\]

A \textbf{right $R$-module} is defined by $(r, m) \mapsto mr$ with the axiom $m(rs) = (mr)s$. If $R$ is commutative, left-modules and right-modules coincide. Otherwise, one can transform using the opposite ring $R^{\mathrm{op}}$.
\end{definition}

\begin{theorem}
Let $M$ be an abelian group. Then $M$ is an $R$-module if and only if there is a ring homomorphism $R \rightarrow \operatorname{End}(M)$.
\end{theorem}

\begin{proof}
$(\Rightarrow)$ Suppose $M$ is an $R$-module. Define $\rho: R \rightarrow \operatorname{End}(M)$ by
\[
\rho(r)(m) = rm, \quad \forall r \in R, m \in M.
\]
For each $r \in R$, the map $\rho(r): M \rightarrow M$ is a group homomorphism since $r(m_1 + m_2) = rm_1 + rm_2$. Moreover, $\rho$ is a ring homomorphism:
\begin{itemize}
    \item $\rho(r + s)(m) = (r + s)m = rm + sm = \rho(r)(m) + \rho(s)(m)$, so $\rho(r + s) = \rho(r) + \rho(s)$.
    \item $\rho(rs)(m) = (rs)m = r(sm) = \rho(r)(\rho(s)(m))$, so $\rho(rs) = \rho(r) \circ \rho(s)$.
    \item $\rho(1)(m) = 1m = m$, so $\rho(1) = \operatorname{id}_M$.
\end{itemize}

$(\Leftarrow)$ Suppose $\rho: R \rightarrow \operatorname{End}(M)$ is a ring homomorphism. Define the scalar multiplication $R \times M \rightarrow M$ by
\[
(r, m) \mapsto \rho(r)(m).
\]
We verify the module axioms:
\begin{itemize}
    \item $r(m_1 + m_2) = \rho(r)(m_1 + m_2) = \rho(r)(m_1) + \rho(r)(m_2) = rm_1 + rm_2$.
    \item $(r + s)m = \rho(r + s)(m) = (\rho(r) + \rho(s))(m) = \rho(r)(m) + \rho(s)(m) = rm + sm$.
    \item $(rs)m = \rho(rs)(m) = (\rho(r) \circ \rho(s))(m) = \rho(r)(\rho(s)(m)) = r(sm)$.
    \item $1m = \rho(1)(m) = \operatorname{id}_M(m) = m$.
\end{itemize}
Hence $M$ is an $R$-module.
\end{proof}

\begin{example}
(1) Let $M$ be an abelian group. Then $M$ is an $\operatorname{End}(M)$-module via
\[
\operatorname{End}(M) \times M \rightarrow M, \quad (\varphi, m) \mapsto \varphi(m).
\]
(Use the theorem: $\operatorname{End}(M) \rightarrow \operatorname{End}(M)$.)

(2) Let $V$ be a vector space over a field $K$, and let $A: V \rightarrow V$ be a linear transformation. Then $V$ is a $K[x]$-module via
\[
K[x] \times V \rightarrow V, \quad (f(x), v) \mapsto f(A)v.
\]

(3) If $|M| = 1$, then $M$ is called the zero module.

(4) Let $ R $ be a ringg with 1 and $ A \in M_n(R) $. The solution of $ AX=0 $ where $ X \in R^n $ is a righ t $ R $-module.   
\end{example}

\begin{definition}[$R$-submodule]
Let $R$ be a ring (with $1$) and let $M$ be a left $R$-module. A subset $N\subseteq M$
is called an \emph{$R$-submodule} of $M$ if:
\begin{enumerate}
  \item $0\in N$;
  \item $n_1,n_2\in N \implies n_1+n_2\in N$;
  \item $n\in N \implies -n\in N$;
  \item $r\in R,\ n\in N \implies rn\in N$.
\end{enumerate}
Equivalently, $N$ is a subgroup of the abelian group $(M,+)$ that is closed under the
$R$-action.
\end{definition}

\textbf{Question:} Are submodules of a finitely generated $R$-module finitely generated?

\textbf{Answer:} No, not in general. This property holds if and only if the ring $R$ is left Noetherian.

\begin{definition} 
Let $ M $ be an $ R $-module. A subgroup $ N $ of $ M $ is called a submodule of $ M $ if $ N $ is closed under scalar multiplication by elements of $ R $. In other words, $ rn \in N $ for all $ r \in R $ and $ n \in N $.
\end{definition}

\begin{example}[Submodule of a f.g. module that is not f.g.]
Let $k$ be a field and consider the polynomial ring $R = k[x_1, x_2, x_3, \dots]$ in infinitely many variables.
\begin{itemize}
    \item The ring $R$ itself can be viewed as a left $R$-module, which we denote by $M=R$. $M$ is finitely generated, for instance by the single element $1 \in R$.
    \item Consider the ideal $I = \langle x_1, x_2, x_3, \dots \rangle$ generated by all the variables. This ideal $I$ is a submodule of $M=R$.
    \item This submodule $I$ is not finitely generated. To see this, suppose for contradiction that $I$ is finitely generated by a set of polynomials $\{f_1, \dots, f_n\}$. Let $N$ be the largest index of any variable appearing in any of the polynomials $f_1, \dots, f_n$. Then any element in the ideal generated by $\{f_1, \dots, f_n\}$ must be a polynomial in variables $x_1, \dots, x_N$. However, the polynomial $x_{N+1}$ is in $I$ but cannot be expressed as a combination of $f_1, \dots, f_n$. This is a contradiction.
\end{itemize}
Therefore, the ideal $I$ is a submodule of the finitely generated $R$-module $R$ which is not itself finitely generated. This shows that the ring $R$ is not Noetherian.
\end{example}

\begin{definition} 
    Let $ M $ be an $ R $-module, $ M_i, i\in I $ are $ R $-submodules of $ M $. Then 
    \[
    \bigcap_{i\in I} M_i
    \]     
is an $ R $-submodule of $ M $.   
\[
\sum_{i\in I} = \{ \sum_{i \in I} \text{ finite} m_i | m_i \in M_i\} 
\]
is an $ R $-submodule of $ M $.
\end{definition}

\begin{definition} 
Let $M$ be an $R$-module, $ S \subseteq M $. 
\[
<S> = \{ \sum_{finite}^{a_is_i} | a_i \in R, s_i \in S \}
\]
or
\[
<S>=\bigcap_{N \text{is an} R -\text{submodule of }  M   \   S \subseteq N } N
\]
is an $ R $-submodule of $ M $, called the $ R $-submodule of $ M $ generated by $ S $.   

If $ M=<S> $, then we say $ M $ is generated by $ S $. In particular, if $ |S|=1 $ then we call $ M $ a cyclic module. If $ |S| < \infty $ then we call $ M $ a finite generated module.    
\end{definition}

\section{Quotient Modules and Homomorphism Theorems}

\begin{proposition}
Let $M$ be an $R$-module, $N$ a submodule of $M$, $M'$ an $R$-module, and $\varphi: M \rightarrow M'$ a morphism such that $N \subseteq \ker(\varphi)$. Then there exists a unique morphism
\[
\varphi_{N}: M/N \rightarrow M'
\]
such that the following diagram commutes. Moreover, $\ker(\varphi_{N}) = \ker\varphi / N$.
\[
\begin{tikzcd}
M \arrow[r, "\varphi"] \arrow[d, "\pi_{N}"'] & M' \\
M/N \arrow[ru, "\varphi_{N}"', dashed]
\end{tikzcd}
\]
\end{proposition}

\begin{remark}
The pair $(M/N, \pi_{N})$ is unique in the following sense: if $(S, \theta)$ is a pair such that
\begin{itemize}
\item $S$ is an $R$-module,
\item $\theta: M \rightarrow S$ is a surjective morphism satisfying the property of $(M/N, \pi_{N})$ described in the above proposition,
\end{itemize}
then there exists a unique isomorphism $\psi: M/N \rightarrow S$ such that the following diagram commutes.
\[
\begin{tikzcd}[row sep=2.8em, column sep=3.8em]
M \arrow[r, two heads, "\theta"] \arrow[d, two heads, "\pi_{N}"'] & S \\
M/N \arrow[ur, dashed, "\exists!\,\psi"'] &
\end{tikzcd}
\]
\end{remark}

\begin{proposition}
There is a unique isomorphism $\bar{\varphi}: \operatorname{coim}\varphi \rightarrow \operatorname{im}\varphi$ such that the following diagram commutes:
\begin{center}
\begin{tikzcd}
M \arrow[rr, "\varphi"] \arrow[dd, "\pi_{\operatorname{ker}(\varphi)}"'] &  & M'                                          \\
                                                                         &  &                                             \\
\operatorname{coim}(\varphi) \arrow[rr, "\overline{\varphi}"]            &  & \operatorname{im}(\varphi) \arrow[uu, hook]
\end{tikzcd}
\end{center}
Recall:
\[
\begin{aligned}
& \varphi: M \rightarrow M', \quad R\text{-morphism} \\
& \operatorname{coker}(\varphi)=M'/\operatorname{im}(\varphi) \\
& \operatorname{coim}(\varphi)=M/\ker \varphi
\end{aligned}
\]
\end{proposition}

\begin{definition}
Let $M$ be an $R$-module and $I$ an ideal of $R$. Denote by $IM$ the submodule of $M$ generated by all elements $am$ where $a \in I$, $m \in M$.
\[
IM = \left\{ \sum_{\text{finite}} a_{i} m_{i} \mid a_{i} \in I, \; m_{i} \in M \right\}
\]
\end{definition}

\begin{proposition}
Let $M$ be an $R$-module. Then $R \rightarrow \operatorname{End}(M)$ is a ring homomorphism.
\begin{enumerate}
\item $R \rightarrow \operatorname{End}(M)$ factors through the natural surjection $R \rightarrow R/I$ if and only if $IM = 0$:
\[
\begin{tikzcd}[row sep=2.8em, column sep=4.2em]
R \arrow[rr, "\rho"] \arrow[dr, two heads, "\pi_{I}"'] && \operatorname{End}(M) \\
& R/I \arrow[ur, dashed, "\exists\,\rho_{I}"'] &
\end{tikzcd}
\]
$\exists \rho_{I}: R/I \rightarrow \operatorname{End}(M)$ such that the diagram commutes.
\item The module $M/IM$ is naturally endowed with a structure of $R/I$-module.
\item Let $\varphi: M \rightarrow M'$ be an $R$-homomorphism. Then $\varphi(IM) \subseteq IM'$, so $\varphi$ induces a morphism of $R/I$-modules
\[
M/IM \rightarrow M'/IM'.
\]
\end{enumerate}
\end{proposition}

\section{Direct Sums and Free Modules}

\begin{definition}
Direct sum \& direct product $\left(\prod_{i \in I} M_i \; \& \; \bigoplus_{i \in I} M_i\right)$
\[
\left\{
\begin{array}{l}
\text{direct product: } \prod_{i \in I} M_i \\
\text{direct sum: } \bigoplus_{i \in I} M_i \text{ (finitely many } \neq 0\text{)}
\end{array}
\right.
\]
Equal when $|I| < \infty$. In a module category, direct product $=$ product, direct sum $=$ coproduct.
\end{definition}

\begin{definition}[free module]
Let $M$ be an $R$-module.
\begin{enumerate}
\item $x_{1}, \dots, x_{n} \in M$ are linearly independent if
\[
a_{1} x_{1} + a_{2} x_{2} + \cdots + a_{n} x_{n} = 0 \Rightarrow a_{1} = a_{2} = \cdots = a_{n} = 0.
\]
\item $M$ is called free if it is generated by a linearly independent subset. This subset is called a basis of $M$.
\item If $M$ is an $R$-module with a basis $x_{1}, \ldots, x_{n}$, then $M = \langle x_{1} \rangle \oplus \cdots \oplus \langle x_{n} \rangle$.
\end{enumerate}
\end{definition}

\begin{example}
$R \oplus \cdots \oplus R = R^{n}$.
\end{example}

\begin{theorem}
Let $R$ be a commutative ring with $1$. Let $M$ be a finitely generated free $R$-module. Then any basis of $M$ has the same cardinality. The number of elements in a basis is called the rank of the free module.
\end{theorem}

\begin{proof}
\textbf{Step 1 (Any basis is finite).}
Let $B=\{x_i\}_{i\in I}$ be a basis of $M$. Since $M$ is finitely generated, choose
generators $y_1,\dots,y_n$ of $M$. Each $y_k$ is a finite $R$--linear combination of elements
of $B$, so there exists a finite subset $I_0\subset I$ such that
$y_1,\dots,y_n\in \sum_{i\in I_0}Rx_i$. Hence
\[
M=\sum_{k=1}^n Ry_k \subseteq \sum_{i\in I_0}Rx_i \subseteq M,
\]
so $M=\sum_{i\in I_0}Rx_i$. If $I\setminus I_0\neq\varnothing$, pick $j\in I\setminus I_0$.
Then $x_j\in M=\sum_{i\in I_0}Rx_i$, contradicting that $B$ is linearly independent.

Thus $I=I_0$ is finite. 

\textbf{Step 2 (Reduce modulo a maximal ideal).}
Let $J\subset R$ be a maximal ideal and set $F:=R/J$, a field. Then $M/JM$ is naturally an $F$-vector space.

\textbf{Step 3 (A basis descends to a basis mod $J$).}
Let $x_1,\dots,x_r$ be an $R$--basis of $M$, and write $\overline{m}$ for the class of
$m\in M$ in $M/JM$, and $\overline{a}$ for the class of $a\in R$ in $F$.

\emph{Spanning.} For any $m\in M$, write $m=\sum_{i=1}^r a_i x_i$. Then
$\overline{m}=\sum_{i=1}^r \overline{a_i}\,\overline{x_i}$, so $\overline{x_1},\dots,\overline{x_r}$
span $M/JM$.

\emph{Linear independence.} Suppose $\sum_{i=1}^r \overline{a_i}\,\overline{x_i}=0$ in $M/JM$.
Then $\sum_{i=1}^r a_i x_i\in JM$, so there exist $a'_1,\dots,a'_r\in J$ such that
$\sum_{i=1}^r a_i x_i=\sum_{i=1}^r a'_i x_i$. Hence $\sum_{i=1}^r (a_i-a'_i)x_i=0$.
By linear independence of $x_1,\dots,x_r$, we get $a_i-a'_i=0$ for all $i$, so $a_i\in J$,
thus $\overline{a_i}=0$ for all $i$. Therefore $\overline{x_1},\dots,\overline{x_r}$ are
$F$-linearly independent, and form an $F$--basis of $M/JM$. 

\textbf{Step 4 (Compare cardinalities).}
Let $x_1,\dots,x_r$ and $y_1,\dots,y_s$ be two $R$--bases of $M$. By Step 3, their images
$\overline{x_1},\dots,\overline{x_r}$ and $\overline{y_1},\dots,\overline{y_s}$ are both $F$--bases
of the same $F$--vector space $M/JM$. Hence $r=\dim_F(M/JM)=s$. Therefore any two
bases of $M$ have the same cardinality.
\end{proof}

\chapter{Lecture 2: Modules and Categories}

\begin{proposition}[Properties of free modules]
Let $L$ be a free $R$-module.
\begin{enumerate}
\item Let $\pi: X \rightarrow Y$ be an epimorphism of $R$-modules. For any $R$-morphism $\varphi: L \rightarrow Y$, there exists a morphism $\tilde{\varphi}: L \rightarrow X$ such that the following diagram commutes (hence a free module is projective):
\[
\begin{tikzcd}
                               & L \arrow[d, "\varphi"] \arrow[ld, "\overline{\varphi}"', dashed] \\
X \arrow[r, "\pi"', two heads] & Y                                                               
\end{tikzcd}
\]
\item In particular, for any $R$-module $M$, any epimorphism $\pi: M \rightarrow L$ has a section, that is, there is a morphism $s: L \rightarrow M$ such that $\pi \circ s = \operatorname{id}_L$. In that case, $s$ is injective, $\operatorname{im}(s) \cong L$ and
\[
M = \ker \pi \oplus \operatorname{im}(s).
\]
\end{enumerate}
\end{proposition}

\begin{proof} 
\textbf{Step 1: Prove the fist property}

Let $B=\{e_i\}_{i\in I}$ be a basis of the free module $L$.

Since $\pi$ is surjective, for each $i\in I$ choose an element $x_i\in X$ such that
\[
\pi(x_i)=\varphi(e_i).
\]
(Here we use the axiom of choice if $I$ is infinite.)

By the defining universal property of a free module: given any $R$-module $X$ and any map
$f:B\to X$, there exists a unique $R$-linear morphism $F:L\to X$ such that $F(e_i)=f(e_i)$
for all $i$. Apply this with $f(e_i):=x_i$. Then there exists a unique $R$-linear map
\[
\widetilde\varphi:L\to X
\quad\text{such that}\quad
\widetilde\varphi(e_i)=x_i\ \ \text{for all }i\in I.
\]

Now check commutativity. Define two $R$-linear maps $L\to Y$:
\[
\psi_1:=\pi\circ \widetilde\varphi,
\qquad
\psi_2:=\varphi.
\]
For each basis element $e_i$ we have
\[
\psi_1(e_i)=(\pi\circ \widetilde\varphi)(e_i)=\pi(x_i)=\varphi(e_i)=\psi_2(e_i).
\]
Since $B$ generates $L$ and both $\psi_1,\psi_2$ are $R$-linear, it follows that
$\psi_1=\psi_2$, i.e.\ $\pi\circ \widetilde\varphi=\varphi$.

\textbf{Step2: Use the first property prove the second one}

Let $\pi:M\to L$ be an epimorphism. Apply (1) with $X=M$, $Y=L$, and take
$\varphi=\operatorname{id}_L:L\to L$. Since $L$ is free, by (1) there exists an $R$-morphism
\[
s:L\to M
\quad\text{such that}\quad
\pi\circ s=\operatorname{id}_L.
\]
So $\pi$ has a section.

\noindent\textit{(i) $s$ is injective.}
If $s(\ell)=0$, then applying $\pi$ gives
\[
\ell=(\pi\circ s)(\ell)=\pi(0)=0,
\]
hence $\ker(s)=0$ and $s$ is injective. In particular, $s$ induces an isomorphism
\[
L \xrightarrow{\ \sim\ } \operatorname{im}(s),\qquad \ell\mapsto s(\ell),
\]
so $\operatorname{im}(s)\cong L$.

\noindent\textit{(ii) $M=\ker\pi\oplus \operatorname{im}(s)$.}
First we show $\ker\pi\cap \operatorname{im}(s)=0$. If $m\in \ker\pi\cap \operatorname{im}(s)$, then
$m=s(\ell)$ for some $\ell\in L$ and also $\pi(m)=0$. Thus
\[
0=\pi(m)=\pi(s(\ell))=(\pi\circ s)(\ell)=\ell,
\]
so $m=s(0)=0$.

Next, we show $\ker\pi+\operatorname{im}(s)=M$. For any $m\in M$, set
\[
m_0 := m - s(\pi(m)).
\]
Then
\[
\pi(m_0)=\pi(m)-\pi(s(\pi(m)))=\pi(m)-(\pi\circ s)(\pi(m))=\pi(m)-\pi(m)=0,
\]
so $m_0\in \ker\pi$, and clearly $s(\pi(m))\in \operatorname{im}(s)$. Hence
\[
m = m_0 + s(\pi(m)) \in \ker\pi + \operatorname{im}(s).
\]
Combining with $\ker\pi\cap \operatorname{im}(s)=0$, we conclude
\[
M=\ker\pi\oplus \operatorname{im}(s).
\]
\end{proof}

\begin{definition}[Direct summand.]
Let $M$ be a module. If $M = M_1 \oplus M_2$, then $M_1$ is called a direct summand of $M$.
\end{definition}

\begin{proposition}
A direct summand of a free module is a projective module.
\end{proposition}

\begin{proof}
Let $F$ be a free $R$-module and suppose $F \cong P\oplus Q$ for some $R$-modules
$P,Q$. Let 
\[
i:P\to F,\qquad p:F\to P
\]
be the canonical inclusion and projection, so that $p\circ i=\operatorname{id}_P$.

To show that $P$ is projective, let $\pi:X\twoheadrightarrow Y$ be an epimorphism
of $R$-modules and let $\varphi:P\to Y$ be any $R$-morphism. Consider the composite
\[
\varphi\circ p: F\to Y.
\]
Since $F$ is free, it is projective; hence there exists an $R$-morphism
\[
\widetilde\Phi:F\to X
\quad\text{such that}\quad
\pi\circ \widetilde\Phi=\varphi\circ p.
\]
Now define
\[
\widetilde\varphi := \widetilde\Phi\circ i : P\to X.
\]
Then
\[
\pi\circ \widetilde\varphi
= \pi\circ \widetilde\Phi\circ i
= (\varphi\circ p)\circ i
= \varphi\circ (p\circ i)
= \varphi\circ \operatorname{id}_P
= \varphi.
\]
Thus every map $\varphi:P\to Y$ lifts along every epimorphism $\pi:X\twoheadrightarrow Y$,
so $P$ is projective.
\end{proof}

\begin{definition}[Tensor product of modules: linearize the bilinear map.]
Let $M, N$ be $R$-modules. The tensor product of $M$ and $N$ is a pair $(M \otimes N, t_{M, N})$ where $M \otimes N$ is an $R$-module and $t_{M, N}: M \times N \rightarrow M \otimes_{R} N$ is a bilinear map, usually denoted $(m, n) \mapsto m \otimes_{R} n$, which satisfies the following universal property: whenever $X$ is an $R$-module and $f: M \times N \rightarrow X$ is a bilinear map, there exists a unique $R$-morphism $\mathcal{F}: M \otimes_{R} N \rightarrow X$ such that the following diagram commutes.
\end{definition}

\textbf{Up to isomorphism, tensor product is unique:} Suppose that $(T_{1}, t_{1}), (T_{2}, t_{2})$ are two tensor products. We have $t_{2}=\varphi_{1} t_{1}$
\[
\begin{aligned}
t_{1} & =\varphi_{2} t_{2} \\
\Rightarrow t_{1} & =\varphi_{2} t_{2}=\varphi_{2} \varphi_{1} t_{1}
\end{aligned}
\]
Now $M \times N \xrightarrow{t_{1}} T_{1}$. Since $t_{1}=\varphi_{2} \varphi_{1} t_{1}$, by uniqueness $\varphi_{2} \varphi_{1}=\operatorname{Id}_{1} \Rightarrow T_{1} \cong T_{2}$. Similarly $\varphi_{1} \varphi_{2}=\operatorname{Id}$.

\textbf{Proof of existence:} Let $F$ be a free $R$-module with basis $M \times N$, say $\{e(m,n)\}_{m \in M, n \in N}$. Let $B$ be the submodule of $F$ generated by
\[
\begin{aligned}
& e\left(m_{1}+m_{2}, n\right)-e\left(m_{1}, n\right)-e\left(m_{2}, n\right), \\
& e\left(m, n_{1}+n_{2}\right)-e\left(m, n_{1}\right)-e\left(m, n_{2}\right), \\
& e(\lambda m, n)-\lambda e(m, n), \\
& e(m, \lambda n)-\lambda e(m, n).
\end{aligned}
\]
Then $M \otimes_{R} N=F / B$. Write $t(m, n):=m \otimes n$. The elements $\{m \otimes_{R} n\}_{m \in M, n \in N}$, called elementary tensors, generate $M \otimes_{R} N$.

Fix $m \in M$, the map $N \rightarrow M \otimes_{R} N, n \mapsto m \otimes_{R} n$ is a morphism of $R$-modules. Its image $m \otimes_{R} N=\{m \otimes_{R} n \mid n \in N\}$ is a submodule of $M \otimes_{R} N$ and is isomorphic to a quotient of $N$. The map might not be injective, i.e., we can have $m \neq 0$ but $m \otimes n = 0$.

For example, take $M=\mathbb{F}_{2}, N=\mathbb{F}_{3}$, then $M \otimes_{\mathbb{Z}} N=0$ even though $M$ and $N$ are non-zero. This can be shown by constructing $M \times N \xrightarrow{\varphi} P$ bilinear such that $\varphi(m, n) \neq 0$.


\begin{proposition}
Properties of tensor product. We have unique isomorphisms
\begin{align*}
& (M_{1} \otimes_{R} M_{2}) \otimes_{R} M_{3} \longrightarrow M_{1} \otimes_{R} (M_{2} \otimes_{R} M_{3}) \text{ associative} \\
& M_{1} \otimes_{R} M_{2} \longrightarrow M_{2} \otimes_{R} M_{1} \text{ commutative} \\
& M_{1} \otimes_{R} (M_{2} \oplus M_{3}) \xrightarrow{\sim} (M_{1} \otimes_{R} M_{2}) \oplus (M_{1} \otimes_{R} M_{3}) \text{ tensor product of} \\
& \text{Moreover,} \\
& R \otimes_{R} M \xrightarrow{\sim} M, \quad (\lambda \otimes_{R} m \mapsto \lambda m) \text{ sum of tensor product.}
\end{align*}

\textbf{Proof:}

(1) $(M_{1} \otimes_{R} M_{2}) \otimes_{R} M_{3} \xrightarrow{\sim} M_{1} \otimes_{R} (M_{2} \otimes_{R} M_{3})$

Define $M_{1} \times M_{2} \xrightarrow{f_{m_{3}}} M_{1} \otimes_{R} (M_{2} \otimes_{R} M_{3})$ by
\[
(m_{1}, m_{2}) \mapsto m_{1} \otimes (m_{2} \otimes m_{3}).
\]

$f_{m_{3}}$ is bilinear, hence $\exists!$ $R$-module homomorphism
\[
\begin{aligned}
\psi_{m_{3}}: M_{1} \otimes_{R} M_{2} & \rightarrow M_{1} \otimes_{R} (M_{2} \otimes_{R} M_{3}) \\
m_{1} \otimes m_{2} & \mapsto m_{1} \otimes (m_{2} \otimes m_{3}).
\end{aligned}
\]

Now we can define $M_{1} \otimes_{R} M_{2} \times M_{3} \xrightarrow{\varphi} M_{1} \otimes_{R} (M_{2} \otimes_{R} M_{3})$ by
\[
(m_{1} \otimes m_{2}, m_{3}) \mapsto m_{1} \otimes (m_{2} \otimes m_{3}) = \psi_{m_{3}}(m_{1} \otimes m_{2}).
\]

$\varphi$ is well-defined as $f_{m}$ is, $\forall m \in M_{3}$.

Then $\exists!$ $R$-morphism $(M_{1} \otimes_{R} M_{2}) \otimes_{R} M_{3} \xrightarrow{\theta} M_{1} \otimes_{R} (M_{2} \otimes_{R} M_{3})$ given by
\[
(m_{1} \otimes m_{2}) \otimes m_{3} \mapsto m_{1} \otimes (m_{2} \otimes m_{3}).
\]

Similarly, $\exists!$ $R$-morphism $M_{1} \otimes_{R} (M_{2} \otimes_{R} M_{3}) \xrightarrow{\rho} (M_{1} \otimes_{R} M_{2}) \otimes_{R} M_{3}$ given by
\[
m_{1} \otimes (m_{2} \otimes m_{3}) \mapsto (m_{1} \otimes m_{2}) \otimes m_{3}.
\]

$\theta \rho = \mathrm{Id}$, $\rho \theta = \mathrm{Id}$. Hence $M_{1} \otimes_{R} (M_{2} \otimes_{R} M_{3}) \simeq (M_{1} \otimes_{R} M_{2}) \otimes_{R} M_{3}$.

(2) $M_{1} \otimes_{R} M_{2} \xrightarrow{\sim} M_{2} \otimes_{R} M_{1}$

Define $M_{1} \times M_{2} \xrightarrow{\varphi} M_{2} \otimes_{R} M_{1}$ bilinear by
\[
(m_{1}, m_{2}) \mapsto m_{2} \otimes m_{1}.
\]

Hence $\exists!$ $R$-module morphism $M_{1} \otimes_{R} M_{2} \xrightarrow{\rho} M_{2} \otimes_{R} M_{1}$ given by
\[
m_{1} \otimes m_{2} \mapsto m_{2} \otimes m_{1}.
\]
If $\varphi \rho$ and $\rho \varphi$ are $\mathrm{Id}$ on $m_{1} \otimes m_{2}$ and $m_{2} \otimes m_{1}$ respectively. Since $m_{1} \otimes m_{2}$ generates $M_{1} \otimes M_{2}$, $m_{2} \otimes m_{1}$ generates $M_{2} \otimes M_{1}$, we have $\varphi \rho = \mathrm{Id}$, $\rho \varphi = \mathrm{Id}$. Hence $M_{1} \otimes M_{2} \xrightarrow{\sim} M_{2} \otimes M_{1}$.

\end{proposition}

\begin{theorem}
Canonical isomorphism.

Let $M, N, P$ be $R$-modules. We have a canonical isomorphism
\[
\begin{aligned}
\phi = \phi_{M, N, P}: \operatorname{Hom}_{R}\left(M \otimes_{R} N, P\right) & \longrightarrow \operatorname{Hom}_{R}\left(M, \operatorname{Hom}_{R}(N, P)\right) \\
f & \longmapsto \phi(f)(m)(n) = f(m \otimes n)
\end{aligned}
\]

The inverse $\phi^{-1}$ is given by
\[
\phi^{-1}(g)(m \otimes_{R} n) = g(m)(n) \quad \longleftarrow \quad g
\]
\end{theorem}

\begin{remark}
Canonical.

If $f: A \rightarrow A^{\prime}$, there exists a commutative diagram
\[
\operatorname{Hom}_{R}\left(A^{\prime} \otimes_{R} B, C\right) \xrightarrow{\phi_{A^{\prime}, B, C}} \operatorname{Hom}_{R}\left(A^{\prime}, \operatorname{Hom}_{R}(B, C)\right)
\]

If $f: A \rightarrow A^{\prime}$, there exists a commutative diagram
\[
\begin{aligned}
& \operatorname{Hom}_{R}\left(A^{\prime} \otimes_{R} B, C\right) \xrightarrow{\phi_{A^{\prime}, B, C}} \operatorname{Hom}_{R}\left(A^{\prime}, \operatorname{Hom}_{R}(B, C)\right) \\
& \quad \downarrow{\operatorname{Hom}_{R}(f \otimes \mathrm{id}, C)} \quad \quad \downarrow{\operatorname{Hom}_{R}(f, \operatorname{Hom}_{R}(B, C))} \\
& \operatorname{Hom}_{R}\left(A \otimes_{R} B, C\right) \xrightarrow{\phi_{A, B, C}} \operatorname{Hom}_{R}\left(A, \operatorname{Hom}_{R}(B, C)\right)
\end{aligned}
\]

Notation. Let $N, M_{i}\ (1 \leq i \leq n)$ be $R$-modules.
\[
L^{n}\left(M_{1}, \ldots, M_{n}; N\right) = \left\{f \mid f: M_{1} \times \cdots \times M_{n} \rightarrow N \text{ is an } n\text{-linear map}\right\}
\]

In particular,
\[
\begin{aligned}
L^{2}\left(M_{1}, M_{2}; N\right) & \simeq L^{1}\left(M_{1} \otimes M_{2}, N\right) \simeq \operatorname{Hom}_{R}\left(M_{1} \otimes M_{2}, N\right) \\
& \simeq L^{1}\left(M_{1}, L^{1}\left(M_{2}, N\right)\right)
\end{aligned}
\]
\end{remark}

\begin{proposition}
$M^{v} \otimes_{R} N \simeq \operatorname{Hom}_{R}(M, N)$.

Let $M, N$ be finitely generated free $R$-modules. Define $M^{v} = \operatorname{Hom}_{R}(M, R)$ (the dual of $M$). Then $M^{v} \otimes N \simeq \operatorname{Hom}_{R}(M, N) = L(M, N)$.
\end{proposition}

\begin{proof}
Define $\varphi: M^{v} \times N \rightarrow \operatorname{Hom}_{R}(M, N)$ by
\[
\begin{aligned}
(f, n) \mapsto \varphi(f, n): M & \rightarrow N \\
m & \mapsto f(m) n
\end{aligned}
\]
On the other hand, let $V_{1}, \dots, V_{n}$ be a basis of $M$. Define $V_{j}^{v}(V_{i}) = \delta_{ij}$. $V_{1}^{v}, \dots, V_{n}^{v}$ is a dual basis of $M^{v}$.

Define $\psi\colon \operatorname{Hom}_{R}(M, N) \rightarrow M^{v} \otimes_{R} N$ by
\[
g \mapsto \sum_{i=1}^{n} V_{i}^{v} \otimes g(V_{i}).
\]
Show $\varphi \psi = \mathrm{id}$, $\psi \varphi = \mathrm{id}$.
\end{proof}

\begin{proposition}
Lift $R$-module to $R^{\prime}$-module.

(1) Assume that $R$ is a subring of a ring $R^{\prime}$. Let $M$ be an $R$-module. Then the bilinear map
\[
\begin{aligned}
R^{\prime} \times (R^{\prime} \otimes_{R} M) &\longrightarrow R^{\prime} \otimes_{R} M \\
(\lambda_{1}^{\prime}, \lambda_{2}^{\prime} \otimes m) &\mapsto \lambda_{1}^{\prime} \lambda_{2}^{\prime} \otimes m
\end{aligned}
\]
defines a structure of $R^{\prime}$-module on the $R$-module $R^{\prime} \otimes_{R} M$.

(2) More generally, let $f\colon R \rightarrow T$ be a ring homomorphism. View $T$ as an $R$-module defined by $R \times T \rightarrow T$, $(\lambda, \mu) \mapsto f(\lambda)\mu$. The bilinear map $T \times (T \otimes_{R} M) \rightarrow T \otimes_{R} M$,
\[
(\mu_{1}, \mu_{2} \otimes m) \mapsto \mu_{1} \mu_{2} \otimes m
\]
defines a $T$-module structure on the $R$-module $T \otimes_{R} M$.

\textbf{Example:} Let $V$ be an $n$-dimensional $\mathbb{Q}$-vector space, $V \simeq \mathbb{Q}^{n}$. Then $\mathbb{C} \otimes_{\mathbb{Q}} V \simeq \mathbb{C} \otimes_{\mathbb{Q}} \mathbb{Q}^{n} \simeq \mathbb{C}^{n}$.
\end{proposition}

\begin{lemma}
$M / IM \simeq (R / I) \otimes_{R} M$.

Let $I$ be an ideal of $R$ and let $M$ be an $R$-module. The morphism of $R$-modules $M \rightarrow R / I \otimes_{R} M$ induces an isomorphism $M / IM \simeq (R / I) \otimes_{R} M$, given by $m \mapsto 1 \otimes m$.
\end{lemma}

\begin{example}
Counterexamples:
\begin{enumerate}
\item There is a module having no linearly independent elements.
\item There is a module, some of whose submodules have no complement.
\item There is a finitely generated module, one of its submodules is not finitely generated.
\end{enumerate}
\end{example}
2. There is a module, some of whose submodules have no complement.

3. There is a finitely generated module, one of its submodules is not finitely generated.

4. There is a module with a linearly dependent subset $S$ such that no element is a linear combination of other elements of $S$.

5. There is a module with a minimal generating set that is not a basis and a maximal linearly independent set that is not a basis.

6. There is a free module with non-free submodules and non-free quotient modules.

7. There is an $R$-module $M$: $\exists r \in R, m \in M$ with $r \neq 0, m \neq 0$, but $rm = 0$.

8. There is a free module such that:
\begin{itemize}
\item a linearly independent subset is not contained in any basis;
\item a generating subset is not contained in any basis.
\end{itemize}

\begin{definition}
\textbf{Category}

A category $\mathcal{C}$ consists of:
\begin{itemize}
\item A class of objects.
\item For each pair $(A, B)$, a set $\operatorname{Hom}(A, B)$ whose elements are called morphisms from $A$ to $B$.
\item For each triple of objects $(A, B, C)$, a map
\[
\operatorname{Hom}(A, B) \times \operatorname{Hom}(B, C) \rightarrow \operatorname{Hom}(A, C)
\]
denoted $(f, g) \mapsto g \circ f$.
\end{itemize}

It is assumed that the objects and morphisms satisfy the following conditions:
\begin{itemize}
\item If $(A, B) \neq (C, D)$, then $\operatorname{Hom}(A, B)$ and $\operatorname{Hom}(C, D)$ are disjoint.
\item If $f \in \operatorname{Hom}(A, B)$, $g \in \operatorname{Hom}(B, C)$, $h \in \operatorname{Hom}(C, D)$, then
\[
h \circ (g \circ f) = (h \circ g) \circ f \quad \text{(associative)}.
\]
\item For every object $A$ we have an element $\operatorname{id}_A \in \operatorname{Hom}(A, A)$ such that
\[
f \circ \operatorname{id}_A = f \text{ for any } f \in \operatorname{Hom}(A, B)
\]
and $\operatorname{id}_A \circ g = g$ for any $g \in \operatorname{Hom}(B, A)$ \text{(unital)}.
\end{itemize}
\end{definition}

\begin{remark}
Objects are in one-to-one correspondence with identity morphisms, which are uniquely determined by the property that they serve as two-sided identities for composition.
\end{remark}

\begin{example}
\begin{enumerate}
\item \textbf{Set}: objects are sets; morphisms are maps.
\item \textbf{Group}: objects are groups; morphisms are group homomorphisms.
\item \textbf{Ab}: objects are abelian groups; morphisms are group homomorphisms.
\item $R$-\textbf{Mod}: objects are $R$-modules; morphisms are $R$-homomorphisms.
\item \textbf{Top}: objects are topological spaces; morphisms are continuous functions.
\end{enumerate}
\end{example}

\paragraph{Lesson 3: Categories}
\textbf{Highlight:} Functor (covariant, contravariant), left exact/right exact, projective/injective/flat module

\textbf{Notation:}
\begin{itemize}
    \item $R\text{-}\mathrm{Mod}$ \ldots left $R$-modules
    \item $R\text{-}\mathrm{mod}$ \ldots finitely generated left $R$-modules
    \item $\mathrm{Mod}\text{-}R$ \ldots right $R$-modules
    \item $\mathrm{mod}\text{-}R$ \ldots finitely generated right $R$-modules
\end{itemize}

\begin{definition}[Subcategories]
A category $\mathcal{D}$ is called a \textbf{subcategory} of a category $\mathcal{C}$ if
\begin{enumerate}[(i)]
    \item the objects of $\mathcal{D}$ form a subclass of the objects of $\mathcal{C}$;
    \item for any objects $A, B$ of $\mathcal{D}$, one has that $\operatorname{Hom}_{\mathcal{D}}(A, B) \subseteq \operatorname{Hom}_{\mathcal{C}}(A, B)$, and the identity morphisms and the composition of morphisms for $\mathcal{D}$ are the same as for $\mathcal{C}$.
\end{enumerate}
The subcategory $\mathcal{D}$ is called \textbf{full} if $\operatorname{Hom}_{\mathcal{D}}(A, B) = \operatorname{Hom}_{\mathcal{C}}(A, B)$.
\end{definition}

\begin{example}
$\mathbf{Ab}$ is a full subcategory of $\mathbf{Grp}$.
\end{example}

\begin{definition}[$\mathcal{C}^{\mathrm{op}}$]
Let $\mathcal{C}$ be a category. Define $\mathcal{C}^{\mathrm{op}}$ as follows:
\begin{itemize}
    \item The objects of $\mathcal{C}^{\mathrm{op}}$ are the objects of $\mathcal{C}$.
    \item $\operatorname{Hom}_{\mathcal{C}^{\mathrm{op}}}(A, B) = \operatorname{Hom}_{\mathcal{C}}(B, A)$.
\end{itemize}
The composition satisfies: if $f \in \operatorname{Hom}_{\mathcal{C}^{\mathrm{op}}}(A, B)$ and $g \in \operatorname{Hom}_{\mathcal{C}^{\mathrm{op}}}(B, C)$, then $g \circ f$ in $\mathcal{C}^{\mathrm{op}}$ equals $f \circ g$ in $\mathcal{C}$. The identity morphism $\mathrm{id}_A$ in $\mathcal{C}^{\mathrm{op}}$ is the same as $\mathrm{id}_A$ in $\mathcal{C}$.

We call $\mathcal{C}^{\mathrm{op}}$ the \textbf{dual category} of $\mathcal{C}$.
\end{definition}

\begin{definition}[(Covariant) Functor]
Let $\mathcal{C}$ and $\mathcal{D}$ be categories. A \textbf{(covariant) functor} $F \colon \mathcal{C} \to \mathcal{D}$ consists of:
\begin{itemize}
    \item For each object $A$ of $\mathcal{C}$, there is a unique object $F(A)$ of $\mathcal{D}$.
    \item For each pair $(A, B)$ of objects of $\mathcal{C}$, $F$ gives a map
    \[
    \begin{aligned}
    \operatorname{Hom}_{\mathcal{C}}(A, B) &\to \operatorname{Hom}_{\mathcal{D}}(F(A), F(B)) \\
    f &\mapsto F(f).
    \end{aligned}
    \]
\end{itemize}
It is assumed that the following conditions are satisfied:
\begin{itemize}
    \item If $g \circ f$ is defined in $\mathcal{C}$, then $F(g) \circ F(f) = F(g \circ f)$.
    \item $F(\mathrm{id}_A) = \mathrm{id}_{F(A)}$.
\end{itemize}
\end{definition}

\begin{definition}[Contravariant Functor]
A \textbf{contravariant functor} from $\mathcal{C}$ to $\mathcal{D}$ is a covariant functor from $\mathcal{C}^{\mathrm{op}}$ to $\mathcal{D}$.
\end{definition}

\begin{example}[$\operatorname{Hom}_{R}(M, -)$ and $\operatorname{Hom}_{R}(-, N)$]
Let $M, N$ be $R$-modules. $\operatorname{Hom}_{R}(M, N) = \{R\text{-homomorphisms from } M \text{ to } N\}$ is an abelian group.
\end{example}

\begin{remark}
If $R$ is commutative, then $\operatorname{Hom}_{R}(M, N)$ is an $R$-module.

In general, if $R$ is noncommutative, $\operatorname{Hom}_{R}(M, N)$ is not an $R$-module in a natural way. For $f \in \operatorname{Hom}_{R}(M, N)$, define $\operatorname{Hom}_{R}(M, -) \colon R\text{-}\mathrm{Mod} \to \mathbf{Ab}$.
\end{remark}
\[
\begin{aligned}
& \Rightarrow f \in \operatorname{Hom}_{R}(M, N).
\end{aligned}
\]

\[
\begin{aligned}
& = f(srm)
\end{aligned}
\]

\[
\begin{aligned}
N & \xrightarrow{\quad} \operatorname{Hom}_{R}(N, M) \\
{\scriptstyle +}\downarrow & \uparrow {\scriptstyle \operatorname{Hom}_{R}(f, M)} \\
N' & \xrightarrow{\quad} \operatorname{Hom}_{R}(N', M)
\end{aligned} \text{ contravariant functor}
\]


\begin{definition}
Tensor product ($R$ non-commutative).

$M$: right $R$-module, $N$: left $R$-module.
$M \otimes_{R} N$: abelian group satisfying $\to$ no $R$-module structure.

\[
\begin{aligned}
(rm) \otimes n = m \otimes (rn) \quad (rs)(m \otimes n) & \to m(rs) \otimes n = r(s(m \otimes n)) \\
& \to m \otimes (rs)n = s(r(m \otimes n))
\end{aligned}
\]

\end{definition}

\begin{definition}
Bimodule (双模).

Let $R$ and $S$ be rings (with $1$). An abelian group $M$ is a left $R$-right $S$-bimodule if $M$ is both a left $R$-module and a right $S$-module for which the two scalar multiplications jointly satisfy
\[
r(xs) = (rx)s \quad \text{for } r \in R, x \in M, s \in S.
\]

Notation:
\[
\begin{aligned}
{}_{R}M &:= \text{left } R\text{-module} \\
M_{R} &:= \text{right } R\text{-module} \\
{}_{R}M_{S} &:= \text{left } R\text{-right } S\text{-bimodule}.
\end{aligned}
\]

\end{definition}

\begin{example}

Let $U := {}_{R}U_{S}$ be a bimodule.
Then for any $R$-module ${}_{R}M$, there are two $S$-modules:
\begin{itemize}
\item $\operatorname{Hom}_{R}(U, M) \ni f$, $s \in S$: $(sf)(u) = f(us)$, $\forall u \in U$ $\to$ left $S$-module.
\item $\operatorname{Hom}_{R}(M, U) \ni g$, $s \in S$: $(gs)(m) = g(m)s$, $\forall m \in M$ $\to$ right $S$-module.
\end{itemize}

\[
\begin{aligned}
\operatorname{Hom}_{R}(U, -): R\text{-}\mathbf{Mod} & \to S\text{-}\mathbf{Mod} \\
M & \mapsto \operatorname{Hom}(U, M)
\end{aligned}
\] covariant functor.

\[
\left(\begin{array}{c}
M \\
\downarrow {\scriptstyle f} \\
M'
\end{array}\right) \longmapsto \begin{aligned}
& \operatorname{Hom}(U, M) \\
& \downarrow {\scriptstyle \operatorname{Hom}(U, f)} \\
& \operatorname{Hom}(U, M')
\end{aligned}
\]

$\operatorname{Hom}_{R}(-, U): R\text{-}\mathbf{Mod} \to \mathbf{Mod}\text{-}S$ contravariant functor.

\[
\begin{aligned}
M & \mapsto \operatorname{Hom}(M, U) \\
\left(\begin{array}{c}
M \\
\downarrow {\scriptstyle f} \\
M'
\end{array}\right) & \mapsto \begin{aligned}
& \operatorname{Hom}(M, U) \\
& \uparrow {\scriptstyle \operatorname{Hom}(f, U)} \\
& \operatorname{Hom}(M', U)
\end{aligned}
\end{aligned}
\]

\end{example}
\begin{proposition}
Tensor product of bimodules: Let ${}_{S}M_{R}$ and ${}_{R}N$ be bimodules. Then $M \otimes_{R} N$ has a left $S$-module structure. If ${}_{R}N_{T}$, then $M \otimes_{R} N \in {}_{S}\textbf{Mod}$. If ${}_{R}N_{R}$, then $M \otimes_{R} N \in \textbf{Mod-}T$.
\end{proposition}

\begin{proof}
We define the scalar multiplication as
\[
s\left(\sum m_{i} \otimes n_{i}\right)=\left(\sum s m_{i}\right) \otimes n_{i}.
\]
We need to show well-definedness:
\[
\sum m_{i} \otimes n_{i}=\sum m_{j}^{\prime} \otimes n_{j}^{\prime} \Rightarrow s\left(\sum m_{i} \otimes n_{i}\right)=s\left(\sum m_{j}^{\prime} \otimes n_{j}^{\prime}\right).
\]
Define $\varphi: M \times N \rightarrow M \otimes_{R} N$ by $(m, n) \mapsto sm \otimes n$. It is easy to check that $\varphi$ is bilinear. Then there exists a unique $\psi: M \otimes_{R} N \rightarrow M \otimes_{R} N$ with $m \otimes n \mapsto sm \otimes n$, i.e., well-definedness is guaranteed.
\end{proof}

\begin{definition}
Let $M_{R}$, $M_{R}^{\prime}$, ${}_{R}N$, ${}_{R}N^{\prime}$ be modules, and let $f: M \rightarrow M^{\prime}$ and $g: N \rightarrow N^{\prime}$ be $R$-homomorphisms. Consider the bilinear map $M \times N \rightarrow M^{\prime} \otimes_{R} N^{\prime}$ given by
\[
(m, n) \mapsto f(m) \otimes g(n).
\]
Then there exists a unique group homomorphism $f \otimes g: M \otimes_{R} N \rightarrow M^{\prime} \otimes_{R} N^{\prime}$ defined by
\[
m \otimes n \mapsto f(m) \otimes g(n).
\]
If ${}_{S}M_{R}$ and ${}_{S}M_{R}^{\prime}$, then $f \otimes g$ is a left $S$-module homomorphism.

Note that $f \otimes g$ satisfies the following properties. For all $M_{R}$, $M_{R}^{\prime}$, ${}_{R}N$, ${}_{R}N^{\prime}$, and $f_{1}, f_{2}, f \in \operatorname{Hom}_{R}(M, M^{\prime})$, $g_{1}, g_{2}, g \in \operatorname{Hom}_{R}(N, N^{\prime})$:
\begin{enumerate}
\item $(f_{1}+f_{2}) \otimes g=(f_{1} \otimes g)+(f_{2} \otimes g)$
\item $f \otimes (g_{1}+g_{2})=f \otimes g_{1}+f \otimes g_{2}$
\item $f \otimes 0=0$, $0 \otimes g=0$
\item $1_{M} \otimes 1_{N}=1_{M \otimes N}$
\end{enumerate}
For all $M_{R}$, $M_{R}^{\prime}$, $M_{R}^{\prime \prime}$, ${}_{R}N$, ${}_{R}N^{\prime}$, ${}_{R}N^{\prime \prime}$, and $f: M \rightarrow M^{\prime}$, $f^{\prime}: M^{\prime} \rightarrow M^{\prime \prime}$, $g: N \rightarrow N^{\prime}$, $g^{\prime}: N^{\prime} \rightarrow N^{\prime \prime}$:
\[
(f^{\prime} \otimes g^{\prime}) \circ (f \otimes g)=(f^{\prime} \circ f) \otimes (g^{\prime} \circ g).
\]
\end{definition}

\begin{example}
Given ${}_{S}U_{R}$ a bimodule. Then $U \otimes_{R}-: R\text{-}\textbf{Mod} \rightarrow S\text{-}\textbf{Mod}$ is a covariant functor.
\end{example}
\[
\begin{aligned}
M & \mapsto U \otimes M \\
\left(\begin{array}{c}
M \\
f_{\downarrow} \\
M^{\prime}
\end{array}\right) & \mapsto\left(\begin{array}{c}
U \otimes M \\
\downarrow \mathrm{id} \otimes f \\
V \otimes M^{\prime}
\end{array}\right)
\end{aligned}
\]

$-\otimes_{S} U: \operatorname{Mod-}S \rightarrow \operatorname{Mod-}R$ is a covariant functor

\[
\begin{aligned}
M & \mapsto M \otimes_{S} U \\
\left(\begin{array}{c}
M \\
f_{\downarrow} \\
M^{\prime}
\end{array}\right) & \mapsto\left(\begin{array}{c}
M \otimes U \\
\downarrow f \otimes \mathrm{id} \\
M^{\prime} \otimes U
\end{array}\right)
\end{aligned}
\]


\begin{definition}[left exact \& right exact]
Let $F: R\text{-}\mathrm{Mod} \rightarrow S\text{-}\mathrm{Mod}/\mathrm{Ab}$ be a covariant functor. We say $F$ is left exact if for any exact sequence
\[ 0 \rightarrow N \xrightarrow{\varphi} M \xrightarrow{\psi} P, \]
the sequence
\[ 0 \rightarrow F(N) \xrightarrow{F(\varphi)} F(M) \xrightarrow{F(\psi)} F(P) \]
is exact.
\end{definition}

Let $F: R\text{-}\mathrm{Mod} \rightarrow S\text{-}\mathrm{Mod}/\mathrm{Ab}$ be a covariant functor. We say $F$ is left exact if for any exact sequence
\[ 0 \rightarrow N \xrightarrow{\varphi} M \xrightarrow{\psi} P, \]
the sequence
\[ 0 \rightarrow F(N) \xrightarrow{F(\varphi)} F(M) \xrightarrow{F(\psi)} F(P) \]
is exact. We say $F$ is right exact if for any exact sequence
\[ N \xrightarrow{\varphi} M \xrightarrow{\psi} P \rightarrow 0, \]
the sequence
\[ F(N) \xrightarrow{F(\varphi)} F(M) \xrightarrow{F(\psi)} F(P) \rightarrow 0 \]
is exact.

Let $G: R\text{-}\mathrm{Mod} \rightarrow \mathrm{Ab}$ be a contravariant functor. It is called left (or right) exact if it is left (or right) exact as a covariant functor $(R\text{-}\mathrm{Mod})^{\mathrm{op}} \rightarrow \mathrm{Ab}$. (i.e., left and right are interchanged in the corresponding definitions for contravariant functors.) A functor is called exact if it is both left exact and right exact.

\begin{theorem}[classical examples of left (right) exact functors]
\begin{enumerate}
\item ${}_{R}U_{S}$: The hom functor $\operatorname{Hom}_{R}(U,-): R\text{-}\mathrm{Mod} \rightarrow S\text{-}\mathrm{Mod}$ is left exact.
\item ${}_{S}U_{R}$: The hom functor $\operatorname{Hom}_{R}(-, U): \mathrm{Mod}\text{-}R \rightarrow \mathrm{Mod}\text{-}S$ is left exact (contravariant functor).
\item ${}_{S}U_{R}$: The functor $U \otimes_{R}-: R\text{-}\mathrm{Mod} \rightarrow S\text{-}\mathrm{Mod}$ is right exact.
\item ${}_{R}U_{S}$: The functor $-\otimes_{R} U: \mathrm{Mod}\text{-}R \rightarrow \mathrm{Mod}\text{-}S$ is right exact.
\end{enumerate}
\end{theorem}
\begin{proof}
(1) Aim: Prove if $0 \rightarrow N^{\prime} \xrightarrow{f} N \xrightarrow{g} N^{\prime \prime}$ is exact, then $0 \rightarrow \operatorname{Hom}_{R}\left(U, N^{\prime}\right) \xrightarrow{\operatorname{Hom}_{R}(U, f)} \operatorname{Hom}_{R}(U, N) \xrightarrow{\operatorname{Hom}_{R}(U, g)} \operatorname{Hom}_{R}\left(U, N^{\prime \prime}\right)$ is exact.

Step 1: $\operatorname{Hom}_{R}(U, f)$ is injective.

If $\operatorname{Hom}_{R}(U, f)(\varphi)=0$ for $\varphi \in \operatorname{Hom}_{R}\left(U, N^{\prime}\right)$, then $f \circ \varphi = 0$. Since $f$ is injective, $\varphi=0$.

Step 2: Prove $\operatorname{Im}\left(\operatorname{Hom}_{R}(U, f)\right) \subseteq \operatorname{Ker}\left(\operatorname{Hom}_{R}(U, g)\right)$.

Let $\varphi \in \operatorname{Hom}_{R}\left(U, N^{\prime}\right)$. Then $\operatorname{Hom}_{R}(U, g)(f \varphi)=g f \varphi=0$ since $g f=0$.

Step 3: Prove $\operatorname{Im}\left(\operatorname{Hom}_{R}(U, f)\right) \supseteq \operatorname{Ker}\left(\operatorname{Hom}_{R}(U, g)\right)$.

Let $\psi \in \operatorname{ker}\left(\operatorname{Hom}_{R}(U, g)\right)$, so $g \psi=0$. Then $\forall x \in U, g(\psi(x))=0 \Rightarrow \psi(x) \in \operatorname{ker}(g)=\operatorname{Im}(f)$. Since $f$ is injective, $\exists! y \in N^{\prime}$ s.t.\ $f(y)=\psi(x)$. Define $\varphi \in \operatorname{Hom}\left(U, N^{\prime}\right)$ by $\varphi(x)=y$. Then $\psi=f \varphi$.

(3) Aim: Prove if $N^{\prime} \xrightarrow{f} N \xrightarrow{g} N^{\prime \prime} \rightarrow 0$ is exact, then $U \otimes_{R} N^{\prime} \xrightarrow{1 \otimes f} U \otimes_{R} N \xrightarrow{1 \otimes g} U \otimes_{R} N^{\prime \prime} \rightarrow 0$ is exact.

Step 1: $1 \otimes g$ is surjective.

Every element of $U \otimes_{R} N^{\prime \prime}$ is of the form $\sum x_{i} \otimes g\left(y_{i}\right)=(1 \otimes g)\left(\sum x_{i} \otimes y_{i}\right) \in \operatorname{Im}(1 \otimes g)$ (since $g$ is surjective).

Step 2: $\operatorname{Im}(1 \otimes f)=\operatorname{ker}(1 \otimes g)$.

$g f=0 \Rightarrow(1 \otimes g)(1 \otimes f)=1 \otimes(g f)=0 \Rightarrow \operatorname{Im}(1 \otimes f) \subseteq \operatorname{ker}(1 \otimes g)$.

$1 \otimes g$ induces $\theta: U \otimes N / \operatorname{Im}(1 \otimes f) \rightarrow U \otimes N^{\prime \prime}$. Then $\operatorname{Im}(1 \otimes f)=\operatorname{ker}(1 \otimes g) \Leftrightarrow \theta$ is an isomorphism.

Let $x \in U$, $y^{\prime \prime} \in N^{\prime \prime}$. Choose $y \in N$ such that $g(y)=y^{\prime \prime}$.

Claim: $\overline{x \otimes y}$ is independent of the choice of $y$.

Let $y_{1}, y_{2} \in N$ with $g\left(y_{1}\right)=g\left(y_{2}\right)=y^{\prime \prime}$. Then $y_{1}-y_{2} \in \operatorname{ker} g=\operatorname{Im} f$. Thus $x \otimes y_{1}-x \otimes y_{2}=x \otimes\left(y_{1}-y_{2}\right) \in \operatorname{Im}(1 \otimes f)$. Hence the claim holds.

Therefore we can define
\begin{align*}
\tilde{\theta}: U \times N^{\prime \prime} & \rightarrow U \otimes_{R} N / \operatorname{Im}(1 \otimes f) \\
\left(x, y^{\prime \prime}\right) & \mapsto \overline{x \otimes y}
\end{align*}
\end{proof}
\[
(x, y'') \mapsto \pi \otimes y^{k} \quad \text{where } g(y)=y''
\]

Check $\theta\theta'=\operatorname{id}_{N''}$, $\theta'\theta=\operatorname{id}_{N}/\operatorname{Im}(1 \otimes f)$, hence $\theta$ is an isomorphism.


\begin{definition}[Projective/Injective/Flat Module]
Let $M$ be a left $R$-module.
\begin{itemize}
\item $M$ is called a \textbf{projective} (投射) $R$-module if $\operatorname{Hom}_{R}(M, -)$ is exact.
\item $M$ is called an \textbf{injective} (内射) $R$-module if $\operatorname{Hom}_{R}(-, M)$ is exact.
\item $M$ is called a \textbf{flat} (平坦) $R$-module if $M \otimes_{R} -$ is exact.
\end{itemize}
\end{definition}

\paragraph{Equivalent Definitions of Projective Module}

(1) Whenever there is an $R$-epimorphism $M \xrightarrow{g} N$ and an $R$-homomorphism $P \xrightarrow{\gamma} N$, there exists an $R$-module homomorphism $\bar{\gamma}\colon P \to M$ such that $\gamma=g\bar{\gamma}$.
\[
\begin{aligned}
& \bar{\gamma} \downarrow \quad \quad \gamma \\
& M \longrightarrow N \longrightarrow 0
\end{aligned}
\]

(2) For each epimorphism $f\colon {}_{R}M \to {}_{R}N$, the map $\operatorname{Hom}_{R}(P, f)\colon \operatorname{Hom}_{R}(P, M) \to \operatorname{Hom}_{R}(P, N)$ is an epimorphism.

(3) For every exact sequence $M' \xrightarrow{f} M \xrightarrow{g} M''$ in $R\text{-}\mathrm{Mod}$, the sequence
\[
\operatorname{Hom}_{R}(P, M') \xrightarrow{\operatorname{Hom}_{R}(P, f)} \operatorname{Hom}_{R}(P, M) \xrightarrow{\operatorname{Hom}_{R}(P, g)} \operatorname{Hom}_{R}(P, M'')
\]
is exact.

(4) $P$ is projective.

(5) Every epimorphism $M \xrightarrow{f} P$ splits, i.e., there exists an $R$-homomorphism $g\colon P \to M$ such that $fg=\operatorname{id}_{P}$.

(6) $P$ is isomorphic to a direct summand of a free left $R$-module.

\begin{proof}
\section*{Split Exact Sequences}

Let $0 \rightarrow A_{1} \xrightarrow{f} B \xrightarrow{g} A_{2} \rightarrow 0$ be exact. TFAE:
\begin{enumerate}
    \item $\exists h: A_{2} \rightarrow B$ with $g \circ h = \mathrm{id}_{A_{2}}$
    \item $\exists K: B \rightarrow A_{1}$ with $K \circ f = \mathrm{id}_{A_{1}}$
    \item The given sequence is isomorphic to $0 \rightarrow A_{1} \xrightarrow{\iota_{1}} A_{1} \oplus A_{2} \xrightarrow{\pi_{2}} A_{2} \rightarrow 0$
\end{enumerate}

It is easy to show $(1) \Leftrightarrow (2)$ are equivalent.

$(1) \Leftrightarrow (4)$

$(4) \Rightarrow (1)$ obvious.

$(1) \Rightarrow (4)$: If $M \rightarrow N \rightarrow 0$ is exact, then $\operatorname{Hom}(P, M) \rightarrow \operatorname{Hom}(P, N) \rightarrow 0$ is exact.

With identity maps on $A_{1}$ and $A_{2}$, in particular, $B \simeq A_{1} \oplus A_{2}$.

Let $0 \rightarrow A_{1} \xrightarrow{f} B \xrightarrow{g} A_{2} \rightarrow 0$ be exact with $\varphi_{1}, \varphi_{2}, \varphi_{3}$ isomorphisms.

$(1) \Leftrightarrow (4)$

$(4) \Rightarrow (1)$ obvious.

$(1) \Rightarrow (4)$: We have $M \rightarrow N \rightarrow 0$ exact implies $\operatorname{Hom}(P, M) \rightarrow \operatorname{Hom}(P, N) \rightarrow 0$ exact, i.e., if $M^{\prime} \xrightarrow{f} M \xrightarrow{g} M^{\prime \prime} \rightarrow 0$ is exact, then consider
\[
\operatorname{Hom}\left(P, M^{\prime}\right) \xrightarrow{f^{*}} \operatorname{Hom}(P, M) \xrightarrow{g^{*}} \operatorname{Hom}\left(P, M^{\prime \prime}\right) \rightarrow 0.
\]

We already have $g^{*}$ is surjective. $g^{*} f^{*}=0$ follows from $g \circ f=0$.

Then it suffices to show $\operatorname{ker} g^{*} \subseteq \operatorname{im} f^{*}$.

Let $\varphi \in \operatorname{ker} g^{*}$, i.e., $g \circ \varphi=0 \Rightarrow \varphi(p) \in \operatorname{ker} g=\operatorname{im} f$ for any $p \in P$.

Then for any $p \in P$, $\exists y \in M^{\prime}$ s.t.\ $\varphi(p)=f(y)$.

We then can define $\theta: P \rightarrow M^{\prime} / \operatorname{ker} f;\, p \mapsto y+\operatorname{ker} f$.

Then apply (2) again, $\theta$ lifts to $\theta^{\prime}: P \rightarrow M^{\prime}$ satisfying $\varphi=f \circ \theta^{\prime}$.

$(3) \Leftrightarrow (4)$ is similar to $(1) \Leftrightarrow (4)$.

Now we prove $(1) \Rightarrow (5) \Rightarrow (6) \Rightarrow (1)$.

$(1) \Rightarrow (5)$

$(5) \Rightarrow (6)$: Every module is a quotient of free modules.

Take $F \xrightarrow{\varphi} P \rightarrow 0$ where $F$ is a free module.

Then by $(6)$, $0 \rightarrow \operatorname{ker} \varphi \hookrightarrow F \rightarrow P \rightarrow 0$ splits, i.e., $F=\operatorname{ker} \varphi \oplus P$.

$(6) \Rightarrow (1)$: It follows since every free module is a projective module.

\section*{Equivalent Definition of Injective Module}

\begin{enumerate}
    \item $U$ is injective.
    \item For each monomorphism $f: K \rightarrow M$ and each homomorphism $\gamma: K \rightarrow U$ there is an $R$-homomorphism $\bar{\gamma}: M \rightarrow U$ such that the diagram
\end{enumerate}
\[
\begin{array}{r}
U \\
\gamma \uparrow \\
0 \rightarrow k \rightarrow \bar{\gamma} \\
0 \rightarrow M
\end{array}
\]
commutes.

(3) For every monomorphism $f: K \rightarrow M$, the map $\operatorname{Hom}(f, U): \operatorname{Hom}(M, U) \rightarrow \operatorname{Hom}(K, U)$ is an epimorphism.

(4) For every exact sequence $M^{\prime} \rightarrow M \xrightarrow{g} M^{\prime \prime}$ the sequence
\[
\operatorname{Hom}\left(M^{\prime \prime}, U\right) \xrightarrow{g^{*}} \operatorname{Hom}(M, U) \rightarrow \operatorname{Hom}\left(M^{\prime}, U\right)
\]
is exact.

(5) Every monomorphism $U \xrightarrow{f} M$ splits, i.e., there exists $g: M \rightarrow U$ such that $g \circ f = \operatorname{Id}_U$.

\end{proof}

\begin{theorem}
(Schanuel's lemma)
Let
\[
0 \rightarrow K \rightarrow P \xrightarrow{\lambda} V \rightarrow 0
\]
\[
0 \rightarrow K^{\prime} \rightarrow P^{\prime} \xrightarrow{\lambda^{\prime}} V \rightarrow 0
\]
be short exact sequences of $R$-modules with $P, P^{\prime}$ projective. Then
\[
K^{\prime} \oplus P \simeq K \oplus P^{\prime}.
\]
\end{theorem}

\begin{proof}
Let
\[
Y=\left\{\left(p, p^{\prime}\right) \in P \oplus P^{\prime} \mid \lambda(p)=\lambda^{\prime}\left(p^{\prime}\right)\right\},
\]
called the pullback of $\left(\lambda, \lambda^{\prime}\right)$. Then $Y$ is an $R$-module.

Define $\mu: Y \rightarrow P$ by $\left(p, p^{\prime}\right) \mapsto p$ and $\mu^{\prime}: Y \rightarrow P^{\prime}$ by $\left(p, p^{\prime}\right) \mapsto p^{\prime}$.

Then
\[
\ker \mu=\left\{\left(0, p^{\prime}\right) \in P \oplus P^{\prime} \mid \lambda^{\prime}\left(p^{\prime}\right)=0\right\} \simeq \ker \lambda^{\prime}=K^{\prime}
\]
and
\[
\ker \mu^{\prime}=\left\{(p, 0) \in P \oplus P^{\prime} \mid \lambda(p)=0\right\} \simeq \ker \lambda=K.
\]

Then we have the short exact sequences
\[
\begin{aligned}
& 0 \rightarrow K^{\prime} \rightarrow Y \rightarrow P \rightarrow 0, \\
& 0 \rightarrow K \rightarrow Y \rightarrow P^{\prime} \rightarrow 0.
\end{aligned}
\]

Since $P, P^{\prime}$ are projective, we obtain $K^{\prime} \oplus P \simeq K \oplus P^{\prime}$.

\end{proof}

\begin{proposition}
Equivalent definitions of Noetherian module and Artinian module.

For an $R$-module $M$, the following are equivalent:
\begin{enumerate}
\item A.C.C. condition holds.
\item Every submodule of $M$ is finitely generated.
\item Every nonempty set of submodules of $M$ has a maximal element with respect to inclusion.
\end{enumerate}
A module satisfying the above conditions is called a Noetherian module.
\end{proposition}
For an $R$-module $M$, TFAE:
\begin{enumerate}
\item[(1)] d.c.c. condition
\item[(2)] every nonempty set of submodules of $M$ has a minimal element with respect to inclusion.
\end{enumerate}
A module satisfying the above conditions is called an \emph{Artinian module}. If $M$ is a Noetherian (or Artinian) module, we also say that $M$ satisfies the maximal (minimal) / increasing (decreasing) chain condition on submodules.

\begin{proposition}
For an $R$-module $M$, TFAE:
\begin{enumerate}
\item[(1)] d.c.c. condition
\item[(2)] every nonempty set of submodules of $M$ has a minimal element with respect to inclusion.
\end{enumerate}
A module satisfying the above conditions is called an Artinian module. If $M$ is a Noetherian (or Artinian) module, we also say that $M$ satisfies the maximal (minimal) / increasing (decreasing) chain condition on submodules.
\end{proposition}

\begin{definition}
Noetherian \& Artinian ring. $R$ with ${}_R R$ is said to be a left Noetherian (or Artinian) ring if ${}_R R$ is Noetherian (or Artinian) as a left module. (Similar for ``right'')
\end{definition}

\begin{lemma}
Properties of Noetherian \& Artinian modules.
\begin{enumerate}
\item[(1)] Let $0 \rightarrow M' \rightarrow M \rightarrow M'' \rightarrow 0$ be a short exact sequence. Then $M$ is Noetherian (or Artinian) $\Leftrightarrow M'$ and $M''$ are Noetherian (or Artinian).
\item[(2)] Let $M = \sum_{i=1}^{n} M_i$. Then $M$ is Noetherian (or Artinian) $\Leftrightarrow M_i$ is Noetherian (or Artinian) for any $i$.
\end{enumerate}
\end{lemma}

\begin{theorem}
$R$ left Noetherian $\Rightarrow$ f.g. $R$-modules Noetherian.

If $R$ is a (left) Noetherian (or Artinian) ring, then every finitely generated (left) $R$-module $M$ is a Noetherian (or Artinian) module.
\end{theorem}

\begin{proof}
$M = R m_1 + \cdots + R m_n$ for $m_1, \ldots, m_n \in M$. We have $\varphi_i: R \rightarrow R m_i$ given by $r \mapsto r m_i$, so every $R m_i$ is a homomorphic image of ${}_R R$. Hence $R m_i$ is Noetherian. Hence $M$ is Noetherian by lemma.
\end{proof}

\begin{definition}
Semisimple module. A module that is the direct sum of simple submodules is called \emph{semisimple} or \emph{completely reducible}.
\end{definition}

\paragraph{Lemma}
Let $R$ be a ring with ${}_R R$. Suppose that $U = S_1 + \cdots + S_n$ is an $R$-module that can be expressed as a finite sum of simple submodules. If $V$ is any submodule of $U$, then there is a subset $I \subseteq \{1, \ldots, n\}$ such that
\[ U = V \oplus \left( \bigoplus_{i \in I} S_i \right). \]
In particular:
\begin{enumerate}
\item[(1)] $V$ is a direct summand of $U$.
\item[(2)] $U$ is a direct sum of some subset of the $S_i$'s and hence semisimple.
\end{enumerate}

\textbf{Proof:} Choose a subset $I$ of $\{1, \ldots, n\}$ maximal subject to the condition that the sum $W = V \oplus \left( \bigoplus_{i \in I} S_i \right)$ is a direct sum.

Claim: $W = U$. If $W \neq U$, then $S_j \not\subseteq W$ for some $j$. Note: $S_j \cap W = 0 \Rightarrow S_j + W = S_j \oplus W$, which contradicts the maximality of $I$. Hence $W = U$.

\begin{corollary}
Equivalent definitions of semisimple $R$-modules.
Let $U$ be an $R$-module. TFAE:
\begin{enumerate}
    \item $U$ is semisimple (true for infinite sum).
    \item $U$ is a sum of simple $R$-submodules. (Use Zorn's lemma)
    \item Any $R$-submodule of $U$ is a direct summand of $U$.
\end{enumerate}

\end{corollary}

\begin{proof}
(1)$\Rightarrow$(3): Same proof as above but apply Zorn's lemma.

(3)$\Rightarrow$(2):
$1^{\circ}$ Every nonzero submodule of $V$ has a simple submodule.
Take $x \neq 0$ in $U$. $Rx$ has a maximal submodule, say $M$ ($R$ has $1_{R}$).
Then $U = M \oplus M'$ for some $M' \subset U$. Then
\[ Rx = Rx \cap U = Rx \cap (M \oplus M') = M \oplus (Rx \cap M') \]
where $Rx \cap M' \simeq Rx / M$ is simple. So $Rx$ has a simple module.

$2^{\circ}$ Let $N$ be the sum of all simple submodules of $U$. Then $U = N \oplus N'$ for some $N'$.
Since $N \cap N' = 0$, $N'$ has no simple submodules $\Rightarrow N' = 0$. So $U = N$.

(2)$\Rightarrow$(1): Apply $V = 0$ in the proof of $(1) \Rightarrow (3)$.
\end{proof}

\begin{definition}[Socle of $U$]
Prove by lemma $\rightarrow$ requires finite sum.
Prove by theorem $\rightarrow$ holds for infinite sum.
Let $U$ be an $R$-module (with composition series of finite length).
\begin{enumerate}
    \item The sum of all simple submodules of $U$ is a semisimple module that is the unique largest semisimple submodule of $U$ (up to isomorphism), called the \textbf{socle} of $U$, denoted by $\operatorname{Soc}(U)$.
    \item The sum of all submodules of $U$ isomorphic to some given simple module $S$ is a submodule isomorphic to a direct sum of copies of $S$. It is the unique largest submodule with this property.
\end{enumerate}
\end{definition}

\begin{corollary}
Let $U = S_{1}^{a_{1}} \oplus \dots \oplus S_{n}^{a_{n}}$ be a semisimple module where the $S_{i}$'s are nonisomorphic simple $R$-modules.
Then each submodule $S_{i}^{a_{i}}$ is uniquely determined and characterized as the unique largest submodule of $U$ expressed as a direct sum of copies of $S_{i}$.
\end{corollary}

\begin{proof}
It suffices to show that $S_{i}^{a_{i}}$ contains every submodule of $U$ isomorphic to $S_{i}$.

Let $T$ be any nonzero submodule of $U$ such that $T \nsubseteq S_{i}^{a_{i}}$.
Let $\pi_{i}: U \rightarrow S_{i}^{a_{i}}$ be the projection. There exists $j \neq i$ such that $\pi_{j}(T) \neq 0$.
Further consider the projection $T \rightarrow S_{j}$.
If $T$ is simple, $T \simeq S_{j} \neq S_{i}$.
\end{proof}
\begin{theorem}
Let $R$ be a ring with $1$. Let ${}_{R}V = V_{1} \oplus V_{2} \oplus \cdots \oplus V_{n}$ be a direct sum of simple $R$-modules $V_{i}$ such that $V_{i} \simeq V_{j}$ for all $i, j$. Then $\operatorname{End}_{R}(V)$ is isomorphic to the full matrix ring of degree $n$ over the division ring $\operatorname{End}_{R}(V_{1})$.
\end{theorem}

\begin{proof}
Let $\pi_{i}: V \rightarrow V_{i}$ be the projection.

Fix $\theta_{i}: V_{1} \xrightarrow{\sim} V_{i}$ an isomorphism. For $b \in \operatorname{End}_{R}(V)$, let $\varphi_{ij}(b) = \theta_{i}^{-1} \pi_{i} b \theta_{j} \in \operatorname{End}_{R}(V_{1})$.

\[
V_{1} \xrightarrow{\theta_{j}} V_{j} \xrightarrow{b} V \xrightarrow{\pi_{i}} V_{i} \xrightarrow{\theta_{i}^{-1}} V_{1}
\]

Then we have $\varphi: \operatorname{End}_{R}(V) \rightarrow M_{n}(E_{1})$ where $E_{1} = \operatorname{End}_{R}(V_{1})$:

\[
b \mapsto \left(\varphi_{ij}(b)\right)_{n \times n}
\]

Check: $\varphi$ is a ring homomorphism.

\[
\text{For } b\beta: \quad V_{i} \xrightarrow{\pi_{i}} V \xrightarrow{b} V \xrightarrow{\beta} V \xrightarrow{\pi_{k}} V_{k}
\]

\[
\sum_{j} \theta_{k}^{-1} \pi_{j} \beta \pi_{i} b \theta_{i} = \sum_{j} \theta_{k}^{-1} \pi_{j} \beta \theta_{j} \theta_{j}^{-1} \pi_{i} b \theta_{i} = \sum_{j} \varphi_{jk}(\beta) \varphi_{ij}(b)
\]

Next, $\varphi$ is an isomorphism.

$1^{\circ}$ Injective: If $\varphi(b) = 0$, then $\theta_{i}^{-1} \pi_{i} b \theta_{j} = 0$ for all $i, j$, which implies $\pi_{i} b(V_{j}) = 0$ for all $i, j$. Then $b(V_{j}) = 0$ for all $j$, so $b = 0$.

$2^{\circ}$ Surjective: Let $(b_{ij}) \in M_{n}(E_{1})$.
\end{proof}

\begin{remark}
In $\varphi_{ij}(b) = \theta_{i}^{-1} \pi_{i} b \theta_{j}$, let $\tau_{i}: V_{i} \hookrightarrow V$ be the inclusion. Then $V_{i} \simeq (0, \ldots, V_{i}, \ldots, 0)$ in $V$. Set $b = \sum_{i,j} \tau_{i} \theta_{i} b_{ij} \theta_{j}^{-1} \pi_{j} \in \operatorname{End}_{R}(V)$ (no need for inclusion). Hence we might not need $\tau_{i}$.
\end{remark}

\begin{definition}
essential
An epimorphism of modules $f: U \rightarrow V$ is called \emph{essential} if no proper submodule of $U$ is mapped surjectively onto $V$ by $f$.

Equivalently, whenever $g: W \rightarrow U$ is a map such that $fg$ is an epimorphism, then $g$ is an epimorphism.

(Proof: Let $W$ be a proper submodule. Then $W \xrightarrow{g} U \xrightarrow{f} V$ with $fg$ epimorphism $\Rightarrow$ induces $\operatorname{im} g \rightarrow V \Rightarrow \operatorname{im} g = U$.)

\end{definition}

\begin{definition}
A \emph{projective cover} of a module $U$ is an essential epimorphism $P \rightarrow U$ where $P$ is projective.
\end{definition}

\begin{proposition}
``Minimality'' of a projective cover.
\begin{enumerate}
\item Suppose that $f: P \rightarrow U$ is a projective cover of a module $U$ and $g: Q \rightarrow U$ is an epimorphism where $Q$ is projective. Then $Q = Q_{1} \oplus Q_{2}$ so that $g$ has the component $g = (g_{1}, 0)$ with respect to the direct sum decomposition and $g_{1}: Q_{1} \rightarrow U$ satisfies that the diagram commutes where $\gamma$ is an isomorphism.
\item If it exists, the projective cover of a module $U$ is unique up to isomorphism, by isomorphisms that commute with the essential epimorphism.
\end{enumerate}
\end{proposition}

\begin{proof}
(1) Consider:
$Q$ projective $\Rightarrow \exists \beta$ such that $g = f\beta$.
$P$ projective $\Rightarrow \exists \alpha$ such that $f = g\alpha$.
$f = g\alpha = f(\beta\alpha)$ epimorphism $\Rightarrow \beta\alpha$ is an epimorphism.
In particular, $\beta$ is an epimorphism.
\[
Q \xrightarrow{\beta} P \rightarrow 0 \quad \text{exact}
\]
$P$ projective $\Rightarrow \beta$ splits.
\[
Q = Q_{1} \oplus Q_{2}, \quad Q_{2} = \ker \beta, \quad \beta|_{Q_{1}}: Q_{1} \rightarrow P \text{ isomorphism.}
\]
$g = (f\beta|_{Q_{1}}, 0)$, $\gamma = \beta|_{Q_{1}}$.

(2) $f: P \rightarrow U$, $g: Q \rightarrow U$ projective covers.
\[
Q = Q_{1} \oplus Q_{2}, \quad g = (g_{1}, 0), \quad g \text{ essential } \Rightarrow Q_{2} = 0.
\]
\end{proof}

\paragraph{Lesson 5 Radical}

\paragraph{$R$: ring with 1}

\begin{definition}
The \emph{radical} of a module $M$ is defined as
\[
\operatorname{Rad}(M) = \bigcap_{N} \{ \text{maximal submodules of } M \}.
\]
If $M$ has no maximal submodule, we set $\operatorname{Rad}(M) = M$.
\end{definition}
\begin{lemma}[Existence of maximal submodule]
Let $W$ be a proper $R$-submodule of an $R$-module $V$. If $V/W$ is finitely generated over $R$, then there exists a maximal $R$-submodule of $V$ containing $W$.

In particular, if $V$ is finitely generated, then there exists a maximal submodule of $V$ containing $W$.

\textbf{Corollary:} For every proper left ideal $I$ of $R$, there exists a maximal left ideal of $R$ containing $I$.
\end{lemma}

\begin{proof}
Let $\bar{V}=V/W=\sum_{i=1}^{n} R\overline{v}_{i}$, where $v_{i} \in V$.

Let $\theta$ denote the set of $R$-submodules of $V$ containing $W$.

\textbf{Claim:} An arbitrary totally ordered subset in $\theta$, say $\{W_{\lambda}\mid \lambda \in \Lambda\}$, has an upper bound.

Let $U=\bigcup_{\lambda \in \Lambda} W_{\lambda}$. If $U=V$, then for each $1 \leq i \leq n$, there exists $W_{\lambda_{i}}$ such that $v_{i} \in W_{\lambda_{i}}$. Let $W_{\alpha}$ be the largest among $W_{\lambda_{i}}$. Then $W_{\alpha}=V$, since $v_{i} \in W_{\alpha}$ for all $i$, which contradicts $W_{\alpha} \in \theta$.

So $U \neq V$, and $U \in \theta$ is an upper bound for $\{W_{\lambda}\}_{\lambda \in \Lambda}$.

Then we can apply Zorn's lemma.
\end{proof}

\begin{lemma}
Let $M, N$ be $R$-modules and $f\colon M \to N$ an $R$-morphism.
\begin{enumerate}
\item $f(\operatorname{Rad}(M)) \subseteq \operatorname{Rad}(N)$.
\item If $f$ is an epimorphism with $\ker(f) \subseteq \operatorname{Rad}(M)$, then $\operatorname{Rad}(N)=f(\operatorname{Rad}(M))$.
\end{enumerate}
\end{lemma}

\begin{lemma}
Let $U$ be an $R$-module.
\begin{enumerate}
\item Suppose that $M_{1}, \ldots, M_{n}$ are maximal submodules of $U$. Then there is a subset $I \subseteq \{1,2,\ldots, n\}$ such that
\[
U/(M_{1} \cap \cdots \cap M_{n}) \simeq \bigoplus_{i \in I} U/M_{i}.
\]
\item Suppose that $U$ is Artinian. Then $U/\operatorname{Rad} U$ is semisimple and $\operatorname{Rad} U$ is a submodule of $U$ with semisimple quotient.
\end{enumerate}
\end{lemma}

\begin{proof}
(1) Let $I$ be a subset of $\{1,2,\ldots, n\}$ maximal with the property that the quotient homomorphisms
\[
U/\bigcap_{i \in I} M_{i} \to U/M_{i}
\]
induce an isomorphism
\[
U/\bigcap_{i \in I} M_{i} \to \bigoplus_{i \in I} U/M_{i}.
\]
We shall show: $\bigcap_{i \in I} M_{i}=M_{1} \cap \cdots \cap M_{n}$.

If not, there exists $j$ such that $\bigcap_{i \in I} M_{i} \not\subseteq M_{j}$.

Consider the homomorphism $f\colon U \to \left(\bigoplus_{i \in I} U/M_{i}\right) \oplus \left(U/M_{j}\right)$ defined by
\[
u \mapsto \left((u+M_{i})_{i \in I}, u+M_{j}\right).
\]
\end{proof}
$f$ has kernel $M_{j} \cap \left(\bigcap_{i \in I} M_{i}\right)$. Then it suffices to show $f$ is surjective.

Let $g: U \rightarrow \left(U / \bigcap_{i \in I} M_{i}\right) \oplus \left(U / M_{j}\right)$.

Observe that $M_{j} + \bigcap_{i \in I} M_{i} = U$.

$\forall x \in U$, then $x = y + z$ for some $y \in \bigcap_{i \in I} M_{i}$, $z \in M_{j}$.

$g(y) = \left(0, x + M_{j}\right)$, \quad $g(z) = \left(x + \bigcap_{i \in I} M_{i}, 0\right)$.

So $g$ is surjective $\Rightarrow f$ is surjective.

(2) $U$ is Artinian then $\operatorname{Rad} U$ is the intersection of finitely many maximal submodules so $U / \operatorname{Rad}(U)$ is semisimple.

If $V$ is a submodule of $U$ with semisimple $U/V$. Then $U/V \simeq S_{1} \oplus \cdots \oplus S_{n}$ where $S_{i}$ is simple.

Let $M_{i}$ be the kernel of the composition $U \rightarrow U/V \rightarrow S_{i}$. Then $M_{i}$ is maximal since $S_{i}$ is simple.

And $V = M_{1} \cap \cdots \cap M_{n} \Rightarrow \operatorname{Rad}(U) \subseteq V$.


\begin{theorem}
(Nakayama's lemma)

Let $U$ be Noetherian, then $U \rightarrow U / \operatorname{Rad} U$ is essential. Equivalently, if $V$ is a submodule of $U$ with $V + \operatorname{Rad} U = U$, then $U = V$.

\textbf{Proof:}
If $V \neq U$, then $\exists$ a maximal submodule $M$ of $U$ such that $V \subseteq M \subseteq U$. Now $V + \operatorname{Rad} U \subseteq M$ and so $V \subseteq M/\operatorname{Rad} U$ which is not equal to $U/\operatorname{Rad} U$ since $(U/\operatorname{Rad} U)/(M/\operatorname{Rad} U) \simeq U/M \neq 0$.
\end{theorem}

\begin{proposition}
(Sufficient conditions of essential epimorphism for Noetherian modules)

(1) Suppose that $f: U \rightarrow V$, $g: V \rightarrow W$ are two $R$-morphisms of $R$-modules. If two of $f, g, gf$ are essential epi, then so is the third.

(2) Let $f: U \rightarrow V$ be a homomorphism of Noetherian modules. Then $f$ is an essential epimorphism if and only if the radical quotient $U/\operatorname{Rad} U \rightarrow V/\operatorname{Rad} V$ is an essential epimorphism. (Semisimple case: isomorphism)

(3) Let $f_{i}: U_{i} \rightarrow V_{i}$ be homomorphisms of Noetherian modules for $i = 1, 2, \cdots, n$. Then $f_{i}$ are all essential epimorphisms $\Leftrightarrow$ $\bigoplus f_{i}: \bigoplus U_{i} \rightarrow \bigoplus V_{i}$ is an essential epimorphism.
\end{proposition}
\begin{proof}
(1) $1^{\circ}$ $f$, $g$ essential $\Rightarrow$ $gf$ essential: easy.

$2^{\circ}$ $f$, $gf$ essential $\Rightarrow$ $g$ essential.

$g$ is surjective since $gf$ is surjective.

Let $V_{0} \subset V$, then $f^{-1}(V_{0}) \not\subset U$ since $f$ is essential.

Then $g(V_{0}) \not\subset W$ since $gf$ is essential.

$3^{\circ}$ $g$, $gf$ essential $\Rightarrow$ $f$ essential.

$f$ is surjective since $gf$ is surjective and $g$ is essential.

Let $U_{0} \subseteq U \Rightarrow gf(U_{0}) \not\subset M \Rightarrow f(U_{0}) \not\subset V$.

(2) Consider the commutative diagram
\[
\begin{tikzcd}
U \arrow[r, "f"] \arrow[d, "\phi_{1}"'] & V \arrow[d, "\phi_{2}"] \\
U/\operatorname{Rad}U \arrow[r, "\bar{f}"'] & V/\operatorname{Rad}V
\end{tikzcd}
\]
$\phi_{1}$, $\phi_{2}$ essential by Nakayama's Lemma, then $f$ essential $\Leftrightarrow$ $\bar{f}$ essential by (1).

(3) $\bigoplus U_{i}/\operatorname{Rad}\left(\bigoplus V_{i}\right) \longrightarrow \bigoplus V_{i}/\operatorname{Rad}\left(\bigoplus V_{i}\right)$

$\bigoplus (V_{i}/\operatorname{Rad}V_{i}) \rightarrow \bigoplus (V_{i}/\operatorname{Rad}V_{i})$
\end{proof}

\begin{corollary}
If $P$ and $Q$ are projective Noetherian modules over a ring, then $P \simeq Q$ iff $P/\operatorname{Rad}P \simeq Q/\operatorname{Rad}Q$.
\end{corollary}

\begin{corollary}
If $P$ and $Q$ are projective Noetherian modules over a ring, then $P \simeq Q$ iff $P/\operatorname{Rad}P \simeq Q/\operatorname{Rad}Q$.

\textbf{Proof:} $P \rightarrow P/\operatorname{Rad}P$ and $Q \rightarrow Q/\operatorname{Rad}Q$ are essential.

Then $P$ (resp. $Q$) is a projective cover of $P/\operatorname{Rad}P$ (resp. $Q/\operatorname{Rad}Q$).

The rest follows by the uniqueness of projective covers.
\end{corollary}

\begin{definition}
A ring $R$ (with $1$) is called \textbf{semisimple} if ${}_{R}R$ is semisimple as a left module (i.e., it is a direct sum of simple left ideals).
\end{definition}

\begin{theorem}
\textbf{Equivalent definition of semisimple Artinian ring.} TFAE for an Artinian ring $R$:

(1) $R$ is a semisimple Artinian ring.

(2) Every $R$-module is semisimple.

\textbf{Proof:} $(1) \Rightarrow (2)$

Let ${}_{R}R = \sum_{\lambda} I_{\lambda}$ be a sum of simple left $R$-submodules.

Then $V = \sum_{v \in V} \sum_{\lambda} I_{\lambda}v$ and let $f: I_{\lambda} \rightarrow I_{\lambda}v$ be an $R$-homomorphism
\[
i \mapsto iv.
\]
Then $f$ is surjective, hence $I_{\lambda}v \simeq I_{\lambda}/\ker f = 0$ or $I_{\lambda}$.

Then $I_{\lambda}v$ is $0$ or simple.
\end{theorem}

\begin{lemma}
$M \simeq \operatorname{Hom}(R, M)$.

Let $M$ be a left $R$-module. Then $M \simeq \operatorname{Hom}_{R}(R, M)$ as left $R$-modules.
\end{lemma}
\begin{proof}
$\operatorname{Hom}(R, M) \rightarrow M$

\[
f \mapsto f(1)
\]

\[
\begin{aligned}
f & \mapsto f^{\prime}(1) \\
f_{m}: r \mapsto r m & \longleftrightarrow m
\end{aligned}
\]

\textbf{Corollary:} $\operatorname{End}_{R}({}_{R}R) \simeq R^{\mathrm{op}}$

For any ring $R$ (with 1), one has $\operatorname{End}_{R}({}_{R}R) \simeq R^{\mathrm{op}}$ as rings.
\end{proof}

\begin{proof}
$\operatorname{End}_{R}(R) \rightarrow R^{\mathrm{op}}$

\[
\begin{aligned}
& f \mapsto f(1) \\
& f_{1} f_{2}(1)=f_{1}\left(f_{2}(1)\right)=f_{1}\left(f_{2}(1) \cdot 1\right)=f_{2}(1) f_{1}(1)
\end{aligned}
\]

\end{proof}

\begin{lemma}
$M_{n}(R)^{\mathrm{op}} \simeq M_{n}\left(R^{\mathrm{op}}\right)$
\end{lemma}

\begin{proof}
Define $\varphi: M_{n}(R)^{\mathrm{op}} \rightarrow M_{n}\left(R^{\mathrm{op}}\right)$

\[
A \mapsto A^{\top}
\]

Then $\varphi(A * B)=(B A)^{\top}=A^{\top} *^{\prime} B^{\top}=\varphi(A) *^{\prime} \varphi(B)$.

When $R$ is not commutative:

\[
\begin{gathered}
(B A)^{\top} \neq A^{\top} B^{\top} \\
\text{in fact } (B A)^{\top}=A^{\top} * B^{\top}
\end{gathered}
\]

\end{proof}

\begin{theorem}[Wedderburn-Artin]
Let $A$ be a semisimple Artinian ring (with 1). Then $A$ is a direct sum of matrix rings over division rings.

Specifically, if $A \simeq S_{1}^{n_{1}} \oplus \cdots \oplus S_{r}^{n_{r}}$ where $S_{1}, \ldots, S_{r}$ are non-isomorphic simple modules occurring with multiplicities $n_{1}, \ldots, n_{r}$, then
\[
A \simeq M_{n_{1}}\left(D_{1}\right) \oplus \cdots \oplus M_{n_{r}}\left(D_{r}\right)
\]
where $D_{i}=\operatorname{End}_{A}\left(S_{i}\right)^{\mathrm{op}}$.

This is a ring isomorphism. Each $M_{n_{i}}\left(D_{i}\right)$ is a 2-sided ideal.
\end{theorem}

\begin{proof}
By Schur's lemma, $D_{i}$ is a division ring.

\[
\begin{aligned}
& {}_{A}A \simeq S_{1}^{n_{1}} \oplus \cdots \oplus S_{r}^{n_{r}} \\
& \operatorname{End}_{A}({}_{A}A) \simeq \operatorname{End}_{A}\left(S_{1}^{n_{1}}\right) \oplus \cdots \oplus \operatorname{End}_{A}\left(S_{r}^{n_{r}}\right)
\end{aligned}
\]

\[
\begin{aligned}
\operatorname{End}_{A}({}_{A}A) & \simeq \operatorname{End}_{A}\left(S_{1}^{n_{1}}\right) \oplus \cdots \oplus \operatorname{End}_{A}\left(S_{r}^{n_{r}}\right) \\
& \quad \left(\operatorname{Hom}_{A}\left(S_{i}^{n_{i}}, S_{j}^{n_{j}}\right)=0 \text{ if } i \neq j\right) \\
\operatorname{End}_{A}\left(S_{i}^{n_{i}}\right) & \simeq M_{n_{i}}\left(D_{i}\right)^{\mathrm{op}} \simeq M_{n_{i}}\left(D_{i}^{\mathrm{op}}\right)
\end{aligned}
\]

\end{proof}

\begin{definition}
A ring is called \emph{simple} if it has no nontrivial 2-sided ideals.
\end{definition}

\begin{example}
$M_n(D)$ is a simple ring if $D$ is a division ring while
\[
M_n(D) M_n(D) = \left\{\begin{pmatrix} * & 0 & \cdots & 0 \\ \vdots & 1 & \vdots \\ * & 0 & \cdots & 0 \end{pmatrix} \cdot \begin{pmatrix} 0 & * & 0 & 0 \\ 1 & 1 & 1 & 1 \\ 1 & * & 1 \\ 0 & * & 0 & 0 \end{pmatrix}\right\} \oplus \cdots \oplus \begin{pmatrix} 0 & \cdots & * \\ 1 & 1 & * \\ 1 & \vdots & 1 \\ 0 & \cdots & 0 \end{pmatrix}
\]

\end{example}

\paragraph{Simple left $M_n(D)$-module}

\begin{definition}
$R$-algebra

Let $R$ be a commutative ring with $1$ and $A$ be a ring with $1$. We call $A$ an $R$-algebra if $A$ is an $R$-module such that
\[
r(xy) = (rx)y = x(ry) \quad \forall x, y \in A, \quad r \in R.
\]

In particular, $(r \cdot 1)y = r(1 \cdot y) = r(y \cdot 1) = y(r \cdot 1)$.

Equivalently there is a ring homomorphism $R \xrightarrow{\varphi} A$ such that the image of $R$ lies in the center of $A$ and $R \times A \rightarrow A$,
\[
(r, a) \mapsto \varphi(r)a.
\]

If moreover $R$ is a field and $A$ is finite dimensional over $R$, then we say that the algebra is finite dimensional $\Rightarrow$ Noetherian \& Artinian.

\end{definition}

\begin{example}
(1) Group algebra: $R$ commutative ring with $1$; $G$ finite group.
\[
RG = \left\{\sum_{g \in G} r_g g \mid r_g \in R\right\}
\]
free $R$-module with basis $\{g \mid g \in G\}$.

(2) Polynomial ring $F[x]$ over a field.

(3) Matrix algebra over a commutative ring.
\[
\begin{aligned}
\varphi: R &\rightarrow M_n(R) \\
r &\mapsto \begin{pmatrix} r & & \\ & \ddots & \\ & & r \end{pmatrix}
\end{aligned}
\]

Let $A$ be a $k$-algebra, where $k$ is a field. Let $S$ be a simple $A$-module. Then $\operatorname{End}_A(S)$ is a division ring which is a $k$-algebra. If $A$ is finite dimensional over $k$, so is $\operatorname{End}_A(S)$. Furthermore, if $A$ is finite dimensional and $k$ is algebraically closed, then $\operatorname{End}_A(S) \simeq k$. ( $f: S \rightarrow S$ has an eigenvalue, $f - \lambda I: S \rightarrow S$ is singular since $S$ is simple, $f - \lambda I = 0 \Rightarrow f = \lambda I$ )
\end{example}

\paragraph{Lesson 6: Idempotents, Jacobson radical, Artinian ring}

Remark. Let $M$ be an $R$-module, $x \in M$. $\operatorname{Ann}_R(x) = \{r \in R \mid rx = 0\}$ is called the annihilator of $x$ in $R$ $\Rightarrow$ left ideal of $R$. $X \subseteq M$: $\operatorname{Ann}_R(X) = \bigcap_{x \in X} \operatorname{Ann}_R(x)$. $N \subseteq M$ submodule: $\operatorname{Ann}_R(N)$ is a two-sided ideal of $R$.

\begin{lemma}
Let $R$ be a ring with $1$. Let $S$ be a simple left $R$-module. Then $S$ is a quotient of the left regular module ${}_R R$. Let $x \in S$, $x \neq 0$. Then $Rx = S$.
Define $\varphi: R \rightarrow S$ by $r \mapsto rx$. Then $S \simeq R/\operatorname{Ann}_R(x)$.

\end{lemma}

\begin{theorem}
A finite-dimensional semisimple $k$-algebra $A$ over a field $k$ with $k=\bar{k}$ is isomorphic to a direct sum of matrix algebras over $k$.
\[
A \simeq M_{n_1}(k) \oplus \cdots \oplus M_{n_r}(k)
\]
\end{theorem}

\begin{corollary}
If $A$ is a finite-dimensional $K$-algebra, then any simple $A$-module is finite-dimensional as a $k$-vector space.
\end{corollary}

\begin{proposition}

\end{proposition}

\begin{proposition}
Let $A$ be a finite-dimensional semisimple $k$-algebra over a field $k$. In any decomposition $A \simeq S_{1}^{n_{1}} \oplus \cdots \oplus S_{r}^{n_{r}}$ where the $S_i$ are pairwise non-isomorphic simple $A$-modules, we have that $S_{1}, \ldots, S_{r}$ is a complete set of representatives of the isomorphism classes of simple $A$-modules. If, moreover, $k=\bar{k}$, then $n_{i}=\operatorname{dim}_{k} S_{i}$ and $\operatorname{dim}_{k} A=\sum n_{i}^{2}$.
\end{proposition}

\begin{definition}[Idempotent]
Let $R$ be a ring with identity. An idempotent in $R$ is a nonzero element $e$ with $e^{2}=e$.

Then $eRe$ is a ring with identity $e$.

Two idempotents $e_{1}, e_{2}$ are called orthogonal if $e_{1} e_{2}=0=e_{2} e_{1}$. More generally, $n$ idempotents $e_{1}, \ldots, e_{n}$ are called orthogonal if any two of them are orthogonal.

(Note that if $e_{1}, \ldots, e_{n}$ are orthogonal, then $e_{1}+\cdots+e_{n}$ is also an idempotent.)

An idempotent decomposition of an idempotent $e$ is to decompose $e$ into a sum of orthogonal idempotents $e=e_{1}+\cdots+e_{n}$. If moreover $e_{1}, \ldots, e_{n}$ are primitive, then it is called a primitive idempotent decomposition.
\end{definition}

\begin{remark}
Let $e \neq 0$ be an idempotent. Then $1=e+(1-e)$ is an idempotent decomposition. Note that $e$ cannot be written as a sum $e=e_{1}+e_{2}$ where $e_{1}, e_{2}$ are orthogonal idempotents unless one of them is zero.
\begin{itemize}
\item $(1-e)^{2}=1-e$
\item $(1-e)e=0=e(1-e)$
\end{itemize}
\end{remark}

\begin{theorem}
Let $e$ be an idempotent of $R$. If $e=e_{1}+\cdots+e_{n}$ is an idempotent decomposition, then $Re=Re_{1} \oplus \cdots \oplus Re_{n}$ (as left ideals or left $R$-modules).

Conversely, if $Re$ is a direct sum of left ideals of $R$, $Re=I_{1} \oplus \cdots \oplus I_{n}$, then there is an idempotent decomposition $e=e_{1}+\cdots+e_{n}$ with $e_{i} \in I_{i}$ and $I_{i}=Re_{i}$.

That is, the idempotent decompositions of $e$ are in bijection with the direct sum decompositions of the left ideal $Re$.
\end{theorem}
\begin{proof}
$\Rightarrow$ Let $e=e_{1}+\cdots+e_{n}$ be an idempotent decomposition. Then $e_{i}=e_{i}e \Rightarrow Re_{i} \subseteq Re$. $\forall a \in R$, $ae=ae_{1}+\cdots+ae_{n} \in Re_{1}+\cdots+Re_{n}$ thus $Re=Re_{1}+\cdots+Re_{n}$. Let $ae=a_{1}e_{1}+\cdots+a_{n}e_{n}$. Then $ae_{i}=aee_{i}=a_{i}e_{i}$, i.e., the expression is unique. Hence $Re=Re_{1} \oplus \cdots \oplus Re_{n}$.

$\Leftarrow$ Let $Re=I_{1} \oplus \cdots \oplus I_{n}$. Let $e=e_{1}+\cdots+e_{n}$. Then $e_{i}^{2}=e_{i}$ and $e_{i}e_{j}=0=e_{j}e_{i}$ for $i \neq j$.
\end{proof}

\begin{corollary}
Equivalent definitions of primitive idempotent. Let $e$ be an idempotent of $R$. TFAE:
\begin{enumerate}
\item $e$ is primitive.
\item $Re$ is indecomposable as a left $R$-module.
\item $e$ is the unique idempotent in $eRe$.
\end{enumerate}
\end{corollary}

\begin{proof}
(1)$\Leftrightarrow$(2) Immediate from the previous result.

(1)$\Rightarrow$(3) Suppose $e^{\prime} \in eRe$, $e^{\prime} \neq 0$, $e^{\prime 2}=e^{\prime}$, $e^{\prime} \neq e$. Then $e=e^{\prime}+(e-e^{\prime})$ is an idempotent decomposition, contradicting (1).

(3)$\Rightarrow$(1) If $e=e_{1}+e_{2}$ is an idempotent decomposition, then $ee_{1}e=e_{1}$ and $ee_{2}e=e_{2}$, so $e_{1}, e_{2} \in eRe$ are idempotents, contradicting (3).
\end{proof}

\begin{proposition}
Characterization of the direct summands of ${}_{R}R$. A left ideal $I$ of a ring $R$ is a direct summand of ${}_{R}R$ if and only if there is an idempotent $e \in R$ such that
\[
I=Re.
\]
Moreover, if $e \in R$ is an idempotent, then so is $1-e$, and $Re$ and $R(1-e)$ are direct complements of each other. That is,
\[
{}_{R}R=Re \oplus R(1-e).
\]
\end{proposition}

\begin{proof}
We have $\rho: R \rightarrow \operatorname{End}({}_{R}R)$ an isomorphism. Moreover, for the projection $\pi: {}_{R}R \rightarrow I$, $\pi$ is an idempotent in $\operatorname{End}({}_{R}R)$ with $I=\operatorname{im}\pi$. Let $e=\rho^{-1}(\pi)$, then we have $I=Re$.
\end{proof}

\begin{theorem}
Let $e, f$ be idempotents of $R$ and $V$ be an $R$-module.
\begin{enumerate}
\item $\operatorname{Hom}_{R}(Re, V) \simeq eV$. In particular, $\operatorname{Hom}_{R}(Re, Rf) \simeq eRf$ (as abelian groups).
\item $\operatorname{End}_{R}(Re) \simeq (eRe)^{\mathrm{op}}$ as rings.
\end{enumerate}
\end{theorem}

\begin{proof}
(1) Let $\varphi \in \operatorname{Hom}_{R}(Re, V)$, then $\varphi(e)=\varphi(ee)=e\varphi(e)$. Consider the map $g: \operatorname{Hom}_{R}(Re, V) \rightarrow eV$,
\[
\varphi \mapsto \varphi(e).
\]
\end{proof}
$\varphi(ae)=a\varphi(e) \quad \forall a \in R \Rightarrow \varphi$ determined by $\varphi(e)$. Hence $g$ is injective.

Let us define $\psi\colon Re \to V\colon xe \mapsto xeV$. Hence $g$ is injective.

Let $V \in eV$, define $\psi\colon Re \to V\colon xe \mapsto xeV$. $\psi$ is an $R$-homomorphism with $g(\psi)=\mathrm{Id}$. Hence $g$ is surjective.

(2) $V=Re$ for $\sigma, z \in \operatorname{End}_{R}(Re)$.
\[
\sigma(z(e))=\sigma(z(e)e)=z(e)\sigma(e)
\]
So $\operatorname{End}_{R}(Re) \simeq (eRe)^{\mathrm{op}}$.

\begin{theorem}
Let $e, f$ be idempotents of $R$. TFAE:
\begin{enumerate}
\item $Re \simeq Rf$ (as left modules)
\item $eR \simeq fR$ (as right modules)
\item $\exists b \in fRe, a \in eRf$ with $ab=e$ and $ba=f$.
\end{enumerate}
\end{theorem}

\begin{proof}
WLOG, it suffices to prove $(1) \Leftrightarrow (3)$.

$(1)\Rightarrow(3)$: $\quad Re \stackrel{\varphi}{\simeq} Rf$.
Let $a=\varphi(e)=e\varphi(e) \in eRf$.
\[
b=\varphi^{-1}(f)=f\varphi^{-1}(f) \in fRe.
\]
Then $f=\varphi(\varphi^{-1}(f))=\varphi(b)=\varphi(be)=b\varphi(e)=ba$.
\[
e=\varphi^{-1}(\varphi(e))=\varphi^{-1}(a)=\varphi^{-1}(af)=a\varphi^{-1}(f)=ab.
\]

$(3)\Rightarrow(1)$: $eaf=eaba=ea=a$, $fbe=fbab=fb=b$.
So $\varphi_{a}\colon Re \to Rf;\, xe \mapsto xf$,
$\varphi_{b}\colon Rf \to Re;\, xf \mapsto xe$
satisfy $\varphi_{a} \circ \varphi_{b} = \mathrm{id}_{Rf}$, $\varphi_{b} \circ \varphi_{a} = \mathrm{id}_{Re}$.
\end{proof}

\begin{definition}
Jacobson radical of a ring.

Let $R$ be a ring. The (Jacobson) radical $J(R)$ of $R$ is the intersection of all maximal left ideals of $R$.
\end{definition}

\begin{remark}
$J(R)=\operatorname{Rad}({}_{R}R)$.
\end{remark}

\begin{theorem}
Properties of $J(R)$.

The following holds for the radical of a ring $R$.
\begin{enumerate}
\item $J(R)$ is the intersection of all the annihilator ideals of simple (left) $R$-modules. Consequently, $J(R)$ is a $2$-sided ideal of $R$.
\item $a \in J(R) \Leftrightarrow 1-xa$ has a left inverse for all $x \in R$, i.e., $\exists b \in R$ s.t.\ $b(1-xa)=1$.
\item $J(R)$ is the largest one among the $2$-sided ideals $I$ satisfying $a \in I \Leftrightarrow 1-a$ is a unit in $R$.
\item $J(R)$ coincides with the intersection of all maximal right ideals of $R$.
\end{enumerate}
\end{theorem}
\begin{proof}
(1) Let $S$ be a simple $R$-module. Then $S \simeq R / \operatorname{Ann}_{R}(x)$ for any $x \in S$. Thus $\operatorname{Ann}_{R}(x)$ is a maximal left ideal of $R$. In particular, $J(R) \subseteq \operatorname{Ann}_{R}(x)$. Hence $J(R) \subseteq \operatorname{Ann}_{R}(S)$.

Conversely, we want to show that for any $x \in R$ such that $x S = 0$ for all simple modules $S$, then $x$ is in every maximal left ideal. Let $I$ be a maximal left ideal, then $R / I$ is a simple $R$-module. $x \in \operatorname{Ann}_{R}(R / I) \Rightarrow x \in I$.

(2) $\Rightarrow$ Let $a \in J(R)$. Suppose $1 - xa$ has no left inverse in $R$ for some $x \in R$. Then $R(1 - xa) \neq R$. Hence there is a maximal left ideal $M$ of $R$ with $1 - xa \in M$. Since $a \in J(R) \subseteq M$, $xa \in M \Rightarrow 1 \in M$, contradiction.

$\Leftarrow$ If $a \notin J(R)$, then $a \notin M$ for some maximal left ideal $M$ of $R$. Thus $M + Ra = R$. So $b + xa = 1$ for some $b \in M$, $x \in R$. Then $1 - xa = b \in M$ has no left inverse, contradiction.

(3) Let $a \in J(R)$. Then $\exists b$ such that $b(1 - a) = 1$. Now $b = 1 + ab$ has a left inverse, i.e., $\exists c \in R$ such that $cb = 1$. So $c = cb(1 - a) = 1 - a \Rightarrow 1 - a$ is a unit.

Let $I$ be a $2$-sided ideal with $a \in I \Leftrightarrow 1 - a$ is a unit of $R$. If $I \not\subseteq J(R)$ there exists a maximal left ideal $M$ of $R$ with $I \not\subseteq M$, whence $M + I = R$. $\exists b \in M$ such that $b + a = 1 \Rightarrow 1 - a \in M$, then $1 - a$ is not a unit, contradiction.
\end{proof}

\begin{theorem}
If $I$ is a nil left ideal, then $I \subseteq J(R)$.
\end{theorem}

\begin{proof}
Note that for any nilpotent element $a$, say $a^{n} = 0$, then $(1 - a)(1 + a + a^{2} + \cdots + a^{n-1}) = 1$; in particular, $1 - a$ is a unit.

Now, for any $a \in I$, $\forall x \in R$, $xa \in I$. Then $1 - xa$ is a unit $\Rightarrow a \in J(R)$.
\end{proof}

\begin{theorem}
If $A$ is Artinian, then $J(A)$ is nilpotent.
\end{theorem}

\begin{proof}
Consider $J(A) \geq J(A)^{2} \geq \ldots$

$\exists n$ such that $J(A)^{n} = J(A)^{n+1} = \cdots$.

Assume $N = J(A)^{n} \neq 0$. Consider the set of left ideals $I$ of $A$ such that $NI \neq 0$ (nonempty since $N^{2} \neq 0$). Let $I_{0}$ be a minimal member of this set. $\exists a \in I_{0}$ such that $Na \neq 0$ by construction. $Na \subseteq I_{0}$. $N(Na) = N^{2}a = Na \neq 0$. So $Na = I_{0}$ by the minimality of $I_{0}$. In particular, $\exists b \in N$ such that $ba = a$, i.e., $(1 - b)a = 0$. Recall that $\forall b \in J(A)$, $1 - b$ is a unit $\Rightarrow a = 0$, contradiction.
\end{proof}

\begin{corollary}
Let $A$ be an Artinian ring. Then $J(A)$ is the maximal nilpotent left ideal of $A$.
\end{corollary}

\begin{corollary}
Let $A$ be an Artinian ring. Then $J(A)$ is the unique left ideal $U$ of $A$ with the property that $U$ is nilpotent and $A/U$ is semisimple.
\end{corollary}

\begin{theorem}
Let $I$ be a 2-sided ideal of $R$ with $I \subset J(R)$. Then $J(R/I)=J(R)/I$.
\end{theorem}

\begin{proof}
$\subseteq$: Let $a+I \in J(R/I)$. Then for $x+I \in R/I$, $\exists b$ s.t.\ $b(1-xa)+I=1+I$. Assume $b(1-xa)=1-c$ for some $c \in I \subset J(R)$. Then $1-c$ is a unit in $R \Rightarrow$ $\exists d$ s.t.\ $db(1-xa)=1 \Rightarrow db$ is a left inverse of $1-xa$. Hence $a \in J(R)$.

$\supseteq$: Let $a \in J(R)$. Then for any $x \in R$, $1-xa$ has a left inverse in $R$, i.e.\ $\exists b$ s.t.\ $b(1-xa)=1$. In particular, $b(1-xa)+I=1+I$. Thus we have shown for any $x+I \in R/I$, $(1-xa)+I$ has a left inverse, i.e.\ $a+I \in J(R/I)$.
\end{proof}

\begin{theorem}
Let $A$ be an Artinian ring. Then $A$ is semisimple $\Leftrightarrow J(A)=0$.
\end{theorem}

\begin{proof}
Let $M_{\lambda}$, $\lambda \in \Lambda$ be all maximal left ideals of $A$. We consider the set of left ideals that are expressed as a finite intersection of some $M_{\lambda}$. Artinian $\Rightarrow$ This set has a minimal element, say $L=M_{\lambda_{1}} \cap \cdots \cap M_{\lambda_r}$. If $L \neq 0$, since $J(A)=0$, there exists $M_{\lambda}$ with $M_{\lambda} \cap L \subsetneq L$, contradicting the minimality of $L$. So $L=0$.

Consider the map $A \rightarrow A/M_{\lambda_{1}} \oplus \cdots \oplus A/M_{\lambda_r}$. This is a monomorphism. Hence ${}_AA$ is semisimple.

$\Leftarrow$: $A=S_{1} \oplus \cdots \oplus S_{r}$, $S_i$ simple. $J(A) \leq \operatorname{Ann}_R(S_{i})$. Then $J(A)A=0 \Rightarrow J(A)=0$.
\end{proof}

\begin{proposition}
For a ring $R$, TFAE:
\begin{enumerate}
\item[(a)] $R$ has a simple left generator, i.e.\ $\exists$ some simple module $T$ and some $\varphi: T^{n} \rightarrow {}_RR$ for some $n$.
\item[(b)] $R$ is simple and Artinian.
\item[(c)] For some simple ${}_RT$, ${}_RR \simeq T^{(n)}$ for some $n$.
\item[(d)] $R$ is simple and ${}_RR$ is semisimple.
\end{enumerate}
\end{proposition}

\begin{remark}
The four equivalent statements are equal to the according four equivalent statements for ``right'', denoted by $(a')$, $(b')$, $(c')$, $(d')$.
\end{remark}
(a)$\Leftrightarrow$(c)

(c) $\Rightarrow$(d) Let $I$ be a proper ideal of $R$. Then $I$ is contained in a maximal left ideal $L$ of $R$ with $R/L \simeq T$. Since $R \simeq T^{n}$ for some $n$, $\operatorname{Ann}_{R}(T)=0$. Now, $IR \subseteq L$. Then
\[
I \subseteq \operatorname{Ann}_{R}(R/L)=\operatorname{Ann}_{R}(T)=0.
\]
Hence $R$ is simple.

(d)$\Rightarrow$(b)


\paragraph{Lesson 7: Radical Series and Socle Series, Lifting Idempotents}

\begin{proposition}
Let $A$ be an Artinian ring, $U$ a f.g.\ $A$-module, $V \subseteq U$ a submodule.

(1) TFAE:
\begin{enumerate}
\item[(i)] $V=\operatorname{Rad} U$
\item[(ii)] $V$ is the smallest submodule of $U$ with semisimple quotient
\item[(iii)] $J(A)U=V$
\end{enumerate}

(2) TFAE:
\begin{enumerate}
\item[(i)] $V=\operatorname{Soc} U$
\item[(ii)] $V$ is the largest semisimple submodule of $U$
\item[(iii)] $V=\{u \in U \mid J(A) \cdot u=0\}$
\end{enumerate}
\end{proposition}

\begin{proof}
(1) Let $V=J(A)U$. Then $U/V$ is an $A/J(A)$-module, hence semisimple. On the other hand, if $V' \subseteq U$ with semisimple quotient $U/V'$, then $J(A)(U/V')=0$, so $J(A)U \subseteq V'$.

(2)
(iii)$\Rightarrow$(i) $\{u \in U \mid J(A)u=0\}$ is the largest submodule of $U$ annihilated by $J(A)$. It is an $A/J(A)$-module, hence semisimple. It is the largest since each simple module is annihilated by $J(A)$.
\end{proof}

\begin{definition}[Radical series \& Socle series]
Let $U$ be a f.g.\ $A$-module with $A$ Artinian. Define inductively:
\[
\operatorname{Rad}^{n}(U)=\operatorname{Rad}\left(\operatorname{Rad}^{n-1}U\right), \quad
\operatorname{Soc}^{n}(U)/\operatorname{Soc}^{n-1}(U)=\operatorname{Soc}\left(U/\operatorname{Soc}^{n-1}(U)\right).
\]
So $\operatorname{Rad}^{n}(U)=J(A)^{n}U$ and
\[
\operatorname{Soc}^{n}(U)=\left\{u \in U \mid J(A)^{n} u=0\right\}.
\]

Chains of submodules (Loewy series):
\[
\begin{array}{ll}
U=\operatorname{Rad}^{0}(U) \supseteq \operatorname{Rad}(U) \supseteq \operatorname{Rad}^{2}(U) \supseteq \cdots & \rightarrow \text{radical series} \\
0 \subseteq \operatorname{Soc}(U) \subseteq \operatorname{Soc}^{2}(U) \subseteq \cdots & \rightarrow \text{socle series}
\end{array}
\]
\end{definition}

\begin{lemma}
\[
\begin{aligned}
& \operatorname{Rad}^{n}(U \oplus V)=\operatorname{Rad}^{n}(U) \oplus \operatorname{Rad}^{n}(V) \\
& \operatorname{Soc}^{n}(U \oplus V)=\operatorname{Soc}^{n}(U) \oplus \operatorname{Soc}^{n}(V)
\end{aligned}
\]
\end{lemma}

\begin{theorem}[Properties of radical series and socle series]
Let $A$ be an Artinian ring. Let $U$ be a f.g.\ $A$-mod.
\begin{enumerate}
\item The radical series of $U$ is the fastest descending series of submodules of $U$ with semisimple quotient.
\item The socle series of $U$ is the fastest ascending series with semisimple quotient.
\item The two series terminate, and if $m$ and $n$ are the minimal integers with $\operatorname{Rad}^{m} U=0$ and $\operatorname{Soc}^{n} U=U$, then $m=n$.
\end{enumerate}
\end{theorem}

\begin{proof}
(1) Suppose that $\ldots \subseteq U_{2} \subseteq U_{1} \subseteq U_{0}$ is a series of submodules of $U$ with semisimple quotient. We show by induction on $r$ that $\operatorname{Rad}^{r}(U) \subseteq U_{r}$.

$r=0$: clear.

Suppose $r>0$, $\operatorname{Rad}^{r-1}(U) \leqslant U_{r-1}$.
\[
\operatorname{Rad}^{r-1}(U) / \operatorname{Rad}^{r-1}(U) \cap U_{r} \simeq \operatorname{Rad}^{r-1}(U)+U_{r} / U_{r} \leq U_{r-1} / U_{r}
\]
So $\operatorname{Rad}^{r-1}(U) \cap U_{r} \geqslant \operatorname{Rad}\left(\operatorname{Rad}^{r-1}(U)\right)=\operatorname{Rad}^{r}(U)$.

(3) $J(A)^{r}=0$ for some $r$.
\[
\operatorname{Rad}^{r}(U)=J(A)^{r} U=0
\]
\[
\operatorname{Soc}^{r}(U)=\left\{u \in U \mid J(A)^{r} u=0\right\}=U
\]
\end{proof}

\begin{lemma}
\begin{enumerate}
\item Let $A$ be a semisimple ring. Then $A$ is Artinian if and only if it is Noetherian.
\item If $A$ is an Artinian ring, then it is Noetherian.
\end{enumerate}
\end{lemma}

\begin{proof}
(2) See $A$ as a left module ${}_{A}A$. Let $J=J(A)$.
\[
J^{2}=J \cdot J=\operatorname{Rad}\left(J^{2}\right), \quad J^{n+1}=\operatorname{Rad}\left(J^{n}\right)
\]
\[
A \supseteq J \geq J^{2} \geq \cdots \geq J^{n}=0
\]
where $J^{n} / J^{n+1}$ is semisimple.
\end{proof}

\begin{definition}
Lift an idempotent $e$ in $R / I$ to an idempotent $f \in R$.

Let $I$ be a $2$-sided ideal of $R$. Let $f \notin I$ be an idempotent of $R$. $e=f+I$ is an idempotent in $R / I$. We say $f$ lifts $e$. If for every idempotent $e$ of $R / I$, there is an idempotent $f \in R$ lifting $e$, then we can say lift idempotents from $R / I$ to $R$.
\end{definition}

\begin{theorem}
Lift idempotents from a nilpotent ideal.

Let $I$ be a nilpotent ideal of $R$. Then we can lift idempotents.
\end{theorem}

\begin{theorem}
Lift idempotents from a nilpotent ideal.

Let $I$ be a nilpotent ideal of $R$. Then we can lift idempotents from $R / I$ to $R$. Moreover, if $e$ is primitive, so is any lift $f$.
\end{theorem}
\begin{proof}
We define (inductively) idempotents $e_i \in R/I^i$ such that
$e_{i} + I^{i-1}/I^{i} = e_{i-1}$ for each $i$ and $e_{1} = e$.
(Use $R/I^{i} / I^{i-1}/I^{i} \simeq R/I^{i-1}$.)
Suppose $e_{i-1}$ is an idempotent of $R/I^{i-1}$.
Pick any element $a \in R/I^{i}$ with $a + I^{i-1} = e_{i-1}$.
$a^{2} - a \in I^{i-1}/I^{i}$. Since $(I^{i-1})^{2} \subseteq I^{i}$, we have
$\left(a^{2} - a\right)^{2} = 0$ in $R/I^{i}$.
Put $e_{i} = 3a^{2} - 2a^{3}$, $e_{i} + I^{i-1} = e_{i-1}$.
$e_i^{2} - e_i = \left(3a^{2} - 2a^{3}\right)^{2} - 3a^{2} + 2a^{3} = -(3-2a)(1+2a)\left(a^{2}-a\right)^{2} = 0$ in $R/I^{i}$.
If $I^{r} = 0$, let $f = e_r$.
Suppose $e$ is primitive and $f$ has an idempotent decomposition $f = f_{1} + f_{2}$.
Then $e = e_{1} + e_{2}$ with $f_{i} + I = e_{i}$.
\[
\begin{aligned}
& e_{1} e_{2} = e_{2} e_{1} = 0 \\
& e_{1}^{2} = e_{1}, \quad e_{2}^{2} = e_{2}
\end{aligned}
\]
$e$ primitive $\Rightarrow e_{1} = 0$ or $e_{2} = 0$.
Then $f_{1} \in I$ or $f_{2} \in I$. Contradiction, since $I$ has no idempotent.
\end{proof}

\paragraph{Corollary}
Let $I$ be a nilpotent (2-sided) ideal of a ring $R$, and let $1 = e_{1} + \cdots + e_n$ be an idempotent decomposition in $R/I$. Then we have an idempotent decomposition $1 = f_{1} + \cdots + f_{n}$ in $R$ with $e_{i} = f_{i} + I$. If $e_{i}$ are primitive, so is $f_{i}$.

\begin{proof}
Induction on $n$.
$n = 1 \quad \checkmark$
Suppose $n > 1$. Suppose it holds for all $m < n$.
Write $1 = e_{1} + E$ in $R/I$, with $E = e_{2} + \cdots + e_{n} \Rightarrow E$ is an idempotent orthogonal to $e_1$.
Then we lift $e_{1}$ to $f_{1} \in R$, $e_{1} = f_{1} + I$.
By induction hypothesis, we have $1 = f_{1} + F$ in $R$ where $F$ is a lift of $E$.
Now $F$ is the identity of the ring $FRF$, $FIF$ is a nilpotent ideal of $FRF$. The composition homomorphism
$FRF \rightarrow R \rightarrow R/I$
has kernel $FRF \cap I = FIF$.
Indeed, obviously $FIF \subseteq FRF \cap I$. On the other hand,
$\forall x \in FRF \cap I$, $x = FxF \in FIF$.
Then we get an inclusion $FRF/FIF \hookrightarrow R/I$
with image $E(R/I)E$ which has identity $E$.
By induction hypothesis, $F = f_{2} + \cdots + f_{n}$, $f_{i} + I = e_{i}$, $i \geqslant 2$.
Hence $1 = f_{1} + \cdots + f_{n}$.
\end{proof}

\begin{corollary}
Let $f$ be an idempotent in a ring $R$ that has a nilpotent ideal $I$. Then $f$ is primitive in $R$ if and only if $f + I$ is primitive in $R/I$.
\end{corollary}

\begin{lemma}
Existence of primitive decomposition in a semisimple Artinian ring.
\end{lemma}
Let $A$ be a semisimple Artinian ring.
\begin{enumerate}
\item There is a primitive idempotent decomposition
$1 = e_{1} + \cdots + e_{n}$. This corresponds to ${}_{A}A = Ae_{1} \oplus \cdots \oplus Ae_{n}$ where $Ae_{i}$ is simple.
\item For any simple $A$-module $S$, there is a primitive idempotent $e$ with $S \simeq Ae$.
\item For any simple $A$-module $S$, there exists $1 \leq i \leq n$ such that $e_{i}S \neq 0$ (where $e_{i}$ is from (1)), and if $T$ is a simple $A$-module with $T \not\simeq S$, then $e_{i}T = 0$.
\end{enumerate}

\begin{proof}
(3) $S = 1 \cdot S = (e_{1} + \cdots + e_{n})S \Rightarrow \exists i$ s.t.\ $e_{i}S \neq 0$.
For any $T$ simple, $T \simeq Ae$. We have $e_{i}Ae \simeq \operatorname{Hom}_{A}(Ae, Ae_{i}) = 0$ if $Ae \not\simeq Ae_{i}$.
\end{proof}

\begin{theorem}
From semisimple Artinian to Artinian.

Let $A$ be an Artinian ring and $S$ a simple $A$-module.
\begin{enumerate}
\item There is an indecomposable projective $A$-module $P_{S}$ with
$P_{S}/\operatorname{Rad}(P_{S}) \simeq S$ of the form $P_{S} \simeq Af$ where $f$ is a primitive idempotent in $A$.
\item The idempotent $f$ has the property that $fS \neq 0$ and if $T$ is any simple module with $T \not\simeq S$, then $fT = 0$.
\item $P_{S}$ is the projective cover of $S$; it is uniquely determined up to an isomorphism by $S$.
\end{enumerate}
\end{theorem}

\begin{proof}
(1) Let $e \in A/J(A)$ be any primitive idempotent with $eS \neq 0$. Let $f$ be any lift of $e$ to $A$. Then $e$ is primitive.
$fS = eS \neq 0$. $fT = eT = 0$ for any simple module $T \not\simeq S$.
Set $P_{S} = Af$: an indecomposable projective module.
$Af \to (A/J(A))(f + J(A)) = S$ is an $A$-module homomorphism
$af \mapsto (a + J(A))(f + J(A))$ with kernel $J(A)Af = \operatorname{Rad}Af$.
Hence $P_{S}/\operatorname{Rad}P_{S} \simeq S$.
\end{proof}

\begin{theorem}
Let $A$ be an Artinian ring. Up to isomorphism, the indecomposable projective $A$-modules are exactly the modules $P_{S}$ that are the projective covers of simple modules, and $P_{S} \simeq P_{T} \Leftrightarrow S \simeq T$.
Each $P_{S}$ appears as a direct summand of ${}_{A}A$ (left regular module), with multiplicity equal to the multiplicity of $S$ as a summand of $A/J(A)$, i.e.,
\[
{}_{A}A \simeq \bigoplus_{S} (P_{S})^{n_{S}}
\]
where $S$ runs through simple $A$-modules up to isomorphism and $n_{S} = \operatorname{dim}_{D}S$ with $D = \operatorname{End}_{A}(S)$.
\end{theorem}

\begin{proof}
(We only prove for finitely generated indecomposable projective modules, but true for all.)
\end{proof}
Let $P$ be an indecomposable projective module and write
\[
P / \operatorname{Rad} P \simeq S_{1} \oplus \cdots \oplus S_{n}.
\]
Then $P \rightarrow S_{1} \oplus \cdots \oplus S_{n}$ is a projective cover. Now $P_{S_{1}} \oplus \cdots \oplus P_{S_{n}} \rightarrow S_{1} \oplus \cdots \oplus S_{n}$ is also a projective cover, and by uniqueness of projective covers we have
\[
P \simeq P_{S_{1}} \oplus \cdots \oplus P_{S_{n}}.
\]
Since $P$ is indecomposable we have $n=1$ and $P \simeq P_{S_{1}}$.

Suppose that each simple $A$-module $S$ occurs with multiplicity.

Let $A$ be an Artinian ring, $V$ a f.g.\ $A$-mod. Then $V$ has a composition series (both Artinian and Noetherian). For any simple $A$-module $S$, the number of the composition factors isomorphic to $S$ is called the multiplicity of $S$ in $V$.

\begin{theorem}
Let $A$ be an Artinian ring, $f \in A$ be a primitive idempotent and $e=f+J(A) \in A / J(A)=\bar{A}$ a primitive idempotent. Then the multiplicity of $\bar{A} e$ in any f.g.\ $A$-module $V$ is equal to the composition length of $eV$ as an $eAe$-module.
\end{theorem}

\begin{proof}
Let $V=V_{0} \supseteq V_{1} \supseteq \cdots \supseteq V_{n}=0$ be a composition series of $V$. Consider $fAf$-modules:
\[
\begin{aligned}
& f V=f V_{0} \supseteq f V_{1} \supseteq \cdots \supseteq f V_{n}=0, \\
& V_{i} / V_{i+1} \simeq \bar{A} e \Leftrightarrow f\left(V_{i} / V_{i+1}\right) \neq 0.
\end{aligned}
\]
In this case, $e \bar{A} e \simeq f\left(V_{i} / V_{i+1}\right)=f V_{i}+V_{i+1} / V_{i+1} \simeq f V_{i} / f V_{i} \cap V_{i+1}=f V_{i} / f V_{i+1}$ which is also simple as $fAf$-module.
\end{proof}

\paragraph{Lesson 8: Projective morphism, integral closed, localization and fractional modules}

2025年4月13日 星期日

\begin{definition}[Cartan matrix]
Let $A$ be an Artinian ring. Let $1=e_{1}+\cdots+e_{n}$ be a primitive idempotent decomposition in $A / J(A)$. Let $1=f_{1}+\cdots+f_{n}$ be the according primitive idempotent decomposition in $A$ with $f_{i}+J(A)=e_{i}$. We denote by $C_{i j}$ the multiplicity of $(A / J(A)) e_{i}$ in $A f_{j}$, which is called the Cartan invariant in $A$. The $n \times n$ matrix $(C_{i j})_{n \times n}$ is called the Cartan matrix of $A$.

Note that we have
\[
C_{i j} \neq 0 \Leftrightarrow f_{i} A f_{j} \neq 0.
\]

For convenience, if $S, T$ are two simple $A$-modules, we denote by $C_{S T}$ the corresponding Cartan invariant, i.e., the multiplicity of $S$ in $P_{T}$.
\end{definition}

\begin{proposition}
Let $A$ be a finite-dimensional $k$-algebra over a field $k$ and $S$ a simple $A$-module with projective cover $P_{S}$. Let $U$ be a finite-dimensional $A$-module.

(1) If $T$ is a simple $A$-module, then
\[
\dim_{k} \operatorname{Hom}_{A}\left(P_{S}, T\right)= \begin{cases}\dim_{k} \operatorname{End}_{A}(S) & \text { if } S \simeq T \\ 0 & \text { otherwise }\end{cases}
\]

(2) The multiplicity of $S$ in $U$ is $\dim_{k}\operatorname{Hom}_{A}\left(P_{S}, U\right) / \dim_{k} \operatorname{End}_{A}(S)$.

(3) If $e \in A$ is an idempotent, then $\dim_{k} \operatorname{Hom}_{A}(Ae, U)=\dim_{k} eU$.

(4) Let $e_{S}$ and $e_{T}$ be idempotents so that $P_{S}=A e_{S}$, $P_{T}=A e_{T}$ are projective covers of $S$ and $T$. Then
\[
C_{ST}=\dim_{k} \operatorname{Hom}_{A}\left(P_{S}, P_{T}\right) / \dim_{k} \operatorname{End}_{A}(S)=\dim_{k} e_{S}Ae_{T} / \dim_{k} \operatorname{End}_{A}(S).
\]
If moreover $k=\bar{k}$, then
\[
C_{ST}=\dim_{k} \operatorname{Hom}_{A}\left(P_{S}, P_{T}\right)=\dim_{k} e_{S} A e_{T}.
\]

\end{proposition}

\begin{proof}
(1) If $P_{S} \rightarrow T$ is nonzero, the kernel must be a maximal submodule containing $\operatorname{Rad}\left(P_{S}\right)$. But $P_{S} / \operatorname{Rad}\left(P_{S}\right) \simeq S \Leftrightarrow S \simeq T$.

\end{proof}

\paragraph{Now let $R$ be a commutative ring with 1.}

For $R$-modules $X, Y$ and $M$, we have an $R$-morphism
\[
\operatorname{Hom}_{R}(X, M) \otimes_{R} \operatorname{Hom}_{R}(M, Y) \rightarrow \operatorname{Hom}_{R}(X, Y)
\]
for $R$-modules $X, Y$ and $M$, we have an $R$-morphism
\[
\begin{aligned}
\operatorname{Hom}_{R}(X, M) \otimes_{R} \operatorname{Hom}_{R}(M, Y) & \rightarrow \operatorname{Hom}_{R}(X, Y) \\
\varphi \otimes \psi & \longmapsto \psi \circ \varphi
\end{aligned}
\]

\begin{proposition}
The image of this morphism is comprised of those morphisms from $X$ to $Y$ which factor through some finite power of $M$.
\[
X \rightarrow M^{n} \rightarrow Y
\]
\end{proposition}

\begin{proof}
An element of the image is $\sum_{i \in I} \psi_{i} \circ \varphi_{i}$ for some finite set $I$.

This induces a morphism $\varphi\colon X \rightarrow M^{|I|}$; $x \mapsto \left(\varphi_{i}(x)\right)_{i \in I}$ and $\psi\colon M^{|I|} \rightarrow Y$; $\left(m_{i}\right)_{i \in I} \mapsto \sum_{i \in I} \psi_{i}\left(m_{i}\right)$.

\end{proof}

\begin{definition}
$\operatorname{Hom}_{R}^{f}(X, Y)$ denotes the set of morphisms from $X$ to $Y$ which factor through a finite power of $M$.
\end{definition}
Denote $X^{*}=\operatorname{Hom}_{R}(X, R)$ dual. Taking $M:=R$, we get a morphism the image factors through $f\text{-}g$. $\to$ projective modules (direct summand of a free module).

\[
Z_{X, Y}\colon X^{*} \otimes_{R} Y \rightarrow \operatorname{Hom}_{R}(X, Y) \quad\left(\operatorname{Hom}_{R}(R, Y) \cong Y\right)
\]

Denote by $\operatorname{Hom}_{R}^{pr}(X, Y)$ the image of $Z_{X, Y}$. $f \mapsto f(1)$ The elements of $\operatorname{Hom}_{R}^{pr}(X, Y)$ are called the projective morphisms from $X$ to $Y$.


\begin{lemma}
For any $R$-module $M$, $\operatorname{Hom}_{R}^{Pr}(M, M)$ is a two-sided ideal of the ring $\operatorname{End}_{R}(M)=\operatorname{Hom}_{R}(M, M)$.
\end{lemma}

\begin{proof}
Easy.
\end{proof}

\begin{theorem}
Let $P$ be an $R$-module. TFAE:
\begin{enumerate}
\item[(1)] $P$ is a $f\text{-}g$. projective module.
\item[(2)] $\operatorname{Hom}_{R}^{Pr}(P, P)=\operatorname{End}_{R}(P)$.
\item[(3)] For any $R$-module $X$, the morphism $Z_{X, P}\colon X^{*} \otimes P \rightarrow \operatorname{Hom}_{R}(X, P)$ is an isomorphism.
\item[(4)] For any $R$-module $X$, the morphism $Z_{P, X}\colon P^{*} \otimes_{R} X \rightarrow \operatorname{Hom}_{R}(P, X)$ is an isomorphism.
\item[(5)] The morphism $Z_{P, P}\colon P^{*} \otimes_{R} P \rightarrow \operatorname{End}_{R}(P)$ is an isomorphism.
\end{enumerate}
\end{theorem}

\begin{proof}
(1)$\Rightarrow$(2), (3)$\Rightarrow$(5), (4)$\Rightarrow$(5) $\checkmark$

(1)$\Rightarrow$(2): Since $\operatorname{Hom}_{R}^{Pr}(P, P)$ is a two-sided ideal, it suffices to show $\operatorname{id} \in \operatorname{Hom}_{R}^{Pr}(P, P)$, which is obvious since $P$ is a direct summand of $R^{n}$ for some $n$.

(2)$\Rightarrow$(3): Assume $\operatorname{id}_{P}=Z_{P, P}\left(\sum_{i \in I} \varphi_{i} \otimes_{R} m_{i}\right)$, i.e., $\forall v \in P, v=\sum_{i} \varphi_{i}(v) m_{i}$.

Consider the morphism $\theta\colon \operatorname{Hom}_{R}(X, P) \rightarrow X^{*} \otimes P$
\[
\alpha \mapsto \sum_{i} \varphi_{i} \alpha \otimes m_{i}
\]

Check: $\theta$ is the inverse of $Z_{X, P}$.

$\operatorname{Hom}_{R}(X, P) \xrightarrow{\theta} X^{*} \otimes_{R} P \xrightarrow{Z_{X, P}} \operatorname{Hom}_{R}(X, P)$
\[
\alpha \longmapsto \sum_{i} \varphi_{i} \alpha \otimes m_{i} \mapsto \left(x \mapsto \sum_{i} \varphi_{i}(\alpha(x)) m_{i}\right)
\]

$X^{*} \otimes_{R} P \xrightarrow{Z_{X, P}} \operatorname{Hom}_{R}(X, P) \xrightarrow{\theta} X^{*} \otimes_{R} P$
\[
\psi \otimes v \longmapsto Z_{X, P}(\psi \otimes v) \mapsto \sum_{i} \varphi_{i} Z_{X, P}(\psi \otimes v) \otimes m_{i}
\]
where $Z_{X, P}(\psi \otimes v)(x)=\psi(x) v$, and $\varphi_{i}(Z_{X, P}(\psi \otimes v)(x))=\varphi_{i}(\psi(x) v)=\varphi_{i}(\psi(x) v)$.
\end{proof}
\[
=\varphi_{i}(v) \psi(x)
\]

Hence $\varphi_{i} \circ \tau_{x} \circ p(\psi \otimes v)=\varphi_{i}(v) \psi$. Then $\psi \otimes V=\psi \otimes\left(\sum_{i} \varphi_{i}(V) m_{i}\right)=\sum_{i} \psi \otimes \varphi_{i}(V) m_{i}$

\[
\begin{aligned}
& =\sum \varphi_{i}(v) \psi \otimes m_{i} \\
& =\sum \varphi_{i} \circ \tau_{x} \circ p(\psi \otimes v) \otimes m_{i}
\end{aligned}
\]

(5)$\Rightarrow$(1) Id $p$ is projective. $p \xrightarrow{\mathrm{Id}_p} p$

\[
\begin{array}{c}
p \\
f \searrow \quad \swarrow \hat{g} \\
R^{n}
\end{array}
\]

We have $\mathrm{Id}_p = gf \Rightarrow P$ is a direct summand of $R^{n}$.


\paragraph{Integrality}

Let $S$ be an integral domain with $1$.

\begin{definition}[integral element over a subring $R$]
Let $R$ be a subring of $S$ containing $1_S$. Let $x \in S$. We say $\alpha$ is an integral element over $R$ if there exists a monic polynomial $f(x) \in R[x]$ such that $f(\alpha)=0$.
\end{definition}

\begin{proposition}[equivalent definition of an integral element]
Let $S$ be an integral domain with $1$. Let $R \subseteq S$ be a subring with $1_S$. Let $\alpha \in S$. TFAE:
\begin{enumerate}
\item $\alpha$ is an integral element over $R$.
\item $R[\alpha]$ is a finitely generated $R$-module.
\item There exists a subring $T$ of $S$ with $R \subseteq T$ such that $T$ is finitely generated as an $R$-module and $\alpha \in T$.
\end{enumerate}
\end{proposition}

\begin{proof}
(1)$\Rightarrow$(2): Obvious.

(2)$\Rightarrow$(3): Take $T=R[\alpha]$.

(3)$\Rightarrow$(1): Let $\varepsilon_{1}, \ldots, \varepsilon_{n}$ be a set of generators of $T$ as an $R$-module. Write
\[
\alpha \varepsilon_{i}=a_{i1} \varepsilon_{1}+\cdots+a_{in} \varepsilon_{n}, \quad a_{ij} \in R.
\]
Then
\[
-a_{i1} \varepsilon_{1}+\cdots+(\alpha-a_{ii}) \varepsilon_{i}+\cdots+(-a_{in}) \varepsilon_{n}=0 \quad \forall i.
\]
\[
\left(\varepsilon_{1}, \ldots, \varepsilon_{n}\right)\begin{pmatrix}
\alpha-a_{11} & -a_{21} & \cdots & -a_{n1} \\
-a_{12} & \alpha-a_{22} & & \vdots \\
\vdots & & \ddots & \\
-a_{1n} & \cdots & \cdots & \alpha-a_{nn}
\end{pmatrix}=\left(0, \ldots, 0\right)
\]
Multiplying by $A^{*}$ gives
\[
\left(\varepsilon_{1}, \ldots, \varepsilon_{n}\right)\begin{pmatrix}
\det A & & \\
& \ddots & \\
& & \det A
\end{pmatrix}=(0, \ldots, 0)
\]
\[
\Rightarrow \sum_{i} \varepsilon_{i} \det A = 0
\]
\[
\Rightarrow (\det A) T=0
\]
\[
\Rightarrow \det A=0.
\]
\end{proof}

\begin{lemma}
If $M_{1}, M_{2}$ are finitely generated $R$-modules, then $R[M_{1}, M_{2}]$ is finitely generated.
\end{lemma}
Let $S, R$ be as above. $R \subseteq M_{i} \subseteq S$ where $M_{i}$, $i=1,2$ are subrings of $S$. If $M_{1}, M_{2}$ are f.g.\ $R$-modules, then the ring $R[M_{1}, M_{2}]$ is f.g.\ as an $R$-module.

\textit{Pf:} It suffices to show if $R[\alpha], R[\beta]$ is f.g.\ then $R[\alpha, \beta]$ is f.g.

Indeed, if $R[\alpha]$ is generated by $1, \alpha, \ldots, \alpha^{n}$, $R[\beta]$ is generated by $1, \beta, \ldots, \beta^{m}$, then $R[\alpha, \beta]$ is generated by $\alpha^{i} \beta^{j}$ ($0 \leqslant i \leqslant n$, $0 \leqslant j \leqslant m$).


\begin{theorem}
Integral elements form a subring.

Let $S, R$ be as above. The integral elements of $S$ over $R$ form a subring of $S$ containing $R$.
\end{theorem}

\begin{proof}
Let $\alpha, \beta$ be integral elements $\Rightarrow R[\alpha], R[\beta]$ f.g.

$\Rightarrow R[\alpha, \beta]$ is f.g.\ by lemma.

$\Rightarrow \alpha - \beta, \alpha \beta \in R[\alpha, \beta]$ are integral elements.
\end{proof}

\begin{definition}
Integral closure.
\end{definition}

\begin{definition}
Integral closed.

Let $S, R$ be as above.
\begin{enumerate}
\item The subring of $S$ consisting of all integral elements of $S$ over $R$ is called the \textbf{integral closure} of $R$ in $S$.
\item $R$ is said to be \textbf{integrally closed} in $S$ if every integral element of $S$ over $R$ is contained in $R$.
\item We say an integral domain $R$ is \textbf{integrally closed} if it is integrally closed in its fraction field.
\item $\mathbb{Z} \subseteq \mathbb{C}$. If $z \in \mathbb{C}$ is integral over $\mathbb{Z}$, we call $z$ an \textbf{algebraic integer}. All algebraic integers form a ring.
\end{enumerate}
\end{definition}

\begin{definition}
Torsion and torsion free elements.

$R$: ring. $M$: $R$-module. $m \in M$ is said to be
\begin{itemize}
\item a \textbf{torsion element} if $\operatorname{Ann}_{R}(m) \neq 0$.
\item a \textbf{torsion free element} if $\operatorname{Ann}_{R}(m) = 0$.
\end{itemize}
\end{definition}

\begin{example}
\begin{enumerate}
\item $\mathbb{Z}$-module $\mathbb{Q}$ has no torsion element but $0$.
\item $R$ commutative. $R$ as an $R$-module has no torsion element but $0$ $\Leftrightarrow R$ is an integral domain.
\item If $I$ is a proper ideal, all elements of $R$-module $R/I$ are torsion.
\item $R/I$ has no torsion element but $0$ $\Leftrightarrow I$ is prime.
\end{enumerate}
\end{example}

\begin{lemma}
Torsion module.

If $R$ is an integral domain, the subset of $M$
\[
\operatorname{Tor}(M) = \left\{m \in M \mid \operatorname{Ann}_{R}(m) \neq 0\right\}
\]
is a submodule of $M$, called the \textbf{torsion submodule} of $M$.

We say $M$ is a \textbf{torsion module} if $M = \operatorname{Tor}(M)$.

$M$ is \textbf{torsion free} if $\operatorname{Tor}(M) = 0$.
\end{lemma}

\begin{example}
(1) $\mathbb{Z}$-modules: $\mathbb{Z}, \mathbb{Q}, \mathbb{Z} / 2 \mathbb{Z}$.
\begin{itemize}
\item $\operatorname{tor}(\mathbb{Q})=0$
\item $\operatorname{tor}(\mathbb{Z} / 2 \mathbb{Z})=2 \mathbb{Z}$
\end{itemize}

(2) $R$: integral domain, $M$ an $R$-module.
$\operatorname{tor}(M^{*})=0$ where $M^{*}=\operatorname{Hom}_{R}(M, R)$.

(3) $R=$ integral domain. Free modules are torsion free.
\end{example}

\begin{proposition}
Let $R$ be an integral domain. Any finitely generated torsion free $R$-module $M$ is isomorphic to a submodule of a free $R$-module of finite rank.
\end{proposition}

\begin{proof}
Suppose $M$ is generated by nonzero elements $x_{1}, \ldots, x_{n}$.

Up to reordering, assume $x_{1}, \ldots, x_{r}$ is a maximal linearly independent subset.

If $r=n$, $M$ is free.

Assume $r<n$. Consider $M_{0}=R x_{1}+\cdots+R x_{r}=R x_{1} \oplus \cdots \oplus R x_{r}$.

For each $j \geqslant r+1$, $\exists \lambda_{j} \in R$ s.t.\ $\lambda_{j} x_{j} \in M_{0}$. Let $\lambda=\lambda_{r+1} \cdots \lambda_{n}$.

$\Rightarrow \lambda M \leq M_{0}$. Consider $M \rightarrow M_{0}$ injective:
\[
m \mapsto \lambda m
\]
\end{proof}

\paragraph{Localization}

\begin{definition}[$S^{-1} R$]
Let $R$ be an integral domain with fraction field $F$.

$S$: nonempty multiplicative closed subset of $R \setminus \{0\}$. Then
\[
S^{-1} R=R\left[S^{-1}\right]=\left\{\left.\frac{\lambda}{s} \right\rvert\, s \in S, \lambda \in R\right\}
\]
is a subring of $F$ containing $R$.

$i: R \rightarrow S^{-1} R$ the natural injection.
\end{definition}

\begin{definition}[$S^{-1} M$]
Let $M$ be an $R$-module. $S^{-1} M$ is the quotient of the set $M \times S$ by the equivalence relation
\[
(m_{1}, s_{1}) \sim (m_{2}, s_{2}) \Leftrightarrow \exists s \in S \text{ s.t.\ } s(s_{2} m_{1}-s_{1} m_{2})=0.
\]
It is an $S^{-1} R$-module with
\[
\left(\frac{\lambda}{s}\right)\left(m, s'\right)=\left(\lambda m, s s'\right).
\]
For $s \in S$, $m \in M$, denote by $\frac{m}{s}$ the equivalence class of $(m, s)$.
\end{definition}

\begin{lemma}
$S^{-1} R \otimes_{R} M \simeq S^{-1} M$

The map $f\colon \frac{\lambda}{s} \otimes m \mapsto \frac{\lambda m}{s}$ defines an isomorphism of $S^{-1} R$-modules between $S^{-1} R \otimes_{R} M$ and $S^{-1} M$.
\end{lemma}

\begin{proof}
It is easy to see $f$ is a surjective $S^{-1} R$-morphism.
\end{proof}
It suffices to show $f$ is injective.

First, show any element of $S^{-1} R \otimes M$ can be written as $\frac{1}{s} \otimes m$ for some $s \in S, m \in M$.

Let $\sum_{i} \frac{\lambda_{i}}{s_{i}} \otimes m_{i} \in S^{-1} R \otimes M$. Let $s = \prod_{i} s_{i}$. For all $i$, let $s_{i}' = \prod_{j \neq i} s_{j}$. Then
\[
\sum_{i} \frac{\lambda_{i}}{s_{i}} \otimes m_{i} = \sum_{i} \frac{\lambda_{i} s_{i}'}{s} \otimes m_{i} = \sum_{i} \frac{1}{s} \otimes \lambda_{i} s_{i}' m_{i} = \frac{1}{s} \otimes \sum_{i} \lambda_{i} s_{i}' m_{i}.
\]

If $f\left(\frac{1}{s} \otimes m\right) = 0$, then $\frac{m}{s} = 0$, so $\exists t \in S$ such that $tm = 0$. Then $\frac{1}{s} \otimes m = \frac{t}{st} \otimes m = \frac{1}{st} \otimes tm = 0$.


\begin{corollary}
$F \otimes M$

Set $FM = F \otimes M = S^{-1} M$ where $S = R \setminus \{0\}$.

The kernel of $R \to FM$, $m \mapsto 1 \otimes m$ is $\operatorname{Tor}(M)$.

\begin{enumerate}
\item $F \otimes \operatorname{Tor}(M) = 0$.
\item TFAE:
\begin{enumerate}
\item[(i)] The natural $R$-morphism $M \to FM$, $m \mapsto 1 \otimes m$ is injective;
\item[(ii)] $M$ is torsion free.
\end{enumerate}
\item If $M$ is a f.g.\ projective $R$-module, then $M \to FM$ is injective.
\end{enumerate}
\end{corollary}

\begin{definition}
rank of a f.g.\ torsion free module

The rank of a f.g.\ torsion free $R$-module is the dimension of the $F$-space $FM$.
\end{definition}

\begin{remark}
\begin{enumerate}
\item A f.g.\ torsion free $R$-module $M$ of rank one is isomorphic to a f.g.\ $R$-submodule of $F$.

Indeed, if $FM = Fe$, then $I = \{\lambda \in F \mid \lambda e \in M\}$ is an $R$-submodule of $F$, and we have $M = Ie$. Embed $M$ into $FM$; $M$ is f.g.\ $\stackrel{\text{iso}}{\Rightarrow}$ $I$ is f.g.

\item Let $M$ be a f.g.\ torsion free $R$-module. The map $M \to FM$, $m \mapsto 1 \otimes m$ is injective. So $M$ can be viewed as an $R$-submodule of $FM$. (Note: $M$ f.g.\ $\Rightarrow$ $FM$ f.d.)

Given an $F$-basis $(e_{i})_{1 \leq i \leq r}$ of $FM$, assume $M$ is generated by $m_{1}, \ldots, m_{n}$. Then
\[
m_{i} = \frac{s_{1}}{t_{1}} e_{1} + \cdots + \frac{s_{r}}{t_{r}} e_{r}
\]
for some $s_{i} \in R$, $t_{i} \in R \setminus \{0\}$ (so $\frac{s_{i}}{t_{i}} \in F$).

$\exists \lambda \in R \setminus \{0\}$ s.t.\ $\lambda M \subseteq Re_{1} \oplus \cdots \oplus Re_{r}$.

Additionally, $M$ generates $FM$ as an $F$-vector space.
\end{enumerate}
\end{remark}

\begin{definition}
fractional $R$-module
Let $V$ be an $F$-vector space of dimension $r$. A fractional $R$-module $M$ in $V$ is an $R$-submodule of $V$ such that
\begin{enumerate}
\item $M$ generates $V$ as a vector space.
\item given any basis $(e_{i})_{1 \leq i \leq r}$ of $V$, there exists $\lambda \in R$, $\lambda \neq 0$, such that $\lambda M \subseteq \bigoplus_{i=1}^{r} Re_{i}$.
\end{enumerate}
In the case $V = F$, fractional modules are called fractional ideals of $R$. It is an $R$-submodule $I$ of $F$ such that $\exists \lambda \in R$, $\lambda \neq 0$, with $\lambda I \subseteq R$ (hence $\lambda I$ is an ideal of $R$).

\end{definition}

\begin{example}
\begin{enumerate}
\item Every finitely generated $R$-submodule $M$ of $F$ is a fractional ideal (torsion free; $R$-submodule of $F \Rightarrow FM$ is of dimension 1). Moreover if $R$ is Noetherian, then these are all fractional ideals.
\item If $I$ is a fractional ideal, $\{\lambda \in F \mid \lambda I \leq R\}$ is again a fractional ideal (see below).
\end{enumerate}
\end{example}

\begin{lemma}[Properties of fractional modules]
Let $M$ be a finitely generated torsion-free $R$-module, let $V = FM$. (So $M$ is identified with an $R$-submodule of $FM$.)
\begin{enumerate}
\item If $\lambda M \subset \bigoplus_{i=1}^{r} Re_{i}$ for some basis $(e_{i})_{1 \leq i \leq r}$ of $V$ and some $\lambda \in R$, $\lambda \neq 0$, then whenever $(f_{i})_{1 \leq i \leq r}$ is a basis of $V$ there exists $\mu \in R$, $\mu \neq 0$, such that $\mu M \subset \bigoplus_{i=1}^{r} Rf_{i}$.
\item Condition (1) in the above definition of fractional modules can be replaced by
\begin{enumerate}
\item[(1')] given a basis $(e_{i})_{1 \leq i \leq r}$ of $V$, there exists $\mu \in R$, $\mu \neq 0$, such that $\bigoplus_{i=1}^{r} Re_{i} \subset M$.
\end{enumerate}
\item If $M_{1}$ and $M_{2}$ are fractional modules in $V$, then $M_{1} + M_{2}$ and $M_{1} \cap M_{2}$ are also fractional modules.
\end{enumerate}
\end{lemma}

\begin{proof}
(2) is obvious. (3) obvious from (2).

(1) We have $(e_{1}, \ldots, e_{r}) = (f_{1}, \ldots, f_{r})A$ for some $A = (a_{ij})$ with non-zero determinant. Then
\[
e_{i} = \sum_{j=1}^{r} a_{ji} f_{j} \text{ where } a_{ij} = \frac{b_{ij}}{c_{ij}} \text{ for some } b_{ij} \in R, c_{ij} \in R \setminus \{0\}.
\]
Then set $C = \prod_{i,j}$. We have $c e_{i}^{T_{1},j} \in \bigoplus_{j=1}^{r} R f_{j}$. Since $\lambda M C \bigoplus_{i=1} R e_{i}$, we have $G M \subseteq C \bigoplus_{i=1}^{\infty} R e_{i} C \bigoplus_{j=1}^{r} R f_{j}$.
\end{proof}

\begin{corollary}
connecting fractional ideals with fractional modules

If $(e_{i})_{1 \leq i \leq r}$ is a basis of $V$ and $(\alpha_{i})_{1 \leq i \leq r}$ is a family of fractional ideals then
\[
M = \bigoplus_{i=1}^{r} \alpha_{i} e_{i}
\]
is a fractional module.

$M = \bigoplus \alpha_{i} e_{i}$ is a fractional module.
\end{corollary}

\begin{lemma}
constructing fractional ideals

Let $I_{1}, I_{2}$ be fractional ideals.
\begin{enumerate}
\item $I_{1} + I_{2}$ (the $R$-submodule of $F$ generated by $I_{1} \cup I_{2}$)
\item $I_{1} \cap I_{2}$
\item $I_{1} I_{2}$ the $R$-submodule of $F$ generated by all $x_{1} x_{2}$, $x_{1} \in I_{1}$, $x_{2} \in I_{2}$
\end{enumerate}
are all fractional ideals.
\end{lemma}

\begin{proposition}
characterization of $I^{*}$

(1) For any fractional ideal, there is a natural identification
\[
I^{*} = \{x \in F \mid x I \subset R\}
\]
which is again a fractional ideal.
$\left(I^{*} = \operatorname{Hom}_{R}(I, R)\right)$

(2) Let $V$ be an $r$-dimensional $F$-space with a basis $(e_{i})_{1 \leq i \leq r}$, and let $(I_{i})_{1 \leq i \leq r}$ be a family of fractional ideals.
Set $M = \bigoplus_{i=1}^{r} I_{i} e_{i}$.
\begin{enumerate}
\item[(a)] The duality between $V$ and $V^{*}$ induces
\[
M^{*} = \bigoplus_{i=1}^{r} I_{i}^{*} e_{i}^{*}
\]
\item[(b)] The fractional module $M$ is a finitely generated projective $R$-module iff for all $i$ ($1 \leq i \leq r$), we have $I_{i}^{*} I_{i} = R$.
In particular a fractional ideal $I$ is a finitely generated projective $R$-module iff $I = R$.
\end{enumerate}
\end{proposition}

\begin{proof}
(1) Let $0 \neq \lambda \in R$ s.t.\ $\lambda I \subset R$.
Let $\varphi \in \operatorname{Hom}_{R}(I, R)$. For $x, y \in I$, $\varphi(\lambda x y) = \lambda \varphi(x y) = \lambda x \varphi(y) = \lambda y \varphi(x)$.
Since $R$ is an integral domain, we then have $\varphi(x y) = x \varphi(y) = y \varphi(x)$ (in $F$).
If $x, y \neq 0$, then $y^{-1} \varphi(y) = x^{-1} \varphi(x)$.
So we may define $h_{\varphi} = x^{-1} \varphi(x)$ for any $x \in I \setminus \{0\}$.
Easy to see the map $\operatorname{Hom}_{R}(I, R) \rightarrow \{x \in F \mid x I \subset R\}$ is an isomorphism
\[
\varphi \mapsto h_{\varphi} = x^{-1} \varphi(x)
\]
Moreover, $I^{*}$ is a fractional ideal since $\forall x \in I, x I^{*} \subset R$.

(2)(a) By assumption, $M \simeq \bigoplus_{i=1}^{b} I_{i}$. Then $M^{*} \simeq \bigoplus_{i=1}^{b} I_{i}^{*}$.

(b) The module $M$ is f.g.\ projective $\Leftrightarrow$ $I_{i}$ is f.g.\ projective for all $i$.

So it suffices to prove a fractional ideal $I$ is f.g.\ projective $\Leftrightarrow I I^{*}=R$.

From above, $I$ is f.g.\ projective $\Leftrightarrow I^{*} \otimes_{R} I \simeq \operatorname{End}_{R}(I)$.

\[
\begin{aligned}
& \Leftrightarrow \exists \sum_{i=1}^{n} b_{i} \otimes a_{i} \in I^{*} \otimes I \text{ s.t. } \sum_{i=1}^{n} x b_{i} a_{i}=x \quad \forall x \in I \\
& \Leftrightarrow \sum_{i=1}^{n} a_{i} b_{i}=1 \\
& \Leftrightarrow I^{*} I=R
\end{aligned}
\]

\end{proof}

\begin{definition}
Dedekind domain \& principal ideal domain

(1) A Dedekind domain is an integral domain $R$ all ideals of which are finitely generated projective.

(2) A principal ideal domain is an integral domain $R$ all ideals of which are free (hence of rank one). ($\Leftrightarrow$ all ideals are principal)
\end{definition}

\begin{definition}
Local ring (probably non-commutative)

We say a ring $R$ (probably non-commutative) is local if the set of nonunits forms an ideal of $R$.

\textbf{An equivalent definition of a local ring}

Let $R$ be a ring with $1$. TFAE:
\begin{enumerate}
\item $R$ is a local ring.
\item The set of nonunits of $R$ coincides with $J(R)$. That is, $J(R)$ is a unique maximal left (or right) ideal of $R$.
\item $R / J(R)$ is a division ring.
\end{enumerate}
\end{definition}

\begin{proof}
Let $M$ denote the set of nonunits of $R$.

$(1) \Rightarrow (2)$, $(2) \Rightarrow (3)$.

$(3) \Rightarrow (1)$: If $\bar{a} \neq 0$ in $R / J(R)$, then $\exists b \in R, c \in J(R)$ s.t.\ $ba = 1 - c$ has a left inverse $\Rightarrow$ $a$ has a left inverse. Similarly, $a$ has a right inverse. Hence $a$ is a unit.
\end{proof}

\begin{proposition}
$R$ is a commutative local ring $\Rightarrow$ any f.g.\ projective module is free.

Let $R$ be a commutative ring. Assume $R$ is local with unique maximal ideal $\mathfrak{m}$. Then any f.g.\ projective $R$-module is free.
\end{proposition}

\begin{proof}
Let $X$ be a f.g.\ projective $R$-module.
Then $X/IX$ is a finite-dimensional $(R/I)$-vector space, say of dimension $r$. Choose a minimal basis $w_{1}, \ldots, w_{n}$ of $X$. Then $\overline{w_{1}}, \ldots, \overline{w_{r}}$ is an $R/I$-basis of $X/IX$.

\textbf{Remark:} an $R/I$-basis of $X/IX$.

Define a surjective map $\varphi: R^{r} \rightarrow X$
\[
\left(a_{1}, \ldots, a_{r}\right) \mapsto \sum_{i=1}^{r} a_{i} w_{i}.
\]
If we set $K = \ker \varphi$, then from the minimal basis property, $\sum_{i=1}^{r} a_{i} w_{i}=0 \Rightarrow a_{i} \in I$ for all $i$. Hence $K \subset I^{r}$.

Since $X$ is projective, there exists $\psi: X \rightarrow R^{r}$ such that $\varphi \psi = 1_{X}$. Then we have $R^{r} = \psi(X) \oplus K$. Then
\[
K = K \cap I^{r} = K \cap I(\psi(X) \oplus K) = K \cap IK.
\]
Hence $K = IK$. Since $K$ is a quotient of $R^{r}$, $K$ is finitely generated over $R$; thus applying Nakayama's lemma again we get $K = 0$. $\Rightarrow X \simeq R^{r}$.

\end{proof}

\begin{corollary}
A local Dedekind domain is a local PID.
\end{corollary}

\begin{remark}
(1) Minimal basis: If $W$ is a set of generators of an $A$-module $M$ which is minimal, in the sense that any proper subset of $W$ does not generate $M$, then $W$ is said to be a minimal basis of $M$. Two minimal bases do not necessarily have the same number of elements.

(2) Local ring: $\{\text{minimal basis of } M\} \Leftrightarrow \{\text{basis of } \bar{M}\}$, hence any minimal basis has the same number of elements.
\end{remark}

\begin{proof}
Use Nakayama's lemma.

Let $\overline{u}_{1}, \ldots, \overline{u}_{n}$ be a basis of $\bar{M}$, then $M = \sum_{i=1}^{n} A u_{i} + \mathfrak{m} M$. Since $M$ is finitely generated, by Nakayama's lemma, $M = \sum_{i=1}^{n} A u_{i}$. The other way around is obvious.
\end{proof}

\begin{proposition}
Localization.

Let $R$ be an integral domain. Let $S$ be a nonempty multiplicatively closed subset of $R \setminus \{0\}$.
\begin{itemize}
\item If $R$ is a PID, so is $S^{-1} R$.
\item If $R$ is a Dedekind domain, so is $S^{-1} R$.
\end{itemize}
\end{proposition}

\begin{definition}
Invertible (fractional ideal).

A finitely generated projective fractional ideal is called invertible.
(Recall that a fractional ideal is finitely generated projective iff $II^{*} = R$.)
\end{definition}

\begin{lemma}
An integral domain $R$ is a Dedekind domain iff all fractional ideals are invertible.
\end{lemma}

\begin{proof}
$(\Leftarrow)$ If all fractional ideals are finitely generated projective, so are in particular the ideals of $R$.
$\leftarrow \pm 1$ all fractional ideals are invertible, so are in particular the ideals of $R$.

$\Rightarrow$ Given a fractional ideal $I$, if $\lambda I \subseteq R$ for some $\lambda \neq 0$, the map $I \rightarrow \lambda I; a \mapsto \lambda a$ is an isomorphism of $R$-modules, hence $I$ is invertible.
\end{proof}

\begin{proposition}
The set of fractional ideals of a Dedekind domain is an abelian group. Let $R$ be a Dedekind domain. The set of fractional ideals of $R$ is an abelian group under multiplication. The identity element is $R$, and the inverse of a fractional ideal is $I^{*}$ (from now on denoted by $I^{-1}$).

This group is called the ideal group of $R$.
\end{proposition}

\begin{definition}[Integral Ideal; $I/J$]
(1) A fractional ideal $I$ of $R$ is called an integral ideal if $I \subseteq R$.
(2) Let $I$ and $J$ be fractional ideals in $R$. We say that $J$ divides $I$ if there is an integral ideal $K$ s.t.\ $I = JK$.
\end{definition}

\begin{lemma}[$J \mid I \Leftrightarrow I \subseteq J$]
Let $I$ and $J$ be fractional ideals in $R$. TFAE:
\begin{enumerate}
    \item $I \subseteq J$
    \item $J \mid I$
\end{enumerate}
\end{lemma}

\begin{proof}
$(2) \Rightarrow (1)$: clear.

$(1) \Rightarrow (2)$: Let $K = J^{-1}I$. Then $I = JK$ with $K \subseteq J^{-1}J = R$.
\end{proof}

\begin{corollary}
Every nonzero prime ideal of a Dedekind domain is maximal.

\textbf{pf:} Assume that $P$ is a nonzero prime ideal and $P \subseteq I$ for some ideal $I$. WTS: $I = P$ or $R$.

By above lemma, $\exists$ an ideal $J$ s.t.\ $P = IJ$.
Since $P$ is prime, we have $I \subseteq P$ or $J \subseteq P$.
\begin{itemize}
    \item $I \subseteq P \Rightarrow I = P$.
    \item $J \subseteq P \Rightarrow P = PI \Rightarrow I = R$ (multiply by $P^{-1}$ on both sides).
\end{itemize}
\end{corollary}

\begin{corollary}
A Dedekind domain is integrally closed.
\end{corollary}

\begin{proof}
Let $R$ be a Dedekind domain and let $F$ be its field of fractions.

Let $x \in F$ be integral over $R$:
\[x^{n} + a_{n-1}x^{n-1} + \cdots + a_{1}x + a_{0} = 0\]
for some $a_{i} \in R$.

Then $x$ belongs to the fractional ideal $I$ generated by $1, x, \ldots, x^{n-1}$.

Check: $I^{2} = I \Rightarrow I = R$ (multiply by $I^{-1}$ on both sides).
\end{proof}

\begin{theorem}[Factorization into Prime Ideals in Dedekind Domain]
Let $R$ be a Dedekind domain. Every nonzero ideal $I$ of $R$ can be written uniquely as a product of prime ideals:
\[I = P_{1}^{e_{1}} P_{2}^{e_{2}} \cdots P_{r}^{e_{r}}\]
where $P_{i}$ are distinct prime ideals and $e_{i} \geq 1$.
\end{theorem}
Let $R$ be a Dedekind domain. Every nonzero proper ideal of $R$ can be written in a unique way as a product of prime ideals.

\textbf{Proof:}

\textbf{Step 1:} Prove any proper ideal is a product of prime ideals. If not, then there exists an ideal $I$ which is maximal among proper ideals which are not products of prime ideals (by Noetherian property of $R$). In particular, $I$ is not prime, hence is not maximal. So $\exists$ an ideal $J$ s.t.\ $I \subset J \subset R$. So $I = JK$ for some ideal $K$ of $R$, $I \subseteq K$. Moreover, $I \neq K$ otherwise $J = R$. Hence $I \subseteq K \subsetneq R$. By maximality of $I$, $J$ and $K$ are products of prime ideals, so is $I$, a contradiction.

\textbf{Step 2:} Uniqueness. Assume $I = \mathfrak{p}_{1} \ldots \mathfrak{p}_{m} = \mathfrak{q}_{1} \ldots \mathfrak{q}_{n}$ where $\mathfrak{p}_{i}$, $\mathfrak{q}_{j}$ are prime. We prove by induction on $m$. Since $\mathfrak{q}_{1} \cdots \mathfrak{q}_{n} \subset \mathfrak{p}_{1}$, $\exists j$ s.t.\ $\mathfrak{q}_{j} \subset \mathfrak{p}_{1}$. WLOG, assume $j=1$. Then $\mathfrak{q}_{1} = \mathfrak{p}_{1}$ by maximality of prime ideals. Then we have $\mathfrak{p}_{2} \ldots \mathfrak{p}_{m} = \mathfrak{q}_{2} \ldots \mathfrak{q}_{n}$ and it follows by induction hypothesis.

\begin{corollary}
Factorization of fractional ideals in a Dedekind Domain.

Let $R$ be a Dedekind Domain and $I$ be any fractional ideal. Then there exist integral ideals $J$ and $K$ such that $I = J K^{-1}$. Thus for every fractional ideal $I$, we can write
\[
I = \prod_{\mathfrak{p} \text{ prime }} \mathfrak{p}^{v_{\mathfrak{p}}(I)}, \quad v_{\mathfrak{p}}(I) \in \mathbb{Z}
\]
We call $v_{\mathfrak{p}}(I)$ the valuation of $I$ at $\mathfrak{p}$.
\end{corollary}

\begin{proposition}
Let $I$ be a fractional ideal. Then there exist integral ideals $J$ and $K$ such that $I = J K^{-1}$.
\end{proposition}

\begin{proof}
Let $a \in (R \cap I^{*}) \backslash \{0\}$. Then $a I \subset R$ and $Ra \subset I$. We can set $J = a I$ and $K = Ra$.
\end{proof}

\begin{definition}
Valuation of a field.
Let $k$ be a field. An (exponential) valuation of $k$ is a real-valued function $v$ defined on $k^{x}=k \setminus \{0\}$ satisfying
\begin{enumerate}
\item $v(ab)=v(a)+v(b)$
\item $v(a+b) \geqslant \min \{v(a), v(b)\}$
\end{enumerate}
Denote by $(k, v)$ the valuation field.
Then $v\colon K^{x} \rightarrow (\mathbb{R}, +)$ is a group homomorphism.
The image $v(K^{x})$ is called the value group of $v$.
We set $v(0)=\infty$.
We have:
\begin{enumerate}
\item $v(1)=0$, $v(-1)=0$, $v(a)=v(-a)$, $v\left(a^{-1}\right)=-v(a)$
\item $v(a)>v(b) \Rightarrow v(a+b)=v(b)$
\end{enumerate}
Let $R=\{a \in K \mid v(a) \geqslant 0\}$. $R$ is a subring of $K$, called the valuation ring of $v$.
Let $R^{-1}=\{a^{-1} \mid a \in R \setminus \{0\}\}=\{a \in k \mid v(a) \leqslant 0\}$.
Then $K=R \cup R^{-1}$ and $K$ is the fractional field of $R$.
The set $P=\{a \in K \mid v(a)>0\}$ is an ideal of $R$, called the valuation ideal of $v$.
The group of units of $R$ is $\{a \in K \mid v(a)=0\}$.
$\leadsto$ $R$ is a local ring with $P=J(R)=\{\text{nonunits of } R\}$.
$R / P$ is a field, called the residue field of $v$.

\end{definition}

\begin{definition}
equivalent valuation \& discrete valuation

Let $\lambda \in \mathbb{R}$, $\lambda>0$, define
\[
v^{\prime}(a)=\lambda v(a), \quad \forall a \in k^{x} \quad (v \sim v^{\prime})
\]
Then $v^{\prime}$ is also a valuation. We say $v^{\prime}$ is equivalent to $v$.

A valuation $v$ of $K$ is called discrete if its value group is isomorphic to $\mathbb{Z}$. In this situation, $R$ is called a discrete valuation ring. In the case $v\left(K^{x}\right)=\mathbb{Z}$, $v$ is said to be normalized.

\end{definition}

\begin{theorem}
$R$ is PID and is integrally closed for $k$.
Let $(k, v)$ be a normalized discrete valuation field. Let $R$ denote the valuation ring and $P$ denote the valuation ideal.
\begin{enumerate}
\item Let $\pi \in R$ with $v(\pi)=1$. Then any $\alpha \in K$ is expressed as $\alpha=\pi^{i} u$ for some $u \in R \setminus P$ where $i=v(\alpha)$.
\item Any nonzero ideal of $R$ is of the form $P^{i}=\left(\pi^{i}\right)$ for some $i \geqslant 0$. In particular, $R$ is a PID.
\item $R$ is integrally closed.
\end{enumerate}
\end{theorem}
\begin{proof}
(1) Let $\alpha \in K$ with $v(\alpha)=i$. Then $v\left(\alpha \pi^{-i}\right)=0$, so $u:=\alpha \pi^{-i} \in R \setminus P$ and $\alpha=\pi^{i} u$.

(2) Let $I$ be a nonzero ideal of $R$. Let $i=\min \{v(a) \mid a \in I\}$. Let $a$ be an element in $I$ such that $v(a)=i$. Then $a=\pi^{i} u$ for some unit $u$, hence $\pi^{i}=a u^{-1} \in I$. For any element $b \in I$ with $v(b)=j$, we have $b=\pi^{j} u^{\prime}=\pi^{i}\left(\pi^{j-i} u^{\prime}\right)$, hence $b \in\left(\pi^{i}\right)$. In particular, $P=(\pi)$ and $I=P^{i}$.

(3) Let $x \in K$ be integral over $R$. Then
\[x^{n}+a_{n-1} x^{n-1}+\cdots+a_{1} x+a_{0}=0\]
for some $a_{i} \in R$. If $x \notin R$, then $x^{-1} \in R$, hence $x=-\left(a_{n-1}+a_{n-2} x^{-1}+\cdots+a_{0} x^{1-n}\right) \in R$.
\end{proof}

\begin{corollary}
Two discrete valuations $v$ and $v^{\prime}$ of $k$ are equivalent if and only if they have the same valuation ideal.
\end{corollary}

\begin{proof}
It suffices to prove $\Leftarrow$.

$R=k \setminus P^{-1}$ is the valuation ring where $P^{-1}=\left\{x^{-1} \mid x \in P \setminus\{0\}\right\}$. Hence they have the same valuation ring. Then $\alpha=\pi^{i} u$ for some $i$ and some $u \in R \setminus P$. Hence if $v$ and $v^{\prime}$ are normalized, they will have the same value on $R$.
\end{proof}

\begin{theorem}
A discrete valuation ring is a local PID.
\end{theorem}

\begin{proof}
It suffices to show $\supseteq$.

Let $R$ be a local PID. Let $P=J(R)=(\pi)$.

\textbf{Claim:} $\bigcap_{i=1}^{\infty} P^{i}=0$.

If $0 \neq a \in \bigcap_{i=1}^{\infty} P^{i}$, then $a=\pi a_{1}=\pi^{2} a_{2}=\cdots$ with $a_{i} \in R$. Then $a_{i}=\pi a_{i+1}$ (since $R$ is an integral domain). This induces an ascending chain
\[\left(a_{1}\right) \subset \left(a_{2}\right) \subset \cdots\]
contradicting the Noetherian property (since $\pi$ is not a unit).

For any $0 \neq a \in R$, there exists $i \in \mathbb{Z}_{\geqslant 0}$ such that $a \in P^{i}$, $a \notin P^{i+1}$. Then $a=\pi^{i} u$ with $u \in R \setminus P$. Set $v(a)=i$. Let $K$ be the fraction field of $R$. Define $v\left(\frac{a}{b}\right)=v(a)-v(b)$. Then $v$ is a discrete valuation.
\end{proof}

\begin{definition}
Multiplicative valuation
Let $v$ be an exponential valuation of a field $k$. Fix a real number $\gamma$, $0<\gamma<1$. Define $\varphi: k \rightarrow \mathbb{R}$; $a \mapsto \gamma^{v(a)}$.

Then:
\begin{enumerate}
\item $\varphi(a) \geqslant 0$ and $\varphi(a)=0 \Rightarrow a=0$.
\item $\varphi(ab)=\gamma^{v(ab)}=\gamma^{v(a)+v(b)}=\gamma^{v(a)}\gamma^{v(b)}=\varphi(a)\varphi(b)$.
\item $\varphi(a+b) \leqslant \max \{\varphi(a), \varphi(b)\}$.
\end{enumerate}
with $\varphi(a)=\varphi(-a)$,
\[
\begin{aligned}
& \varphi(a)=\varphi\left(a^{-1}\right)^{-1}, \\
& \varphi(a)<\varphi(b) \Rightarrow \varphi(a+b)=\varphi(b).
\end{aligned}
\]

The valuation ring $R=\{a \in K \mid \varphi(a) \leq 1\}$. The valuation ideal $P=\{a \in k \mid \varphi(a)<1\}$.

Conversely, if we are given a function $\varphi: K \rightarrow \mathbb{R}$ satisfying (1), (2), (3) and define $v(a)=\log_{\gamma} \varphi(a)$ for any $\gamma$, then $v$ is a valuation on $K$. $\varphi$ is called a \emph{multiplicative valuation} of $k$.

Two multiplicative valuations $\varphi_{1}$ and $\varphi_{2}$ are called \emph{equivalent} if $\exists \gamma \in \mathbb{R}$, $\gamma \geqslant 0$, s.t.\ $\varphi_{2}(a)=\varphi_{1}(a)^{\gamma}$ $\forall a \in K$.

\end{definition}

\begin{definition}
Metric space induced by multiplicative valuation.

Define $d(a, b)=\varphi(a-b)$, a metric on $k$. $K$ is a Hausdorff space. Moreover,
\[
\begin{aligned}
& k \times k \rightarrow k,\quad (a, b) \mapsto a+b, \\
& k^{\times} \times k^{\times} \rightarrow k,\quad (a, b) \mapsto ab^{-1}
\end{aligned}
\]
are continuous.

We say $\lim_{n \rightarrow \infty} a_{n}=a$ if $\forall \varepsilon>0$, $\exists N$ s.t.\ $d\left(a_{n}, a\right)<\varepsilon$ for all $n \geqslant N$.

If every Cauchy sequence in $K$ converges, we say $K$ is \emph{complete}, or a complete field. The associated valuation ring is said to be a complete valuation ring.
\end{definition}

\paragraph{Valuation ring.}
Given a valuation field $K$, let $C$ denote the set of Cauchy sequences in $K$. Define
\[
\left(a_{n}\right)+\left(b_{n}\right)=\left(a_{n}+b_{n}\right), \quad \left(a_{n}\right)\left(b_{n}\right)=\left(a_{n}b_{n}\right).
\]
Then $C$ is a commutative ring.
Define $\left(a_{n}\right)+\left(b_{n}\right)=\left(a_{n}+b_{n}\right)$. Then $C$ is a commutative ring with multiplication $\left(a_{n}\right)\left(b_{n}\right)=\left(a_{n} b_{n}\right)$.

If $I$ denotes the subset of $C$ consisting of the sequences converging to $0$, then $\widetilde{k}=C/I$ is a field. For any $\left(\widetilde{a_{n}}\right) \neq 0$, we may choose a representative such that $a_{n} \neq 0$ for all $n$, and then its inverse is $\left(\frac{1}{a_{n}}\right)$. Indeed, it suffices to find a lower bound for $\varphi\left(a_{n}\right)$: if $\left(\widetilde{a_{n}}\right) \neq 0$, then $\left(a_{n}\right)$ does not converge to $0$, i.e., there exists $\varepsilon>0$ such that for all $N$, there exists $n>N$ with $\varphi\left(a_{n}\right)>\varepsilon$. On the other hand, there exists $N_{\varepsilon}$ such that for all $n, m>N_{\varepsilon}$, $\varphi\left(a_{n}-a_{m}\right)<\frac{\varepsilon}{2}$. Take any $m>N_{\varepsilon}$ such that $\varphi\left(a_{m}\right)>\varepsilon$. Since $\varphi\left(a_{m}\right)>\varphi\left(a_{n}-a_{m}\right)$, we have $\varphi\left(a_{n}\right)=\varphi\left(a_{m}\right)>\varepsilon$.

Then for all $\gamma>0$, there exists $N$ such that for all $n, m>N$,
\[
\varphi\left(\frac{1}{a_{n}}-\frac{1}{a_{m}}\right)=\varphi\left(a_{n}\right)^{-1} \varphi\left(a_{m}\right)^{-1} \varphi\left(a_{m}-a_{n}\right)<\frac{\gamma}{\varepsilon^{2}}.
\]
In particular, $I$ is a maximal ideal, and $\widetilde{k}$ is called the \emph{completion} of $k$.

If $\left(a_{n}\right)$ is a Cauchy sequence, then $\left|\varphi\left(a_{m}\right)-\varphi\left(a_{n}\right)\right| \leq \varphi\left(a_{m}-a_{n}\right)$. Hence $\left(\varphi\left(a_{n}\right)\right)$ is a Cauchy sequence in $\mathbb{R}$ and $\lim _{n \rightarrow \infty} \varphi\left(a_{n}\right)$ exists. Where $\left(x_{n}\right) \sim\left(y_{n}\right)$ if $\left(a_{n}\right)-\left(b_{n}\right) \in I$, then $\lim _{n \rightarrow \infty} \varphi\left(a_{n}\right)=\lim _{n \rightarrow \infty} \varphi\left(b_{n}\right)$ (i.e., if $d\left(x_{n}, y_{n}\right) \rightarrow 0$).

Define $\tilde{\varphi}: \widetilde{k} \rightarrow \mathbb{R}$ by $\tilde{\varphi}\left(\widetilde{\left(a_{n}\right)}\right) = \lim _{n \rightarrow \infty} \varphi\left(a_{n}\right)$ for $\widetilde{\left(a_{n}\right)}=\left(a_{n}\right)+I$.

Check: $\tilde{\varphi}$ is a multiplicative valuation of $\tilde{k}$.

For $a \in k$, the sequence $(a)=(a, a, \ldots)$ is a Cauchy sequence and $\tilde{a} \neq \tilde{b}$ if $a \neq b$, with $\tilde{\varphi}(\tilde{a})=\varphi(a)$. Thus we can regard $k$ as a subfield of $\tilde{k}$ by identifying $a$ with $\tilde{(a)}$, and $\tilde{\varphi}$ is an extension of $\varphi$. The field $k$ is dense in $\tilde{k}$.
\begin{lemma}
$\widetilde{K}$ is a complete field.
\end{lemma}

\begin{theorem}
$\tilde{R} / \tilde{P} \simeq R / P$

Let $\tilde{R}$ and $\tilde{P}$ be the valuation ring and valuation ideal of $\tilde{\varphi}$ respectively. Then $\tilde{R} / \tilde{P} \simeq R / P$.
\end{theorem}

\begin{proof}
It suffices to show $R+\widetilde{P}=\widetilde{R}$ (we have $R / P \rightarrow \widetilde{R} / \widetilde{P} ; a+P \mapsto a+\widetilde{P}$). Let $a \in \widetilde{R}$ with $a \neq 0$. Then $0 \leq \widetilde{\varphi}(a) \leq 1$. Since $K$ is dense in $\widetilde{K}$, $\exists a_n \in K$ such that $a=\lim a_n$. Then $\lim \tilde{\varphi}\left(a_{n}-a\right)=0 \Rightarrow \exists m$ s.t.\ $\tilde{\varphi}\left(a_{m}-a\right)<\tilde{\varphi}(a)$.

$\varphi\left(a_{m}\right)=\tilde{\varphi}\left(a_{m}-a+a\right)=\tilde{\varphi}(a) \leqslant 1 \Rightarrow a_{m} \in R$.

$\tilde{\varphi}\left(a_{m}-a\right)<\tilde{\varphi}(a) \leq 1 \Rightarrow a_{m}-a \in \widetilde{P}$.

Hence $a=a_{m}-\left(a_{m}-a\right) \in R+\widetilde{P}$.
\end{proof}

\begin{proposition}
Let $v$ be a normalized discrete valuation of $K$ with valuation ring $R$ and valuation ideal $P=(\pi)$. Then $\left\{P, P^{2}, \cdots\right\}$ is a fundamental neighbourhood system of $0$ since $\bigcap_{i=1}^{\infty}\left(\pi^{i}\right)=0$.

Therefore $\left(a_{n}\right)$ is a Cauchy sequence $\Leftrightarrow$ for any $l \in \mathbb{Z}_{>0}$, $\exists N$ s.t.\ $a_{m}-a_{n} \in P^{l}$ for all $m, n \geqslant N$.

Also, $\lim _{n \rightarrow \infty} a_{n}=a \Leftrightarrow$ given $l$, $\exists N$ s.t.\ $a_{n}-a \in P^{l}$ for all $n \geqslant N$.
\end{proposition}

\begin{lemma}
Let $\varphi_{1}$ and $\varphi_{2}$ be discrete multiplicative valuations of $K$. TFAE:
\begin{enumerate}
\item $\varphi_{1}$ and $\varphi_{2}$ are equivalent.
\item $\left\{a \in K \mid \varphi_{1}(a)<1\right\}=\left\{a \in K \mid \varphi_{2}(a)<1\right\}$.
\item The topologies on $K$ induced by $\varphi_{1}$ and $\varphi_{2}$ are the same.
\end{enumerate}
\end{lemma}

\begin{proof}
$(1) \Leftrightarrow (2)$: Clear.

$(2) \Rightarrow (3)$: Clear.

$(3) \Rightarrow (2)$: Suppose $\varphi_{1}(a)<1$. Then $\lim _{n \rightarrow \infty} \varphi_{1}\left(a^{n}\right)=\lim _{n \rightarrow \infty} \varphi_{1}(a)^{n}=0$.

So $\lim _{n \rightarrow \infty} \varphi_{2}(a)^{n}=0 \Rightarrow \varphi_{2}(a)<1$.

Similarly $\varphi_{2}(a)<1 \Rightarrow \varphi_{1}(a)<1$.
\end{proof}
\begin{definition}
Norm on $K$-vector space w.r.t.\ $\varphi$.

Let $\varphi$ be a multiplicative valuation of $K$. Let $V$ be a vector space over $K$. A real-valued function $\|\cdot\|$ defined on $V$ is said to be a norm on $V$ if
\begin{enumerate}
\item $\|v\| \geqslant 0$ and $\|v\|=0 \Leftrightarrow v=0$.
\item $\|av\|=\varphi(a)\|v\| \quad \forall a \in K, v \in V$.
\item $\|u+v\| \leq \max \{\|v\|,\|u\|\}$.
\end{enumerate}
If we set $d(u, v)=\|u-v\|$, then $V$ becomes a metric space.
\[
\begin{cases}
V \times V \to V, & (u, v) \mapsto u+v \text{ are continuous} \\
K \times V \to V, & (a, v) \mapsto av
\end{cases}
\]
\end{definition}

\begin{lemma}
Induced norm on $K$-vector space w.r.t.\ $\varphi$. Let $u_{1}, \ldots, u_{n}$ be a basis of a finite-dimensional $K$-vector space $V$.

For $v=a_{1} u_{1}+\cdots+a_{n} u_{n} \in V$, define $\|v\|_{0}=\max \left\{\varphi\left(a_{i}\right) \mid 1 \leq i \leq n\right\}$.

Then $\|\cdot\|_{0}$ is a norm on $V$. The following holds for the metric induced by $\|\cdot\|_{0}$:
\begin{enumerate}
\item Let $v_{r}=a_{r1} u_{1}+\cdots+a_{rn} u_{n}$ be given for $r=1,2, \ldots$. Then the sequence $\{v_{r}\}_{r}$ is a Cauchy sequence if and only if the sequence $(a_{ri})_{r}$ is a Cauchy sequence in $K$ for all $i=1,2, \ldots, n$. Also
\[
\lim _{r \rightarrow \infty} v_{r}=\sum_{i=1}^{n} a_{i} u_{i} \Leftrightarrow \lim _{r \rightarrow \infty} a_{ri}=a_{i} \quad \forall i.
\]
\item If $K$ is complete, then so is $V$.
\end{enumerate}
\end{lemma}

\begin{proof}
Immediate.
\end{proof}

\begin{lemma}
$K$ complete $\Rightarrow$ unique norm on $V$ (up to equivalence).
\end{lemma}
Suppose that $k$ is complete. Then for any norm $\|v\|$ on $V$, there exists constants $\lambda, \mu$, s.t.\ $\|v\| \leqslant \lambda\|v\|_{0}$, $\|v\|_{0} \leqslant \mu\|v\|$ for all $v \in V$.
Thus the topologies induced by these two norms are the same.

\begin{proof}
(1) $\lambda$: Let $\lambda=\max \left\{\left\|u_{i}\right\| \mid 1 \leq i \leq n\right\}$.
Then for $v=\sum_{i=1}^{n} a_{i} v_{i}$, $\|v\| \leq \max \left\{\varphi\left(a_{i}\right)\left\|u_{i}\right\|\right\} \leq \lambda\|v\|_{0}$.

(2) $\mu$: We will show the existence of $\mu$ by induction on $n$.

$n=1$: trivial.

Let $n>1$. For $v=a_{1} u_{1}+\cdots+a_{n} u_{n}$, set $\varphi_{i}(v)=\varphi\left(a_{i}\right)$.

Claim: $\varphi_{i}(v) \leq \mu_{i}\|v\|$ for some $\mu_{i}$. If it is true, take $\mu=\max\left\{\mu_{i}\right\}$.

Let $V^{\prime}=K u_{1} \oplus \cdots \oplus K u_{n-1}$. By induction, $\|\cdot\|$ and $\|\cdot\|_{0}$ induce the same topologies on $V^{\prime}$. In particular, $V^{\prime}$ is complete w.r.t.\ both $\|\cdot\|$ and $\|\cdot\|_{0}$.

Suppose otherwise there is no $\mu_{n}$ with $\varphi_{n}(v) \leq \mu_{n}\|v\|$, $\forall v \in V$.
Then for any $j \in \mathbb{Z}_{>0}$, there exists $v_{j} \in V \backslash V^{\prime}$ such that $\varphi_{n}\left(v_{j}\right)>j\left\|v_{j}\right\|$.
This holds also if $v_{j}$ is replaced by its nonzero scalar multiples. Hence we may assume $v_{j}=a_{j1} u_{1}+\cdots+a_{j,n-1} u_{n-1}+u_{n}$.
Thus $\left\|v_{j}\right\|<\left(\frac{1}{j}\right) \varphi_{n}\left(v_{j}\right)=\frac{1}{j} \Rightarrow \lim _{j} v_{j}=0$ w.r.t.\ $\|\cdot\|$.
So $\lim \left(v_{j}-u_{n}\right)=-u_{n}$. But $v_{j}-u_{n} \in V^{\prime}$ for all $j$ and $V^{\prime}$ is closed in $V$. This shows $-u_{n} \in V^{\prime}$, a contradiction.
\end{proof}

\begin{definition}
extension $(K, \varphi) \subseteq(L, \Phi)$ (consider $L$ as an $K$-space)
\end{definition}

\begin{definition}
extension $(k, \varphi) \subseteq(L, \Phi)$ (consider $L$ as an $k$-space)

Let $k, L$ be discrete valuation fields with multiplicative valuations $\varphi$ and $\Phi$ respectively. If $k \subseteq L$ and $\Phi$ is an extension of $\varphi$, then we say $(L, \Phi)$ is an extension of $(k, \varphi)$, and write $(k, \varphi) \subseteq(L, \Phi)$.
Consider $L$ as an $K$-space. Let $\|x\|=\Phi(x)$ for $x \in L$. $\|\cdot\|$ is a norm of $L$ as a $K$-vector space.
\end{definition}

\begin{theorem}
Extension of a complete $(K, \varphi)$ is unique up to equivalence and complete.
\end{theorem}
\begin{theorem}
Let $\varphi$ be a complete discrete multiplicative valuation of $K$ and $L$ be a finite extension of $K$. If there exists an extension $(L, \Phi)$ of $(K, \varphi)$, then $\Phi$ is unique up to equivalence. Furthermore, $(L, \Phi)$ is complete.
\end{theorem}

\begin{remark}
In fact, existence can be proved.
\end{remark}

\begin{proof}
By above lemma (see $\Phi$ as a norm of $L$ as a vector space over $K$).
\end{proof}

\begin{definition}[Metric on an $R$-module]
Let $R$ be a ring, $V$ an $R$-module. Let $I$ be an ideal of $R$. Assume that $\bigcap_{i} I^{i} V = 0$. Fix $r \in R$, $0 < \gamma < 1$. Define $\|\cdot\|_{I}$ on $V$:
\[
\begin{aligned}
& \|0\|_{I} = 0 \\
& \|v\|_{I} = \gamma^{i} \text{ for } v \in I^{i} V \setminus I^{i+1} V
\end{aligned}
\]
Define $d_{I}(u, v) := \|u - v\|_{I}$.

Check: $d_{I}(u, v) \leq \max \left\{d_{I}(u, w), d_{I}(w, v)\right\}$. So $d_{I}$ is a metric on $V$. In this situation, we say $d_{I}$ is defined on $V$.

A sequence $(u_{n})$ is a Cauchy sequence $\Rightarrow$ for any integer $m > 0$, there exists $N > 0$ s.t.\ $v_{i} - v_{j} \in I^{m} V$ for all $i, j \geqslant N$.

In particular, $v = \lim_{n \rightarrow \infty} v_{n} \Leftrightarrow$ for any integer $m > 0$, there exists $N > 0$ s.t.\ $v_{i} - v \in I^{m} V$ for all $i \geqslant N$.

The maps
\begin{enumerate}
\item $V \times V \rightarrow V;\, (a, b) \mapsto a + b$
\item $R \times V \rightarrow V;\, (r, a) \mapsto ra$
\item $f \in \operatorname{Hom}_{R}(V, W)$ where $d_{I}$ is defined on $V$, $W$
\end{enumerate}
are continuous.
\end{definition}

\begin{theorem}
$R$ complete $\Rightarrow$ finitely generated $V$ complete.

Assume $d_{I}$ is defined on $R$. Assume $V$ is finitely generated. If $R$ is complete (w.r.t.\ $d_{I}$), then $V$ is complete (w.r.t.\ $d_{I}$).
\end{theorem}

\begin{proof}
Let $\left\{v_{1}, \ldots, v_{s}\right\}$ be a set of generators of $V$. Let $\left\{w_{i}\right\}$ be a Cauchy sequence in $V$. There exists a sequence $\{m(n)\}$ of nonnegative integers s.t.\ $w_{i} - w_{j} \in I^{m(n)} V$ for all $i, j > n$, with $m(n) \leq m(n+1)$, $\lim m(n) = \infty$.

Let $w_{1} = \sum_{t=1}^{s} b_{1t} v_{t}$ for $b_{1t} \in R$.

$w_{n+1} - w_{n} = \sum_{t=1}^{s} a_{nt} v_{t}$ with $a_{nt} \in I^{m(n)}$.

So $w_{n} = \sum_{t=1}^{s} b_{nt} v_{t}$ with $b_{nt} = b_{1t} + \sum_{j=1}^{n-1} a_{jt}$.

$b_{n+k, t} - b_{nt} = \sum_{j=n}^{n+k-1} a_{jt} \in I^{m(n)}$.

Then $\left(b_{nt}\right)_{n}$ is a Cauchy sequence in $R$ $\Rightarrow$ $b_{t} = \lim_{n \rightarrow \infty} b_{nt}$ exists in $R$.

Let $w = \sum_{t=1}^{s} b_{t} v_{t}$.

Check: $\lim_{n \rightarrow \infty} w_{n} = w$. Hence $V$ is complete.
\end{proof}

\begin{lemma}
Suppose that $d_P$ is defined for all f.g.\ $R$-modules. $R$ is complete w.r.t.\ $d_P$. If $V$ is a f.g.\ $R$-module. Then every submodule of $V$ is closed.
\end{lemma}

\begin{proof}
Let $W$ be a submodule of $V$.

$f: V \rightarrow V / W$ is continuous. $W=f^{-1}(0)$ is closed.

\end{proof}

\paragraph{$R$-algebra where $R$ is a complete discrete valuation ring.}

\begin{example}
Let $I=P$. Then $d_I=v$. Complete w.r.t.\ $d_I$.

Let $R$ be a complete discrete valuation ring with valuation $v$, and let $P=(\pi)$ and $F=R / P$ be the valuation ideal and residue field of $v$ respectively. ($P=J(R)$).

Let $V$ be a f.g.\ $R$-module. Then $V=R u_{1} \oplus \cdots \oplus R u_{r}$ is a direct sum of cyclic submodules. ($R$ is a PID $\Rightarrow$ f.g.\ $R$-module is a direct sum of cyclic $R$-modules).

$\left(\pi^{n}\right) V=\left(\pi^{n}\right) u_{1} \oplus \cdots \oplus\left(\pi^{n}\right) u_{r} \Rightarrow \bigcap_{n=1}^{\infty} \pi^{n} V=0$.

Then $d_P$ is defined for all f.g.\ $R$-modules. $R$ is complete $\Rightarrow V$ is complete. In particular,

In particular, $R \subseteq C(A)$.

Let $A$ be an $R$-algebra, which is f.g.\ as an $R$-module.

Fact: $A \times A \rightarrow A,\ (a, b) \mapsto a b$ is continuous.

\[
A^{*}=A / \pi A \text{ is a f.d.\ } F \text{-algebra.}
\]

\end{example}

\begin{theorem}
Let $A$ be a f.g.\ $R$-algebra. Then
\begin{enumerate}
\item $A\pi \subseteq J(A)$
\item $J(A)^{n} \subseteq A\pi$ for some $n>0$
\end{enumerate}
\end{theorem}

\begin{proof}
(1) Let $V$ be a simple $A$-module.

Then either $\pi V=V$ or $\pi V=0$. Since $V=A v$ for some $0 \neq v \in V$,

$A$ is f.g.\ over $R$, $V$ is f.g.\ over $R$.

Thus if $\pi V=V$, then $V=0$ by Nakayama's lemma.

Hence $\pi V=(A\pi) V=0 \Rightarrow A\pi \subset J(A)$.

(2) $A^{*}=A / A\pi$ is f.d.\ $F$-algebra $\Rightarrow A^{*}$ Artinian $\Rightarrow J\left(A^{*}\right)=J(A) / A\pi$ is nilpotent. Hence $J(A)^{n} \subset A\pi$ for some $n \geqslant 1$.

\end{proof}

\begin{theorem}
Lifting idempotents in f.g.\ $R$-algebra ($R$ complete discrete valuation ring). Let $A$ be a f.g.\ $R$-algebra. Let $I$ be an ideal of $A$ with $I \subseteq J(A)$ and $\bar{A}=A / I$.

\begin{enumerate}
\item Let $\bar{e}$ be an idempotent in $\bar{A}$ and let $\bar{e}=\bar{e}_{1}+\cdots+\bar{e}_{n}$ be an idempotent decomposition in $\bar{A}$. Then $\bar{e}$ lifts to an idempotent $e$ of $A$ and there exist orthogonal idempotents $e_{1}, \ldots, e_{n}$ of $A$ satisfying $\overline{e_{i}}=\bar{e}_{i}$ and $e=e_{1}+\cdots+e_{n}$.
\item An idempotent $e$ of $A$ is primitive $\Leftrightarrow \bar{e}$ is primitive in $\bar{A}$.
\item Let $e, f$ be primitive idempotents of $A$, then
\[
Ae \simeq Af \Leftrightarrow \bar{A}\bar{e} \simeq \bar{A}\bar{f}
\]
\end{enumerate}

\end{theorem}

\begin{proof}
$I \subseteq J(A) \Rightarrow I^{m} \subseteq Az$ for some $m$. Hence $I^{mr} \leqslant A\pi^{r}$ for all $r \geqslant 1$.

Set $\bar{A}^{(r)} = A/I^{mr}$.

(1) $I^{m}$ is a nilpotent ideal in $\bar{A}^{(1)}$. By $A/I^{m}/I/I^{m} \simeq A/I$, $\bar{e}$ lifts to an idempotent of $\bar{A}^{(1)}$, namely, $\exists e^{(1)} \in A$ s.t.\ $\left(e^{(1)}\right)^{2} \equiv e^{(1)} \pmod{I^{m}}$, $e^{(1)} \equiv e \pmod{I}$.

Next, $I^{m}/I^{2m}$ is a nilpotent ideal in $\bar{A}^{(2)}$. Similarly, $e^{(1)}$ lifts to $e^{(2)} \in A$ s.t.\ $\left(e^{(2)}\right)^{2} \equiv e^{(2)} \pmod{I^{2m}}$, $e^{(2)} \equiv e^{(1)} \pmod{I^{m}}$.

Continue—we get $e^{(r)} \in A$ for each $r$ s.t.\ $\left(e^{(r)}\right)^{2} \equiv e^{(r)} \pmod{I^{rm}}$, $e^{(r)} \equiv e^{(r-1)} \pmod{I^{(r-1)m}}$.

Then $\left(e^{(r)}\right)_{r}$ is a Cauchy sequence and $e = \lim_{r \to \infty} e^{(r)}$ exists since $R$ is complete.

$e$ is an idempotent of $A$ that lifts $\bar{e}$ to $A$ (use $\left(e^{(r)}\right)^{2} \equiv e^{(r)} \pmod{I^{rm}}$).

Similarly, there exists $e_{i}^{(r)} \in A$ for each $r$ s.t.\ $e_{i}^{(r)} \equiv e_{i}^{(r-1)} \pmod{I^{(r-1)m}}$, $\left(e_{i}^{(r)}\right)^{2} \equiv e_{i}^{(r)} \pmod{I^{m}}$ and $e^{(r)} \equiv e_{1}^{(r)} + \cdots + e_{n}^{(r)} \pmod{I^{rm}}$ is an idempotent decomposition in $\bar{A}^{(r)}$.

Each $\left(e_{i}^{(r)}\right)_{r}$ is a Cauchy sequence $\Rightarrow$ Let $e_{i} = \lim_{r \to \infty} e_{i}^{(r)} \Rightarrow e_{1}, \ldots, e_{n}$ are the desired idempotents.

(2) Similarly as before (use $J(A)$ has no idempotents).

(3) $Ae$ is the projective cover of $\bar{A}\bar{e}$.
\end{proof}

\begin{remark}
Let $I = Az$, $\bar{A} = A/I = A^{*}$.

\[
\begin{cases}
K \otimes_{R} A : K\text{-alg} & (K \text{ valuation field}) \\
A^{*} : K\text{-alg} \\
A^{*} : F\text{-alg} & (F \text{ residue field})
\end{cases}
\]
Let $A$ be an algebra over a commutative ring $R$. Let $U$ be an $A$-module.

If $R \rightarrow R^{\prime}$ is a homomorphism to another commutative ring $R^{\prime}$, we can form the $R^{\prime}$-algebra $R^{\prime} \otimes_R A$, and $R^{\prime} \otimes_R U$ becomes an $R^{\prime} \otimes_R A$-module.

\begin{itemize}
\item If $R$ is a subring of $R^{\prime}$, we say that the module $R^{\prime} \otimes_R U$ is obtained from $U$ by extending the scalars from $R$ to $R^{\prime}$.

\item If $R^{\prime}=R / I$ for some ideal $I$ of $R$, then applying $R^{\prime} \otimes_R -$ to an $R$-module $U$ is the same as reducing $U$ modulo $I$:
\[
R/I \otimes_R U \simeq U / I U.
\]

\item When $U$ is free as an $R$-module, we can identify $U$ with $R \otimes_R U$, and an $R$-basis of $U$ becomes an $R^{\prime}$-basis of $R^{\prime} \otimes_R U$.
\end{itemize}

\textbf{Counterexample:} $\mathbb{Q} \otimes_{\mathbb{Z}} \mathbb{Z}/2\mathbb{Z}$. We have
\[
1 \otimes 1 = \frac{1}{2} \otimes 2 = 0.
\]

If an $R^{\prime} \otimes_R A$-module $V$ has the form $R^{\prime} \otimes_R U$ for some $A$-module $U$, we say that $V$ is \emph{defined over} $R$ and that $V$ can be \emph{lifted} to $U$ ($U$ is a lift of $V$).

\end{remark}

\begin{proposition}

\end{proposition}

\paragraph{w.r.t.\ $E$}

Let $F \subseteq E$ be fields where every element of $E$ is algebraic over $F$. Let $A$ be a finite-dimensional $F$-algebra. Let $V$ be a finite-dimensional $E \otimes_F A$-module. Then there exists a field $K$ with $F \subseteq K \subseteq E$ of finite degree over $F$, such that $V$ is defined over $K$.

\begin{proof}
Let $a^1, \ldots, a^n$ be an $F$-basis of $A$ and let $a^t$ (identifying $A$ with $F \otimes_F A$) act on $V$ with matrix $(a_{ij}^t)$ with respect to some $E$-basis of $V$. Let
\[
K = F\left[a_{ij}^t \mid 1 \leq t \leq n,\ 1 \leq i, j \leq d\right]
\]
(where $\dim_E V = d$). Then $[K:F] < \infty$ since $E$ is algebraic over $F$.

Then $A$ acts by matrices over the field $K$.
\end{proof}

\begin{definition}
absolutely simple
Let $A$ be a finite-dimensional $F$-algebra where $F$ is a field. A simple $A$-module $U$ is said to be \emph{absolutely simple} if $U$ is simple as an $E \otimes_F A$-module for all extension fields $E$ of $F$.

We say an extension field $E$ of $F$ is a \emph{splitting field} for $A$ if every simple $E \otimes_F A$-module is absolutely simple. ($E \otimes_F A$ is the new ``$A$'' which is a finite-dimensional $E$-algebra; now take any extension $K$ of $E$, then for any simple $A \otimes_F E$-module $U$, $K \otimes_E U$ is simple as an $A \otimes_F K$-module.)

\end{definition}

\begin{definition}
\end{definition}

\begin{proposition}
Equivalent definition of absolutely simple.

Let $A$ be a finite-dimensional $F$-algebra where $F$ is a field. Let $U$ be a simple $A$-module. TFAE:
\begin{enumerate}
\item $U$ is absolutely simple.
\item $\operatorname{End}_A(U) = F$.
\item The matrix algebra summand of $A/J(A)$ corresponding to $U$ has the form $M_n(F)$ where $n = \dim_F U$.
\end{enumerate}
\end{proposition}

\begin{proof}
(1) $\Rightarrow$ (2): Will not prove.

(2) $\Rightarrow$ (3): $\checkmark$

(3) $\Rightarrow$ (1): $A \to A/J(A) \simeq M_n(F) \oplus \cdots$. Then $U = F^n$. Let $F \subseteq E$ be an extension. $E \otimes_F A$ acts on $E \otimes_F U \simeq E^n$.
\[
E \otimes_F A \to E \otimes_F M_n(F) \simeq M_n(E).
\]
$E^n$ is a simple $M_n(E)$-module $\Rightarrow$ $E \otimes_F U$ is a simple $E \otimes_F A$-module.
\end{proof}

\paragraph{Representation theory of finite groups}

\begin{example}
The Galois group of $x^4 - 2$ over $\mathbb{Q}$.

$x^4 - 2$ has four roots in $\mathbb{C}$: $\alpha = \sqrt[4]{2}$, $i\alpha$, $-\alpha$, $-i\alpha$.
\[
\begin{cases}
\sigma: \alpha \leftrightarrow -i\alpha,\, -\alpha \leftrightarrow i\alpha \Rightarrow \langle \sigma, \tau \rangle \simeq D_8 \\
\tau: i\alpha \leftrightarrow -i\alpha
\end{cases}
\]
preserve the square. $D_8$ is ``represented'' as a subgroup of $\operatorname{GL}_2(\mathbb{R})$:
\[
\sigma \mapsto \begin{pmatrix}
0 & -1 \\
1 & 0
\end{pmatrix}, \quad \tau \mapsto \begin{pmatrix}
1 & 0 \\
0 & -1
\end{pmatrix}.
\]
\end{example}

\paragraph{General representation}
Let $\mathcal{C}$ be a category. For example:
\begin{itemize}
\item $\mathcal{C}$ is the category $\mathbf{Top}$ of topological spaces.
\item $\mathcal{C}$ is the category $\mathbf{Vect}_k$ of vector spaces over a field $k$, or the full subcategory $\mathbf{vec}_k$ of $\mathbf{Vect}_k$ whose objects are of finite dimension.
\item $\mathcal{C}$ is the category $\mathbf{Set}$ of sets, or its full subcategory $\mathbf{set}$ whose objects are of finite cardinalities.
\end{itemize}

If $X$ is an object of $\mathcal{C}$, the group $\operatorname{Aut}(X)$ of automorphisms of $X$ consists of all morphisms from $X$ to $X$ which admit an inverse morphism.

\begin{definition}
Let $G$ be a group. A representation of $G$ on $\mathcal{C}$ (or a $\mathcal{C}$-representation of $G$) is a pair $(X, \rho)$ where
\begin{itemize}
\item $X$ is an object of $\mathcal{C}$,
\item $\rho: G \rightarrow \operatorname{Aut}(X)$ is a group homomorphism.
\end{itemize}

Note that:
\begin{itemize}
\item A representation $(X, \rho)$ is said to be faithful if $\rho$ is injective.
\item A morphism between two representations $(X, \rho)$ and $(X', \rho')$ of $G$ on the same category $\mathcal{C}$ is a morphism $f: X \rightarrow X'$ in $\mathcal{C}$ such that for all $g \in G$, the following diagram commutes:
\[
\begin{CD}
X @>{\rho(g)}>> X \\
@V{f}VV @VV{f}V \\
X' @>{\rho'(g)}>> X'
\end{CD}
\]
\end{itemize}

We can define the category $\operatorname{Rep}(G, \mathcal{C})$ of $\mathcal{C}$-representations of $G$. We can further define isomorphisms between two representations. If $(X, \rho)$ is a $\mathcal{C}$-representation of $G$, then we say $G$ ``acts on'' $X$.
\end{definition}

\begin{example}
$\rho: D_8 \rightarrow GL_2(\mathbb{R}) \subset \operatorname{Aut}(\mathbb{R}^2)$ where $\mathbb{R}^2 \in \mathbf{Vect}_{\mathbb{R}}$.

i.e., $(\mathbb{R}^2, \rho)$ is a $\mathbf{Vect}_{\mathbb{R}}$-representation of $D_8$.
\end{example}

\begin{lemma}
\begin{enumerate}
\item A morphism $f: (X, \rho) \rightarrow (X', \rho')$ in $\operatorname{Rep}(G, \mathcal{C})$ is an isomorphism if and only if the morphism $f: X \rightarrow X'$ is an isomorphism in $\mathcal{C}$.
\item Two representations $(X, \rho)$ and $(X, \rho')$ on the same object $X$ are isomorphic if there exists an automorphism $f$ of $X$ such that for all $g \in G$,
\[
\rho'(g) = f \rho(g) f^{-1}.
\]
We write $(X, \rho) \simeq (X, \rho')$ and say: $\rho$ and $\rho'$ are \emph{uniformly conjugate}.
\end{enumerate}
\end{lemma}

\begin{definition}
The group $\operatorname{Aut}(X) \times \operatorname{Aut}(Y)$ acts on the set $\operatorname{Mor}_{\mathcal{C}}(X, Y)$ by
\[
(\alpha, \beta) \cdot f = \beta \circ f \circ \alpha^{-1} \quad \text{for any } \alpha \in \operatorname{Aut}(X), \beta \in \operatorname{Aut}(Y), f \in \operatorname{Mor}(X, Y).
\]

Composed with the diagonal morphism $G \rightarrow \operatorname{Aut}(X) \times \operatorname{Aut}(Y)$, $g \mapsto (\rho(g), \sigma(g))$, we can define...
\end{definition}
\[
\frac{\text{of } G \text{ on } \operatorname{Mor}(X, Y)}{g \cdot f = \rho(g) f \sigma(g)^{-1}} \quad \forall g \in G, f \in \operatorname{Mor}(X, Y)
\]
commutes.


\begin{lemma}
See $\operatorname{Mor}(X, Y)$ from the action of $G$ on $\operatorname{Mor}(X, Y)$.

The set of morphisms $(X, \rho) \rightarrow (Y, \sigma)$ in $\operatorname{Rep}(G, \mathcal{C})$ is the set of fixed points of $G$ in its action on $\operatorname{Mor}_{\mathcal{C}}(X, Y)$.
\end{lemma}

\paragraph{Representation in Set Category}

We consider the case: $\mathcal{C}$: objects containing elements, morphisms are maps.

$(X, \rho)$: $\mathcal{C}$-rept of $G$.

For $x \in X$, write $gx = \rho(g)(x)$.

A morphism $(X, \rho) \rightarrow (Y, \sigma)$ is a morphism $f: X \rightarrow Y$ in $\mathcal{C}$ such that
\[
f(gx) = gf(x) \quad \forall x \in X, g \in G.
\]

Denote the fixator of $x$ in $G$ as $G_x = \{g \in G \mid gx = x\}$.

\begin{lemma}
Fixator change regarding $G$-action and isomorphisms.

Let $(X, \rho)$ be a $\mathcal{C}$-rept of $G$.

(1) $G_{gx} = gG_x g^{-1}$

(2) If $f: (X, \rho) \rightarrow (Y, \sigma)$ is an isomorphism of repts, then $G_{f(x)} = G_x$ for all $x \in X$.
\end{lemma}

\begin{definition}
$\operatorname{Aut}(G)$ acts on $G$-repts.

Let $\operatorname{Aut}(G)$ denote the group of automorphisms of $G$.

For $g \in G$, consider
\[
\operatorname{Inn}(g): G \rightarrow G; \quad h \mapsto ghg^{-1} \quad \text{(inner automorphism)}.
\]

Denote by $\operatorname{Inn}(G)$ the group of $\operatorname{Inn}(g)$, $g \in G$.

We have $\operatorname{Inn}(G) \triangleleft \operatorname{Aut}(G)$ as $a^{-1}\operatorname{Inn}(g)a = \operatorname{Inn}(aga^{-1})$ $\forall a \in \operatorname{Aut}(G)$, $g \in G$.

The group $\operatorname{Aut}(G)$ acts on $\mathcal{C}$-repts:
${}^{a}(X, \rho) = (X, {}^{a}\rho = \rho \circ a^{-1})$ where $a \in \operatorname{Aut}(G)$, $(X, \rho)$ is any $\mathcal{C}$-rept.

If $f: X \rightarrow X'$ induces a morphism of reps $(X, \rho) \rightarrow (X', \rho')$ then it also induces a morphism of repts
\[
{}^{a}(X, \rho) \rightarrow {}^{a}(X', \rho').
\]

So $\operatorname{Aut}(G)$ acts on the set of isomorphism classes of $\mathcal{C}$-repts. For $g \in G$, we have $\operatorname{Inn}(g)(X, \rho) \simeq (X, \rho)$ through
i.e. $\operatorname{Inn}(G)$ acts trivially on the isomorphism classes of $\mathcal{C}$-reps.

Thus the group $\operatorname{Out}(G) = \operatorname{Aut}(G)/\operatorname{Inn}(G)$, of outer automorphisms, acts on the isomorphism classes of $\mathcal{C}$-repts.
\end{definition}

\begin{definition}
Set repts: $(X, \rho) \amalg (Y, \sigma)$ and $(X, \rho) \times (Y, \sigma)$.

Let $\mathcal{C}$ denote the Set Category.
\end{definition}

\begin{definition}
$\operatorname{Rep}(G, \mathbf{Set}) \left\{\begin{array}{l}
\text{object: } (X, \rho), X \text{ is a set, } \rho: G \rightarrow \operatorname{Sym}(X) \\
\text{morphism: } (X, \rho) \rightarrow (Y, \sigma) \quad X \rightarrow Y
\end{array}\right.$

Given sets $X$ and $Y$, we can define $X \sqcup Y$, $X \times Y$.

If $Z$ is a set, then $Z \times (X \sqcup Y) = (Z \times X) \sqcup (Z \times Y)$.

Let $(X, \rho)$, $(Y, \sigma)$ be $\mathbf{Set}$-reps of $G$. Define $\mathbf{Set}$-reps:
\begin{itemize}
\item $(X, \rho) \sqcup (Y, \sigma) = (X \sqcup Y, \rho \sqcup \sigma)$ where $(\rho \sqcup \sigma)(g)(z) = \begin{cases} \rho(g)(z) & \text{if } z \in X \\ \sigma(g)(z) & \text{if } z \in Y \end{cases}$
\item $(X, \rho) \times (Y, \sigma) = (X \times Y, \rho \times \sigma)$ where $(\rho \times \sigma)(g)(x, y) = (\rho(g)(x), \sigma(g)(y))$
\end{itemize}
\end{definition}

\begin{definition}
Subrepresentation

Let $(X, \rho)$ be a $\mathbf{Set}$-rep of $G$. If $X_{1} \subseteq X$ is stable under $\rho(g)$ $\forall g$, then $(X_{1}, \rho_{1})$ is a subrepresentation of $(X, \rho)$ where $\rho_{1}(g)$ is the restriction of $\rho(g)$ to $X_{1}$.

Let $X_{1}'$ be the complement of $X_{1}$ in $X$, i.e.\ $X_{1}' = X \setminus X_{1}$. Then $X_{1}'$ is also stable by $G$. Let $\rho_{1}'(g)$ be the restriction of $\rho(g)$ to $X_{1}'$, then $(X, \rho) = (X_{1}, \rho_{1}) \sqcup (X_{1}', \rho_{1}')$.
\end{definition}

\begin{example}
Orbit

$(X, \rho)$: $\mathbf{Set}$-rep of $G$.

For any $x \in X$, the subset of $X$ defined by $X_{1} = \{\rho(g)(x) \mid g \in G\}$ is stable by $G$.
\end{example}

\begin{definition}
Irreducible \& indecomposable rep (same in $\mathbf{Set}$-rep)

Let $(X, \rho)$ be a $\mathbf{Set}$-rep of $G$, $X \neq \varnothing$. TFAE:
\begin{enumerate}
\item $(X, \rho)$ is ``simple'' or ``irreducible'', i.e.\ if any subrep $(X_{1}, \rho_{1})$ satisfies either $X_{1} = \emptyset$ or $X_{1} = X$.
\item $(X, \rho)$ is ``indecomposable'', i.e.\ if $(X, \rho) \simeq (X_{1}, \rho_{1}) \sqcup (X_{1}', \rho_{1}')$ then either $X_{1} = \emptyset$ or $X_{1}' = \emptyset$.
\item $(X, \rho)$ is ``transitive'', i.e.\ for any $x, y \in X$, there exists $g \in G$ s.t.\ $y = \rho(g)(x)$.
\end{enumerate}
\end{definition}

\begin{example}
Left multiplication on left cosets of some subgroup $H$

Let $H$ be a subgroup of $G$. Then $G$ acts on the set $G/H$, the set of left $H$-cosets in $G$, by left multiplication, denoted by $(G/H, \rho_{H})$:
\[
\rho_{H}(g)(xH) = gxH \quad \forall g, x \in G.
\]
\begin{enumerate}
\item $(G/H, \rho_{H})$ is transitive.
\item The stabilizer of $xH$ in $G$ is $xHx^{-1}$.
\end{enumerate}
\end{example}
\begin{definition}[Equivalence relation (orbit)]
Let $(X, \rho)$ be a Set-representation of $G$. Define an equivalence relation on $X$: $x \sim_{G} y \Leftrightarrow \exists g \in G$ s.t.\ $\rho(g)(x)=y$. Denote by $X/{\sim_G}$ the set of equivalence classes of this relation. Then:
\begin{enumerate}
\item Any equivalence class $\Omega$ in $X/{\sim_G}$ defines a transitive subrepresentation $(\Omega, \rho|_{\Omega})$ of $(X, \rho)$.
\item $(X, \rho)=\bigsqcup_{\Omega \in X/{\sim_G}}\left(\Omega, \rho|_{\Omega}\right)$.
\end{enumerate}
\end{definition}

\begin{theorem}[Characterization of transitive representations]
\begin{enumerate}
\item Let $(X, \rho)$ be a transitive representation of $G$. For $x \in X$, let $G_x$ denote the stabilizer of $x$ in $G$. Then the map $G/G_x \rightarrow X$
\[
gG_x \mapsto \rho(g)(x)
\]
defines an isomorphism from $(G/G_x, \rho_{G_x})$ to $(X, \rho)$.
\item For subgroups $H$ and $H'$ of $G$, the representations $(G/H, \rho_H)$ and $(G/H', \rho_{H'})$ are isomorphic if and only if $H$ and $H'$ are conjugate.
\end{enumerate}
\end{theorem}

\begin{definition}[Degree of a representation]
If $(X, \rho)$ is a Set-representation of $G$, the cardinality of $X$ is called the degree of the representation $(X, \rho)$.
\end{definition}

\begin{corollary}
If $G$ is a finite group, then the degree of any transitive Set-representation divides the order of $G$.
\end{corollary}

Let $k$ be a field.

\paragraph{$k$-linear representations and matrix representations}

\begin{definition}[$k$-linear representations]
A vector space representation $(M, \rho)$ of $G$ is called a (f.d.)\ ($k$-linear) representation of $G$, that is, a pair $(M, \rho)$ where $M$ is a $k$-vector space of finite dimension and $\rho: G \rightarrow \operatorname{GL}(M)$ is a group morphism, where $\operatorname{GL}(M)$ is the set of invertible linear transformations $M \rightarrow M$. I.e.,
\begin{itemize}
\item Objects of $\operatorname{Rep}(G, \mathbf{vect}_k)$: $(M, \rho)$ $k$-linear representations.
\item Morphisms: $f: (M, \rho) \rightarrow (N, \psi)$ with $f(gm)=gf(m)$ for all $g \in G$.
\end{itemize}
The dimension of $M$ is called the degree of the representation $(M, \rho)$. The trivial representation $\rho_{\mathrm{tr}}: G \rightarrow \operatorname{GL}(k) \cong k^\times$, $g \mapsto 1$.
\end{definition}

\begin{definition}[Matrix representations]
Assume $\dim M = r$. Fix a basis of $M$. There is a group isomorphism $\operatorname{GL}(M) \rightarrow \operatorname{GL}_r(k)$, where $\operatorname{GL}_r(k) = \{\text{invertible matrices of degree } r \text{ over } k\}$. Then a $k$-linear representation of $G$ can be seen as a pair $(r, \rho)$ where $r \in \mathbb{Z}_{\geqslant 1}$ and $\rho: G \rightarrow \operatorname{GL}_r(k)$ is a group morphism.
\end{definition}

\begin{remark}
Two matrix representations $(r, \rho)$ and $(s, \psi)$ are isomorphic if and only if $r=s$ and there is an element $f \in \operatorname{GL}_r(k)$ s.t.\ $\psi(g)=f \rho(g) f^{-1}$ for all $g \in G$.
\end{remark}
i.e.\ the diagram


\begin{definition}
Representation induced from extension of scalars

Assume $k$ is a subfield of a field $k'$. A $k$-rep $(M, \rho)$ of $G$, its extension of scalars is a $k'$-rep $(k' \otimes M, k' \otimes \rho)$ where $k' \otimes \rho$ is the morphism $G \to GL(k' \otimes M)$
\[
(k' \otimes \rho)(g) = \operatorname{Id}_{k'} \otimes \rho(g) \quad \forall g \in G.
\]

Note that a basis of $M$ (over $k$) can be identified with a basis of $k' \otimes M$ (over $k'$). Then from the point of view of a matrix rep $G \to GL_r(k)$, extending the scalar means viewing the matrix entries of $\rho(g)$ as elements of $k'$. Conversely, a $k'$-rep $(M', \rho')$ of $G$ is said to be \emph{rational over} $k$ if there is a $k$-rep $(M, \rho)$ s.t.\ $(M', \rho') \simeq (k' \otimes M, k' \otimes \rho)$. From the point of view of matrix reps, a matrix rep $(r, \rho')$ of $G$ over $k'$ is rational over $k$ if there exists $U \in GL_r(k')$ s.t.\ for all $g \in G$, the entries of $U \rho'(g) U^{-1}$ lie in $k$.

\end{definition}

\begin{definition}
$k$-reps: $(M, \rho) \oplus (N, \psi)$ \& $(M, \rho) \otimes (N, \psi)$

Let $(M, \rho), (N, \psi)$ be $k$-reps of $G$. Analog of disjoint union in Set-rep:
\[
(M, \rho) \oplus (N, \psi) = (M \oplus N, \rho \oplus \psi): (\rho \oplus \psi)(g)(m, n) = (\rho(g)(m), \psi(g)(n))
\]
\[
(M, \rho) \otimes_k (N, \psi) = (M \otimes_k N, \rho \otimes_k \psi): (\rho \otimes \psi)(g) = \rho(g) \otimes \psi(g)
\]
Analog of product in Set-rep.

\end{definition}

\begin{definition}
Subrepresentation

If $M_1$ is a subspace of $M$ which is stable under the action of all $\rho(g)$, $g \in G$, we get a subrepresentation $(M_1, \rho_1)$ of $(M, \rho)$ where $\rho_1(g)$ is the restriction of $\rho(g)$ to $M_1$.

In this case, we have the quotient rep $(M/M_1, \bar{\rho})$ where $\bar{\rho}(g)$ is an automorphism of $M/M_1$ induced by $\rho(g)$.

\end{definition}

\begin{definition}
Indecomposable \& irreducible
$(M, \rho)$ is called indecomposable if $M \neq 0$ and $(M, \rho) \simeq\left(M_{1}, \rho_{1}\right) \oplus\left(M_{2}, \rho_{2}\right)$ implies either $M_{1}=0$ or $M_{2}=0$.
$(M, \rho)$ is called simple or irreducible if $M \neq 0$ and the only $G$-stable subspaces of $M$ are $M$ and $0$.

\end{definition}

\begin{theorem}
(Krull-Schmidt Theorem)

Let $(M, \rho)$ be a representation of $G$ of finite degree.
(1) It is a direct sum of indecomposable representations.
(2) If $G$ is finite and if $(M, \rho) \simeq \bigoplus_{i \in I}\left(M_{i}, \rho_{i}\right) \simeq \bigoplus_{j \in J}\left(M_{j}^{\prime}, \rho_{j}^{\prime}\right)$ where $\left(M_{i}, \rho_{i}\right),\left(M_{j}^{\prime}, \rho_{j}^{\prime}\right)$ are indecomposable, then there is a bijection $\alpha: I \to J$ such that for all $i \in I$, $\left(M_{i}, \rho_{i}\right) \simeq\left(M_{\alpha(i)}^{\prime}, \rho_{\alpha(i)}^{\prime}\right)$.

\end{theorem}

\begin{definition}
from $\operatorname{Rept}(G, \mathbf{Set})$ to $\operatorname{Rept}(G, \mathbf{Vect}_k)$

Let $(\Omega, \rho)$ be a set-representation of $G$.
Denote by $k \Omega$ the $k$-vector space with basis $\Omega$.
Then $\rho$ induces a group morphism $k \rho: G \rightarrow GL(k \Omega)$, a $\mathbf{Vect}_k$-representation of $G$: $(k \Omega, k \rho)$ which is called the permutation representation. This gives a functor: $\operatorname{Rept}(G, \mathbf{Set}) \rightarrow \operatorname{Rept}(G, \mathbf{Vect}_k)$.
In particular, this construction transforms:
\begin{itemize}
\item disjoint union of set-representations into direct sum of $\mathbf{Vect}_k$-representations
\item product of set-representations into tensor product of $\mathbf{Vect}_k$-representations
\end{itemize}
$\begin{pmatrix} \text{tensor product} \rightarrow \text{linear} \\ \text{direct product} \rightarrow \text{bilinear} \end{pmatrix}$

\end{definition}

\begin{remark}
$(\Omega, \rho) \in \operatorname{Rept}(G, \mathbf{Set})$ irreducible $\Rightarrow$ $(k \Omega, k \rho) \in \operatorname{Rep}(G, \mathbf{Vect}_k)$ irreducible $(|\Omega| \geqslant 2)$. Indeed, (transitive)

$s \Omega:=\sum_{\omega \in \Omega} \omega \in k \Omega$ is invariant under all $\rho(g), g \in G$
$\Rightarrow s \Omega$ spans a $G$-stable line.

\end{remark}

\begin{definition}
regular representation

The group $G$ acts on itself by left multiplication. The corresponding $k$-linear representation is called the regular representation of $G$.

\end{definition}

\begin{lemma}

Assume that $k$ has characteristic $p$ dividing the order of $G$. Then the line $k(s_G)$ has no $G$-stable complement in the regular representation of $G$.
\end{lemma}
\begin{proof}
Let $k\sigma = k(s\sigma) \oplus V'$ as $k$-vector spaces for some subspace $V'$. Then we have a projection $\pi: kG \rightarrow k(sG)$. In particular, $\pi(k(SG)) = k(SG)$. If $V'$ is $G$-stable, then $\pi(gx) = g\pi(x)$ for all $x \in kG$, $g \in G$. This implies $\pi(g \cdot 1) = g\pi(1)$. Since $\pi(1) \in k(SO)$, we have $g\pi(1) = \pi(1)$ for all $g \in G$. Thus $\pi(SG) = |G| \cdot \pi(1) = 0$, a contradiction.
\end{proof}

\begin{definition}[Ways to obtain new representations]
Let $(M, \rho)$, $(N, \sigma)$ be $k$-representations of $G$.
\begin{itemize}
\item $(M \otimes N, \rho \otimes \sigma)$
\item $(M \oplus N, \rho \oplus \sigma)$
\item $\operatorname{Hom}_{k}(M, N)$
\end{itemize}

For all $\alpha \in \operatorname{Hom}_{k}(M, N)$, $g \in G$, define $g \cdot \alpha = \sigma(g)\alpha\rho(g)^{-1}$. Then $\operatorname{Hom}_{k}(M, N)$ becomes a representation of $G$.

An element $\alpha \in \operatorname{Hom}_{k}(M, N)$ is invariant under all $g \in G$ if and only if $\alpha$ is a morphism $(M, \rho) \rightarrow (N, \sigma)$.

In particular, there is a representation associated with $M^{*} = \operatorname{Hom}_{k}(M, k)$: for all $\alpha \in \operatorname{Hom}_{k}(M, k)$, $g \cdot \alpha = \alpha \rho(g)^{-1}$, which is called the \emph{contragredient representation} of $(M, \rho)$.
\end{definition}

\begin{definition}[Inflation]
Let $(M, \rho)$ be a $k$-representation of $G/N$ where $N \triangleleft G$. Then
\[
\rho: G/N \rightarrow GL(M) \leadsto \operatorname{Inf}_{G/N}^{G} \rho: G \rightarrow GL(M).
\]
Denote the inflation by $(M, \operatorname{Inf}_{G/N}^{G} \rho)$.
\end{definition}

\begin{proposition}
$\operatorname{Hom}(G, k^{\times}) \simeq \operatorname{Hom}\bigl(G/[G,G], k^{\times}\bigr)$.

Consider representations of degree 1:
\[
G \rightarrow GL(k) = k^{\times}.
\]
Its kernel contains $[G, G]$. Then we have $\operatorname{Hom}(G, k^{\times}) \simeq \operatorname{Hom}\bigl(G/[G,G], k^{\times}\bigr)$ (group homomorphism).
\end{proposition}

\begin{example}
$G = D_{8} = \langle \sigma, \tau \mid \sigma^{2} = \tau^{2} = (\sigma\tau)^{2} = 1 \rangle$.

\[
[G,G] = Z(G) = \langle (\sigma\tau)^{2} \rangle \Rightarrow G/[G,G] \cong \mathbb{Z}/2\mathbb{Z} \times \mathbb{Z}/2\mathbb{Z}.
\]

So $G$ has 4 degree 1 representations over a field $k$ with $\operatorname{char} k \neq 2$:
\begin{align*}
\rho_{1}: \sigma &\mapsto 1, \tau \mapsto -1 \\
\rho_{2}: \sigma &\mapsto 1, \tau \mapsto 1 \\
\rho_{3}: \sigma &\mapsto -1, \tau \mapsto -1 \\
\rho_{4}: \sigma &\mapsto -1, \tau \mapsto 1
\end{align*}
\end{example}

\begin{lemma}
$(M, \rho) \otimes (N, \sigma)$
\end{lemma}
Let $(S, \rho)$ be an irreducible representation of $G$ and let $(k, \sigma)$ be a representation of degree $1$. Then the tensor product representation $(S, \rho) \underset{k}{\otimes} (k, \sigma) = (S \underset{k}{\otimes} k, \rho \underset{k}{\otimes} \sigma) \simeq (S, \rho \otimes \sigma)$ is still irreducible.

\begin{proof}
Scalar representation preserves invariant spaces.
\end{proof}

\begin{remark}
$(S, \rho)$ and $(S, \rho \otimes \sigma)$ may not be isomorphic.
\end{remark}

\begin{definition}
Group algebra

Let $G$ be a finite group. Let $k$ be a field. Define $k G$ to be the group algebra of $G$ over $k$, the $k$-vector space with basis $G$ endowed with multiplication defined by the multiplication of $G$: for $x = \sum x_{s} s$, $y = \sum y_{t} t \in k G$,
\[
xy = \left(\sum_{s \in G} x_{s} s\right)\left(\sum_{t \in G} y_{t} t\right) = \sum_{g \in G}\left(\sum_{st = g} x_{s} y_{t}\right) g
\]
\end{definition}

\begin{lemma}
Let $z: k G \rightarrow k$ be the linear form on $k G$ defined by
\[
z(g) = \begin{cases} 0 & g \neq 1 \\ 1 & g = 1 \end{cases}
\]
\begin{enumerate}
\item $z$ is a ``class function'': for all $x, y \in k G$, $z(xy) = z(yx)$.
\item The map $\hat{\operatorname{im}ath}: k G \rightarrow k G^{*} = \operatorname{Hom}(k G, k)$
\[
x \mapsto (y \mapsto z(xy))
\]
is a $k$-linear isomorphism.
\end{enumerate}
\end{lemma}

\paragraph{Class function}

Identifying $\{ \text{functions } G \rightarrow k \} \longleftrightarrow \operatorname{Hom}_{k}(k G, k)$ linear forms. Define by $F(G, k)$ the $k$-vector space of functions $G \rightarrow k$:
\[
\operatorname{Hom}_{k}(k G, k) = (k G)^{*}
\]
For $f \in F(G, k)$, the inverse image of $f$ under the isomorphism $\hat{\operatorname{im}ath}$ is
\[
f^{\circ} = \sum_{g \in G} f(g^{-1}) g
\]
Let $f \in F(G, k)$ be a class function, that is, $f(gh) = f(hg)$ for all $g, h \in G$. Identify $f$ with its corresponding element in $\operatorname{Hom}_{k}(k G, k)$: $f(xy) = f(yx)$ for all $x, y \in k G$. Denote by $CF(G, k)$ the space of class functions on $G$ (or $k G$).

\begin{remark}
Let $a \in k G$, $f \in F(G, k)$. Denote by $a \cdot f$ the element of $F(G, k)$ defined by
\[
(a \cdot f)(x) = f(xa) \text{ for all } x \in k G
\]
Then $(a \cdot f)^{\circ} = a \cdot f^{\circ}$.
\end{remark}

\begin{definition}
$\gamma_{C}$ and $S_{C}$ for a conjugacy class $C$
Let $\operatorname{Cl}(G)$ be the set of conjugacy classes of $G$, and for $C \in \operatorname{Cl}(G)$ set $\gamma_{C}$ to be the characteristic function of $C$, i.e.
\[
\gamma_{C}(g)= \begin{cases}0 & \text { if } g \notin C \\ 1 & \text { if } g \in C\end{cases}
\]

Set $S_{C}$ to be the sum of elements in $C$, i.e.
\[
S_{C}=\sum_{g \in C} g \in k G .
\]

For $f \in \operatorname{CF}(G, k)$,
\[
\begin{aligned}
& f=\sum_{C \in \operatorname{Cl}(G)} f\left(g_{C}\right) \gamma_{C} \\
& f^{0}=\sum_{C \in \operatorname{Cl}(G)} f\left(g_{C}^{-1}\right) S_{C}
\end{aligned}
\]
where $g_{C}$ is an element in $C$.

\end{definition}

\begin{lemma}
$\operatorname{CF}(G, k) \cong Z(k G)$

Let $Z(k G)$ be the center of $k G$.
\begin{enumerate}
\item The family $\left(\gamma_{C}\right)_{C \in \operatorname{Cl}(G)}$ is a basis of $\operatorname{CF}(G, k)$ and the family $\left(S_{C}\right)$ is a basis of $Z(k G)$.
$\Rightarrow \operatorname{dim}_{k} \operatorname{CF}(G, k)=|\operatorname{Cl}(G)|$
\item The isomorphism $\hat{\operatorname{im}ath}$ restricts to a $k$-linear isomorphism $\operatorname{CF}(G, k) \rightarrow Z(k G)$.
$\Rightarrow \operatorname{dim}_{k} Z(k G)=|\operatorname{Cl}(G)|$
\end{enumerate}
\end{lemma}

\begin{proposition}
$\operatorname{Rep}\left(G, \operatorname{Vect}_{k}\right) \leftrightarrow k G\text{-module}$

Let $(M, \rho)$ be a $k$-representation of $G$. $\rho: G \rightarrow \operatorname{GL}(M)$ extends linearly to an algebra morphism $\rho: k G \rightarrow \operatorname{End}_{k}(M) \Rightarrow M$ becomes a $k G$-module. Conversely, given an algebra morphism $\rho: k G \rightarrow \operatorname{End}_{k}(M)$ (sending 1 to 1), it restricts to a group morphism $\rho: G \rightarrow \operatorname{GL}(M)$.

Thus $\operatorname{Rep}\left(G, \operatorname{Vect}_{k}\right) \leftrightarrow k G\text{-module}$.
\end{proposition}

\begin{lemma}
The regular representation is faithful.

Let $\rho_{k G}: k G \rightarrow \operatorname{End}_{k}(k G)$ denote the morphism associated to the regular representation. Then $\rho_{k G}$ is injective, i.e. the regular representation is faithful.
\end{lemma}

\begin{lemma}
(Schur's lemma)

Let $(S, \rho)$ and $\left(S^{\prime}, \rho^{\prime}\right)$ be irreducible representations of $G$.
\begin{enumerate}
\item If $(S, \rho)$ and $\left(S^{\prime}, \rho^{\prime}\right)$ are not isomorphic, then $\operatorname{Hom}_{k G}\left(S, S^{\prime}\right)=0$.
\item $\operatorname{End}_{k G}(S)$ is a finite-dimensional division $k$-algebra.
\item If $k$ is algebraically closed, then $\operatorname{End}_{k G}(S)=k \cdot \operatorname{Id}$.
\end{enumerate}
\end{lemma}

\begin{remark}
The assumption that $k$ is algebraically closed is in general too strong. We often replace the hypothesis by the hypothesis that $k$ is splitting for $G$. (Recall that $k$ is splitting for an algebra $A$ if
\[
\operatorname{End}_{A}(S) \simeq K \text{ for any simple } A\text{-module } S
\]

If $K$ is splitting for the group algebra $KG$, then we say $K$ is splitting for $G$.

\end{remark}

\begin{theorem}
(Brauer)

If $k$ contains an $\exp(G)$-th root of unity, then $k$ is splitting for $G$.
$(\exp(G) = \operatorname{lcm}\{\mathrm{o}(g), g \in G\}$, so $\exp(G) \mid |G|)$.
\end{theorem}

\begin{proposition}
If $\operatorname{char} k \mid |G|$, then $kG$ is not semisimple.
\end{proposition}

\begin{proof}
Let $w = \sum_{g \in G} g$. For any $x \in G$, $xw = w = wx$, i.e. $w \in Z(kG)$.

Now $w^{2} = w \sum_{g \in G} g = |G| w = 0$ since $\operatorname{char} k \mid |G|$. Then the ideal $kGw$ satisfies $(kGw)^{2} = (kG)^{2} w^{2} = (0)$. $kGw \subset J(kG)$ hence $kG$ is not semisimple.
\end{proof}

\begin{theorem}
(Maschke)

If $\operatorname{char} k \nmid |G|$, then $kG$ is semisimple.
\end{theorem}

\begin{proof}
Let $V$ be a $kG$-module, $W$ a $kG$-submodule of $V$.

Aim: there is a $kG$-submodule $U$ of $V$ with $V = W \oplus U$.

Let $\iota: W \hookrightarrow V$ be the inclusion. Then $\iota$ is a $kG$-homomorphism.

As $k$-vector space: $V = W \oplus W^{\prime}$ for some subspace $W^{\prime}$ of $V$.

Let $p: V \rightarrow W$ be the projection. Then $p \circ \iota = \operatorname{Id}_{W}$.

Consider $\pi: V \rightarrow W$
\[
\pi(v) = \frac{1}{|G|} \sum_{g \in G}\left(\rho(g)^{-1} \circ p \circ \rho(g)\right)(v)
\]

$\forall h \in G$, $\rho(h)^{-1} \pi \rho(h) = \frac{1}{|G|} \sum_{g \in G} \rho(gh)^{-1} \circ p \circ \rho(gh) = \pi$.

Hence $\pi$ is a $kG$-morphism, $\pi \in \operatorname{Hom}_{kG}(V, W)$.

$\pi \circ \iota = \frac{1}{|G|} \sum_{g \in G}\left(\rho(g)^{-1} \circ p \circ \rho(g)\right) \circ \iota = \frac{1}{|G|} \sum_{g \in G} \rho(g)^{-1} \circ (p \circ \iota) \circ \rho(g) = \operatorname{Id}_{W}$.

Let $U = \ker \pi$, then $V = W \oplus U$.
\end{proof}

\begin{corollary}

\end{corollary}

\begin{corollary}
Let $\operatorname{char} k \nmid |G|$. Then any finite-dimensional $k$-representation of $G$ is a direct sum of irreducible representations of $G$.

$\operatorname{char} k \nmid |G| \Rightarrow$ ordinary representation.

$\operatorname{char} k \mid |G| \rightarrow$ modular representation.
\end{corollary}

\begin{proposition}
decomposition of $KG$.
Let $G$ be a finite group. $K$ a field with $\operatorname{char} K \nmid |G|$.
(1) As a ring, $KG$ is a direct sum of matrix algebras over division rings.
(2) Suppose in addition $k$ is splitting for $G$ (for example, $k=\bar{k}$). Let $S_{1}, \ldots, S_{r}$ be pairwise non-isomorphic simple $kG$-modules and let $d_{i}=\operatorname{dim}_{k} S_{i}$. Then $d_{i}$ equals the multiplicity with which $S_{i}$ is a summand of the regular representation of $G$, and $|G| \geqslant d_{1}^{2}+\cdots+d_{r}^{2}$ with equality if and only if $S_{1}, \ldots, S_{r}$ is a complete set of representatives of isomorphism classes of simple $kG$-modules.

\end{proposition}

\begin{definition}
Character

$(M, \rho)$ representation of $G$. $\rho: G \rightarrow GL_{r}(k)$.
$\chi_{M}: G \rightarrow \mathbb{C}$ is called the character afforded by $(M, \rho)$.
\[
g \mapsto \operatorname{tr}(\rho(g))
\]
\end{definition}

\begin{proposition}
Properties of characters.

(1) $\chi_{M}$ is a class function, i.e. $\chi_{M}(g)=\chi_{M}(h)$ if $g, h$ are $G$-conjugate.
(2) $\chi_{M}\left(g^{-1}\right)=\overline{\chi_{M}(g)}$.
(3) If $M \simeq N$ as $\mathbb{C}G$-modules, then $\chi_{M}=\chi_{N}$.
(4) $\chi_{M \oplus N}=\chi_{M} + \chi_{N}$, $\chi_{M^{*}}=\overline{\chi}_{M}$, $\chi_{M \otimes N}=\chi_{M} \chi_{N}$, $\chi_{\operatorname{Hom}(M, N)}=\overline{\chi}_{M} \chi_{N}$.
(5) $\chi_{M}(1)=\operatorname{dim}_{\mathbb{C}} M$.
$\left|\chi_{M}(g)\right| \leq \chi_{M}(1)$ with equality iff the action of $g$ on $M$ is a scalar.
$\chi_{M}(g)=\chi_{M}(1)$ iff $g$ acts trivially on $M$.

\end{proposition}

\begin{definition}
$\langle \cdot, \cdot \rangle$

\[
k=\mathbb{C}. \quad CF(G, \mathbb{C})=CF(G)
\]
$\forall f_{1}, f_{2} \in CF(G)$, define $\left\langle f_{1}, f_{2}\right\rangle:=\frac{1}{|G|} \sum_{g \in G} f_{1}(g) \overline{f_{2}(g)}$
\[
\begin{aligned}
& =\frac{1}{|G|} \sum_{C \in Cl(G)} \frac{|G|}{\left|C_{G}(g_{c})\right|} f_{1}\left(g_{c}\right) \overline{f_{2}\left(g_{c}\right)} \quad\left(g_{c} \in C\right) \\
& =\sum_{C \in Cl(G)} \frac{1}{\left|C_{G}\left(g_{c}\right)\right|} f_{1}\left(g_{c}\right) \overline{f_{2}\left(g_{c}\right)}
\end{aligned}
\]
$\langle \cdot, \cdot \rangle$ is a Hermitian form on the $\mathbb{C}$-vector space $CF(G)$.
If $f_{1}$ and $f_{2}$ are characters of $G$, then $\left\langle f_{1}, f_{2}\right\rangle=\left\langle f_{2}, f_{1}\right\rangle$.
\end{definition}

\begin{lemma}
$V^{G}$

Let $V$ be a $\mathbb{C}G$-module. Define $V^{G}=\{v \in V \mid g v=v \text{ for all } g \in G\}$.
Then $V^{G}$ is a submodule of $V$ and $\operatorname{dim} V^{G}=\left\langle \chi_{V}, 1_{G}\right\rangle=\frac{1}{|G|} \sum_{g \in G} \chi_{V}(g)$.
Then $V^{0}$ is a submodule of $V$ and $\operatorname{dim} V^{0}=\left\langle \chi_{V}, 1 \mid G\right\rangle=\frac{1}{|G|} \sum_{g \in G} \chi_{V}(g)$

\textbf{Pf:}
Let $\pi: V \rightarrow V$

\[
v \mapsto \frac{1}{|G|} \sum_{g \in G} g v
\]

\begin{itemize}
\item $\pi \in \operatorname{Hom}_{\mathbb{C} G}(V, V)$
\end{itemize}

$\forall h \in G, v \in V, \pi(h v)=\frac{1}{|G|}\left(\sum_{g \in G} g h\right) v=\pi(v)=\frac{1}{|G|} \sum_{g \in G} h g v=h \pi(v)$

\begin{itemize}
\item $\pi^{2}=\pi$
\end{itemize}

$\pi^{2}=\left(\frac{1}{|G|} \sum_{g \in G} g\right)\left(\frac{1}{|G|} \sum_{h \in G} h\right)=\frac{1}{|G|^{2}} \sum_{g, h \in G} g h=\pi$

Then $V=\operatorname{ker} \pi \oplus \operatorname{Im} \pi$

\begin{itemize}
\item $V=\operatorname{ker} \pi \oplus \operatorname{Im} \pi$
\end{itemize}

If $v \in V^{G}, \pi(v)=\frac{1}{|G|} \sum_{g \in G} g v=v \Rightarrow V^{G} \subseteq \operatorname{Im} \pi$.

On the other hand, $\forall h \in G, h\left(\frac{1}{|G|} \sum_{g \in G} g v\right)=\frac{1}{|G|} \sum_{g \in G} g v$ hence $\operatorname{Im} \pi \subseteq V^{G}$.

Thus $V^{G}=\operatorname{Im} \pi$ and $V=\operatorname{ker} \pi \oplus V^{G}$

\begin{itemize}
\item $\operatorname{tr}(\pi)=\operatorname{dim} V^{G}$
\item $\operatorname{tr}(\pi)=\operatorname{tr}\left(\frac{1}{|G|} \sum_{g \in G} \rho(g)\right)=\frac{1}{|G|} \sum_{g \in G} \operatorname{tr}(\rho(g))=\frac{1}{|G|} \sum_{g \in G} \chi_{V}(g)$
\end{itemize}

\end{lemma}

\begin{lemma}
$\operatorname{Hom}_{k}(V, W)^{G}=\operatorname{Hom}_{k G}(V, W)$

\end{lemma}

\begin{theorem}
(1${}^{\text{st}}$ orthogonality relation)

Let $V, W$ be simple $k G$-modules. Then

\[
\left\langle \chi_{V}, \chi_{W}\right\rangle= \begin{cases}1 & \text{if } V \simeq W \\ 0 & \text{otherwise}\end{cases}
\]

\end{theorem}

\begin{proof}
$\operatorname{dim}_{\mathbb{C}} \operatorname{Hom}_{\mathbb{C}}(V, W)=\operatorname{dim}_{\mathbb{C}} \operatorname{Hom}_{\mathbb{C}}(V, W)^{G}=\left\langle \chi_{\operatorname{Hom}_{\mathbb{C}}(V, W)}, 1\right\rangle$

\[
\begin{aligned}
& =\left\langle\bar{\chi}_{V} \chi_{W}, 1\right\rangle \\
& =\frac{1}{|G|} \sum_{g \in G} \bar{\chi}_{V}(g) \chi_{W}(g) \\
& =\left\langle \chi_{V}, \chi_{W}\right\rangle .
\end{aligned}
\]

Then apply Schur's lemma.

Let $\operatorname{Irr}(G)$ denote the set of irreducible complex characters of $G$.
\end{proof}

\begin{corollary}
$\operatorname{Irr}(G)$ is $\mathbb{C}$-linearly independent.

\end{corollary}

\begin{corollary}
$\left\langle \chi_{V}, \chi_{V}\right\rangle=1 \Leftrightarrow V$ is irreducible (use complete reducibility)

\end{corollary}

\begin{corollary}
$\chi_{V}=\chi_{W} \Leftrightarrow V \simeq W$.

\end{corollary}

\begin{theorem}
$\operatorname{Irr}(G)$ is a $\mathbb{C}$-basis of $CF(G)$

\end{theorem}
\begin{proof}
$\mathbb{C} G \simeq M_{n_{1}}(\mathbb{C}) \oplus \ldots \oplus M_{n_{r}}(\mathbb{C}), \quad r=|\operatorname{Irr}(G)|$
\[Z(\mathbb{C} G)=Z(M_{n_1}(\mathbb{C})) \oplus \cdots \oplus Z(M_{n_r}(\mathbb{C})) \simeq \underbrace{\mathbb{C} \times \cdots \times \mathbb{C}}_{r}, \quad \dim Z(\mathbb{C} G)=r.\]
\[\dim Z(\mathbb{C} G)=|\mathrm{Cl}(G)|=|\mathrm{CF}(G)|\]
\[Z(\mathbb{C} G)=\bigoplus_{C \in \mathrm{Cl}(G)} \mathbb{C}(\hat{C})\]
Therefore, $|\operatorname{Irr}(G)|=|\mathrm{Cl}(G)|$.
\end{proof}

\begin{corollary}
\end{corollary}

\begin{corollary}
Let $V$ be a $\mathbb{C} G$-module with $V=S_{1}^{n_{1}} \oplus \ldots \oplus S_{k}^{n_{k}}$ where $S_{1}, \ldots, S_{k}$ are non-isomorphic simple $\mathbb{C} G$-modules. Then $n_{i}=\left\langle \chi_{V}, \chi_{S_i}\right\rangle$.
\end{corollary}

\begin{theorem}
($2^{\text{nd}}$ orthogonal relation)

Let $X$ be the character table of $G$, regarded as a matrix, i.e.
\[
\begin{aligned}
& \operatorname{Irr}(G)=\{\chi_{1}, \ldots, \chi_{r}\} \\
& \mathrm{Cl}(G)=\{C_{1}, \ldots, C_{r}\}
\end{aligned}
\]
Let $C=\left(\begin{array}{lll}
\left|C_{G}(x_{1})\right| & & \\
& \ddots & \\
& & \left|C_{G}(x_{r})\right|
\end{array}\right)$ where $x_{1}, \ldots, x_{r}$ are representatives of conjugacy classes of the elements of $G$.

Then $\bar{X}^{\top} X=C$.

Namely, $\sum_{i=1}^{r} \chi_{i}(g) \overline{\chi_{i}(h)}= \begin{cases}\left|C_{G}(g)\right| & \text { if } g \text { and } h \text { are conjugate in } G \\ 0 & \text { otherwise }\end{cases}$
\end{theorem}

\begin{proof}
By $1^{\text{st}}$ orthogonal relation,
\[
X\left(\frac{1}{|G|} \operatorname{diag}\left(|C_{1}|, \ldots, |C_{r}|\right)\right)\bar{X}^{\top}=I
\]
i.e. $X C^{-1} \bar{X}^{\top}=I$.
\[\Rightarrow C^{-1} \bar{X}^{\top} X=I\]
\[\Rightarrow \bar{X}^{\top} X=C\]
\end{proof}

\begin{theorem}
Irreducible representations of abelian groups

TFAE for a finite group $G$:
\begin{enumerate}
\item $G$ is abelian.
\item All irreducible complex representations of $G$ have degree $1$.
\end{enumerate}
\end{theorem}

\begin{proof}
$(1) \Rightarrow (2)$: $\mathbb{C} G \simeq M_{n_{1}}(\mathbb{C}) \oplus \cdots \oplus M_{n_{r}}(\mathbb{C})$. $G$ abelian $\Rightarrow n_{1}=\cdots=n_{r}=1$.

$(2) \Rightarrow (1)$: $|G|=\sum_{\chi \in \operatorname{Irr}(G)} \chi(1)^{2}=|\operatorname{Irr}(G)|=|\mathrm{Cl}(G)| \Rightarrow$ abelian.
\end{proof}

\begin{lemma}
Product of irreducible characters is still irreducible.
\end{lemma}
Let $G_{1}, G_{2}$ be two finite groups. Let $V_{i}$ be a $k G_{i}$-module, $i=1,2$. We can define for all $g_{i} \in G_{i}$, $v_{i} \in V_{i}$:
\[
(g_{1}, g_{2})(v_{1} \otimes v_{2}) = g_{1}v_{1} \otimes g_{2}v_{2}.
\]
Then $V_{1} \otimes V_{2}$ is a $k(G_{1} \times G_{2})$-module. If $\chi_{1} \in \operatorname{Irr}(G_{1})$, $\chi_{2} \in \operatorname{Irr}(G_{2})$, then $\chi_{1} \times \chi_{2}$ (afforded by $V_{1} \otimes V_{2}$) $\in \operatorname{Irr}(G_{1} \times G_{2})$, where $(\chi_{1} \times \chi_{2})(g_{1}, g_{2}) = \chi_{1}(g_{1})\chi_{2}(g_{2})$ for all $g_{i} \in G_{i}$, $i=1,2$.

\begin{proof}
$\langle \chi_{1} \times \chi_{2}, \chi_{1} \times \chi_{2} \rangle = \frac{1}{|G_{1}||G_{2}|} \sum_{g_{i} \in G_{i}} \chi_{1}(g_{1})\chi_{2}(g_{2})\overline{\chi_{1}(g_{1})}\overline{\chi_{2}(g_{2})} = 1$.
\end{proof}

\begin{corollary}
The character table of $G_{1} \times G_{2}$ is the tensor product of their character tables.
\end{corollary}

\begin{proof}
(1) The conjugacy classes in $G_{1} \times G_{2}$ are precisely the products of conjugacy classes.
(2) $\langle \chi_{I_{1}} \otimes \chi_{I_{2}}, \chi_{J_{1}} \otimes \chi_{J_{2}} \rangle = \langle \chi_{I_{1}}, \chi_{J_{1}} \rangle \langle \chi_{I_{2}}, \chi_{J_{2}} \rangle$
\[
= \begin{cases} 1 & I_{1} \simeq J_{1}, \quad I_{2} \simeq J_{2} \\ 0 & \text{otherwise} \end{cases}
\]
Hence the irreducible representations of $G_{1} \times G_{2}$ are precisely the tensor products of irreducible representations.
\end{proof}

\paragraph{The degree of an irreducible representation of $G$ (over $\mathbb{C}$) divides $|G|$}

\begin{definition}[Integral over $\mathbb{Z}$]
Let $R$ be a commutative ring and let $x \in R$. We say that $x$ is \emph{integral over $\mathbb{Z}$} if there exists $n \geqslant 1$ and $a_{1}, \ldots, a_{n} \in \mathbb{Z}$ such that
\[
x^{n} + a_{1}x^{n-1} + \cdots + a_{n} = 0.
\]
\begin{itemize}
\item A complex number which is integral over $\mathbb{Z}$ is called an \emph{algebraic integer}.
\item If $x \in \mathbb{Q}$ is an algebraic integer, then $x \in \mathbb{Z}$: write $x = \frac{q}{p}$ where $(p, q) = 1$. Then
\[
p^{n} + a_{1}qp^{n-1} + \cdots + a_{n}q^{n} = 0.
\]
Thus $q \mid p^{n} \Longrightarrow q = 1$.
\item $x$ is integral over $\mathbb{Z}$ $\Leftrightarrow$ the subring $\mathbb{Z}[x]$ of $R$ generated by $x$ is finitely generated as a $\mathbb{Z}$-module. (Since $\mathbb{Z}$ is Noetherian.) $\Leftrightarrow$ There exists a finitely generated $\mathbb{Z}$-submodule of $R$ which contains $\mathbb{Z}[x]$.
\end{itemize}
\end{definition}

\begin{proposition}
If $C_{1}, \ldots, C_{r}$ are conjugacy classes of $G$ and $\lambda_{1}, \ldots, \lambda_{r}$ are algebraic integers in $\mathbb{C}$, then $\sum_{i=1}^{r} \lambda_{i}(SC_{i})$ is integral over $\mathbb{Z}$.
\end{proposition}

\begin{proof}
Note that each product $(SC_{i})(SC_{j})$ is a linear combination with integer coefficients of the $SC_{k}$. The subgroup $\mathbb{Z} SC_{1} \oplus \ldots \oplus \mathbb{Z} SC_{r}$ of $Z(\mathbb{C}G)$ is hence a subring. Since it is finitely generated over $\mathbb{Z}$, each of its elements is integral over $\mathbb{Z}$.

\end{proof}

\begin{lemma}

Let $\rho_{1}, \ldots, \rho_{r}$ be all simple representations (up to isomorphism). $\rho_{i}: G \rightarrow \operatorname{GL}_{d_{i}}(\mathbb{C}) \Rightarrow \rho_{i}^{\prime}: \mathbb{C} G \rightarrow M_{d_{i}}(\mathbb{C})$
$\mathbb{C} G \simeq M_{d_{1}}(\mathbb{C}) \oplus \cdots \oplus M_{d_{r}}(\mathbb{C}) \Rightarrow p_{i}^{\prime}: \mathbb{C} G \rightarrow M_{d_{i}}(\mathbb{C})$ projection.
Fix $i$, let $x \in Z(\mathbb{C}G)$, then $p_{i}(x)=\lambda I$ for some $\lambda \in \mathbb{C}$. In fact,

\[
p_{i}(x)=\frac{1}{d_{i}} \operatorname{tr}\left(p_{i}(x)\right) \cdot I
\]

Write $x=\sum_{g \in G} a_{g} g$, then

\[
p_{i}(x)=\frac{1}{d_{i}} \sum_{g \in G} a_{g} \chi_{i}(g) \cdot I .
\]

\end{lemma}

\begin{theorem}
The degree of a simple representation of $G$ over $\mathbb{C}$ divides $|G|$.

\end{theorem}

\begin{proof}
Let $x=\sum_{g \in G} \overline{\chi_{i}(g)} g \in Z(\mathbb{C}G)$ where $\overline{\chi_{i}(g)}$ is an algebraic integer.
Then $x$ is integral over $\mathbb{Z} \Rightarrow p_{i}(x)$ is integral over $\mathbb{Z}$

\[
p_{i}(x)=\frac{1}{d_{i}} \sum_{g \in G} \overline{\chi_{i}(g)} \chi_{i}(g) \cdot I=\frac{|G|}{d_{i}} I \Rightarrow \frac{|G|}{d_{i}} \text{ is integral over } \mathbb{Z} \Rightarrow d_{i} \mid |G|
\]

\end{proof}

\paragraph{Induction \& Restriction}

\begin{definition}[Restriction]

Let $W$ be a $KG$-module. $H \leq G$. Then $H \hookrightarrow G \rightarrow \operatorname{GL}(W)$
So $W$ is also a $kH$-module, denoted by $\operatorname{Res}_{H}^{G}(W)$: the restriction of $W$ to $H$.

\end{definition}

\begin{definition}[Induction]

Let $V$ be a $kH$-module.
Define a $kG$-module $\operatorname{Ind}_{H}^{G}(V)=kG \otimes_{kH} V$: induced from $H$ to $G$.

\end{definition}

\begin{remark}
$\operatorname{Res}_{H}^{G}: KG\text{-module} \to KH\text{-module}$ functors.
$\operatorname{Ind}_{H}^{G}: KH\text{-module} \to KG\text{-module}$
\end{remark}

\begin{proposition}
Characterization of $\operatorname{Ind}_{H}^{G}(V)$
Let $H \leq G$. Let $V$ be an $RH$-module ($R$ is a commutative ring with $1$).
$G/H = \{g_{1}H, \dots, g_{r}H\}$ where $r=[G:H]$ and $g_{1} \in H$.
Denote the set of representatives as $[G/H]$.
Then $\operatorname{Ind}_{H}^{G}(V)=\bigoplus_{i=1}^{r}(g_{i} \otimes V)$ as $R$-modules, where
\[
g_{i} \otimes V=\{g_{i} \otimes v \mid v \in V\} \simeq V \text{ as an } R\text{-module.}
\]
If $V$ is free, then $\operatorname{rank}_{R} \operatorname{Ind}_{H}^{G}(V)=[G:H] \operatorname{rank}_{R} V$.
If $x \in G$, then $x(g_{i} \otimes V)=g_{j} \otimes V$ where $x g_{i}=g_{j} h$ for some $h \in H$.
So the $R$-submodules $g_{i} \otimes V$ of $\operatorname{Ind}_{H}^{G}(V)$ are permuted transitively under the action of $G$, and $\operatorname{Stab}_{G}(g_{1} \otimes V)=H$.
\end{proposition}

\begin{proof}
$RG_{RH}=\bigoplus_{i=1}^{r} R(g_{i}H)=\bigoplus_{i=1}^{r} g_{i} RH \simeq (RH)^{r}$ as an $R$-module.

\[
RG \otimes_{RH} V=\bigoplus_{i=1}^{r} g_{i}(RH) \otimes_{RH} V=\bigoplus_{i=1}^{r}\left(g_{i} RH \otimes_{RH} V\right)=\bigoplus_{i=1}^{r}\left(g_{i} \otimes V\right) \text{ as } R\text{-mod.}
\]

If $x \in G$, $x(g_{i} \otimes V)=g_{j} \otimes hV$ (where $xg_{i}=g_{j}h, h \in H$),
i.e., $x(g_{i} \otimes V) \subseteq g_{j} \otimes V$. Similarly $x^{-1}(g_{j} \otimes V) \subseteq (g_{i} \otimes V)$, hence $x(g_{i} \otimes V)=g_{j} \otimes V$.
Thus $G$ permutes transitively on $g_{i} \otimes V$, $i=1, \dots, r$.
If $x \in H$, $x(g_{1} \otimes V)=xg_{1} \otimes V=g_{1}(g_{1}^{-1}xg_{1}) \otimes V=g_{1} \otimes V \Rightarrow \operatorname{Stab}_{G}(g_{1} \otimes V)=H$.
If $x \notin H$, $x g_{1}=g_{i} h$ for some $g_{i} \notin H, h \in H$.
\end{proof}

\begin{proposition}
Represented as an induced module.
Let $M$ be an $RG$-module that has an $R$-submodule $V$ such that $M$ is a direct sum of $R$-submodules $\{gV \mid g \in G\}$. Then $M \simeq \operatorname{Ind}_{H}^{G}(V)$ where $H=\{g \in G \mid gV=V\}$.
\end{proposition}

\begin{proof}
Define a map of $R$-modules: $RG \otimes_{RH} V \rightarrow M$, an $RG$-morphism
\[
g \otimes v \mapsto gv
\]
We have $\{gV \mid g \in G\} \longleftrightarrow G/H$
\[
gV \longleftrightarrow gH
\]
\end{proof}
$RG \otimes V$ and $M$ are direct sums of $[G:H]$ $R$-submodules $g \otimes V$, $gV$ respectively:
$RH \cong gV \cong V \cong gV$ as $R$-modules

\[
g \otimes V \leftrightarrow V \leftrightarrow gV
\]


\begin{proposition}
Character of induced module

$H \leqslant G$. Let $\chi$ be a character of $H$. Then

\[
\begin{aligned}
\operatorname{Ind}_{H}^{G}(\chi)(g) &= \frac{1}{|H|} \sum_{\substack{t \in G \\ t^{-1} g t \in H}} \chi\left(t^{-1} g t\right) \\
&= \sum_{\substack{t \in [G/H] \\ t^{-1} g t \in H}} \chi\left(t^{-1} g t\right)
\end{aligned}
\]

\end{proposition}

\begin{proof}
(1) Second ``$=$'': If $t^{-1} g t \in H$, $h \in H$, then $(th)^{-1} g(th) = h^{-1}(t^{-1} g t)h \in H$, hence $\chi\left(t^{-1} g t\right) = \chi\left((th)^{-1} g(th)\right)$.

(2) First ``$=$'': Recall $\operatorname{Ind}_{H}^{G}(V) = \bigoplus_{i=1}^{n} (g_{i} \otimes V)$.

Idea: $g$ acts on $\operatorname{Ind}_{H}^{G}(V)$:
\[
\begin{cases}
g g_{i} \otimes V = g_{i} \otimes V \Rightarrow \text{consider how } g \text{ acts on } g_{i} \otimes V \\
g g_{i} \otimes V \neq g_{i} \otimes V \Rightarrow \text{no contribution to the trace}
\end{cases}
\]

\[
g(t \otimes V) = t \otimes V \Leftrightarrow t^{-1} g t \otimes V = 1 \otimes V \Leftrightarrow t^{-1} g t \in \operatorname{Stab}_{G}(1 \otimes V) = H
\]

Note that $\chi$: $H$ acts on $V$. If $t^{-1} g t \in H$, then $g$ acts on $t \otimes V$ by

\[
g(t \otimes v) = gt \otimes v = t(t^{-1} g t) \otimes v = t \otimes t^{-1} g t\, v \Rightarrow g: t \otimes v \mapsto t \otimes t^{-1} g t\, v
\]

Hence trace of $g$ on $t \otimes V$ = trace of $t^{-1} g t$ on $V$ = $\chi\left(t^{-1} g t\right)$.

Recall:
\begin{itemize}
\item $V$: $A$-module, $W$: $B$-module, $T$: $A$-$B$-bimodule.
\[
\operatorname{Hom}_{A}(W \otimes_{B} T, V) \simeq \operatorname{Hom}_{B}\left(W, \operatorname{Hom}_{A}(T, V)\right)
\]
\item $A$, $B$ rings. $\exists\, \varphi: A \rightarrow B$ ring homomorphism, $1_{A} \mapsto 1_{B}$.
If $W$ is a $B$-module, then it can be viewed as an $A$-module through $a w := \varphi(a) w$, denoted by $W|_{A}$.
\end{itemize}

\end{proof}

\begin{lemma}
Let $A \rightarrow B$ be a ring homomorphism with $A \rightarrow B$. Let $V$ be an $A$-module, $W$ a $B$-module. Then
\begin{enumerate}
\item $\operatorname{Hom}_{B}(B \otimes_{A} V, W) \simeq \operatorname{Hom}_{A}(V, W \downarrow_{A})$
\item $\operatorname{Hom}_{A}(W \downarrow_{A}, V) \simeq \operatorname{Hom}_{B}(W, \operatorname{Hom}_{A}(B, V))$
\item If $B \rightarrow C$ is another ring homomorphism, then $C \otimes_{B} (B \otimes_{A} V) \simeq C \otimes_{A} V$
\end{enumerate}
\textbf{Pf:}
(1) \& (2) $\operatorname{Hom}_{B}(B \otimes_{A} B, W) \simeq W \downarrow_{A} \simeq B \otimes_{B} W$ as $A$-modules, $\omega \mapsto 1 \otimes \omega$.

\end{lemma}

\begin{corollary}
Let $H \leq K \leq G$. Let $V$ be an $RH$-module, $W$ an $RG$-module.
\begin{enumerate}
\item (Frobenius reciprocity)
\[
\operatorname{Hom}_{RG}(\operatorname{Ind}_{H}^{G} V, W) \simeq \operatorname{Hom}_{RH}(V, \operatorname{Res}_{H}^{G} W) \text{ as } R\text{-modules}
\]
\[
\operatorname{Hom}_{RG}(W, \operatorname{Ind}_{H}^{G} V) \simeq \operatorname{Hom}_{RH}(\operatorname{Res}_{H}^{G} W, V)
\]
\item $\operatorname{Ind}_{K}^{G}(\operatorname{Ind}_{H}^{K} V) \simeq \operatorname{Ind}_{H}^{G} V$ as $RG$-modules
\item $\operatorname{Res}_{H}^{K}(\operatorname{Res}_{K}^{G} W) \simeq \operatorname{Res}_{H}^{G}(W)$ as $RH$-modules
\item $\operatorname{Ind}_{H}^{G}(V) \otimes_{R} W \simeq \operatorname{Ind}_{H}^{G}(V \otimes_{R} \operatorname{Res}_{H}^{G}(W))$ as $RG$-modules
\item $\operatorname{Ind}_{H}^{G}(V) \simeq \operatorname{Hom}_{RH}(RG, V)$ as $RG$-modules
\end{enumerate}
\end{corollary}

\begin{proof}
(4) $(RG \otimes_{RH} V) \otimes_{R} W \rightarrow RG \otimes_{RH} (V \otimes_{R} W)$
\[
g \otimes v \otimes w \mapsto g \otimes v \otimes g^{-1} w
\]
\[
g \otimes v \otimes g w \leftarrow g \otimes v \otimes w
\]
(5) $RG \otimes_{RH} V \simeq \operatorname{Hom}_{RH}(RG, V)$
\[
g \otimes v \mapsto \phi_{g,v}: x \mapsto \begin{cases} (xg)v & \text{if } xg \in H \\ 0 & \text{otherwise} \end{cases}
\]
\[
\sum_{g \in [G/H]} g \otimes \theta(g^{-1})
\]
\end{proof}

\begin{corollary}
Frobenius reciprocity.

$H \leq K \leq G$. $\chi$: character of $H$. $\psi$: character of $G$. Note that
\begin{enumerate}
\item (Frobenius reciprocity) $\langle \chi_V, \chi_W \rangle = \dim_{\mathbb{C}} \operatorname{Hom}_{\mathbb{C}G}(V, W)$
\[ \langle \operatorname{Ind}_{H}^{G} \chi, \psi \rangle_{G} = \langle \chi, \operatorname{Res}_{H}^{G} \psi \rangle_{H}, \quad \langle \psi, \operatorname{Ind}_{H}^{G} \chi \rangle_{G} = \langle \operatorname{Res}_{H}^{G} \psi, \chi \rangle_{H}. \]
\item $\operatorname{Ind}_{K}^{G}(\operatorname{Ind}_{H}^{K}(\chi)) = \operatorname{Ind}_{H}^{G}(\chi)$
\item $\operatorname{Res}_{H}^{K}(\operatorname{Res}_{K}^{G}(\psi)) = \operatorname{Res}_{H}^{G}(\psi)$
\item $\operatorname{Ind}_{H}^{G}(\chi) \cdot \psi = \operatorname{Ind}_{H}^{G}(\chi \cdot \operatorname{Res}_{H}^{G}(\psi))$
\end{enumerate}
\end{corollary}

\paragraph{Mackey decomposition formula}

$H, K \leqslant G$. $V$ a $K$-module. $\operatorname{Res}_{H}^{G}(\operatorname{Ind}_{K}^{G}(V))$

If $\Omega$ is a left $G$-set, write $G \backslash \Omega$ for the set of $G$-orbits on $\Omega$ and $[G \backslash \Omega]$ for the representatives. Similarly for a right $G$-set $\Omega$, write $\Omega / G$, $[\Omega / G]$.

\begin{definition}
$(H, K)$-double coset
\[ HgK = \{hgk \mid h \in H, k \in K\} \quad \forall g \in G. \]
\end{definition}

\begin{lemma}
Properties of $(H, K)$-double cosets.

$H, K \leqslant G$.
\begin{enumerate}
\item Each $(H, K)$-double coset is a disjoint union of right $H$-cosets and left $K$-cosets.
\item Any two $(H, K)$-double cosets either coincide or are disjoint. The $(H, K)$-double cosets partition $G$.
\item The set of $(H, K)$-double cosets is in bijection with the orbits $H \backslash (G/K)$ under the mapping $HgK \mapsto H(gK) \in H \backslash (G/K)$.
\end{enumerate}
\end{lemma}

\begin{proof}
(2) If $h_{1} g_{1} k_{1} = h_{2} g_{2} k_{2} \in Hg_{1}K \cap Hg_{2}K$, then $g_{1} = h_{1}^{-1} h_{2} g_{2} k_{2} k_{1}^{-1} \in Hg_{2}K \Rightarrow Hg_{1}K \subset Hg_{2}K$.
\end{proof}

\begin{definition}
$H \backslash G / K$

We denote the set of $(H, K)$-double cosets in $G$ by $H \backslash G / K$, and the representatives by $[H \backslash G / K]$.
\end{definition}

\begin{lemma}
Let $H, K \leqslant G$, $g \in G$. Then
\begin{enumerate}
\item $H g K / K \simeq H / H \cap {}^{g}K$ as left $H$-sets
\[
{}^{g}K = g^{-1} K g
\]
\item $H \backslash H g K \simeq (H^{g} \cap K) \backslash K$ as right $K$-sets
\[
H^{g} = g H g^{-1}
\]
\end{enumerate}
So $[G: K]=\sum_{g \in[H \backslash G / K]}[H: H \cap {}^{g}K]$

\[
[G: H]=\sum_{g \in[H \backslash G / K]}[H^{g} \cap K]
\]

\end{lemma}

\begin{proof}
then $\Omega \simeq G / \operatorname{Stab}_{G}(w)$ for $w \in \Omega$.

$H g K$ is a single $H$-orbit on left $K$-cosets with $\operatorname{Stab}_{H}(g K)=\{h \in H \mid h g K=g K\}=\{h \in H \mid g^{-1} h g \in K\}=H \cap {}^{g}K$.
\end{proof}

\begin{definition}
${}^{g}V$ (a ${}^{g}H$-module)
\[
\begin{cases}
\text{as sets } {}^{g}V = V \\
\text{for } v \in V, h \in H, \underbrace{{}^{g}h \cdot v = h v}_{{}^{g}h \text{ on } {}^{g}V \text{ vs } h \text{ on } V}
\end{cases}
\]
i.e. If $H \xrightarrow{\varphi} GL(V)$ is the rep, then the rep of ${}^{g}V$ is given by
\[
\begin{aligned}
& {}^{g}H \rightarrow H \xrightarrow{\varphi} GL(V) \\
& {}^{g}h \mapsto h
\end{aligned}
\]
\[
g h (g \otimes v) = g h g^{-1} g \otimes v = g \otimes h v
\]
When $g \otimes V$ is identified with $V$ via $g \otimes v \leftrightarrow v$, the action of ${}^{g}H$ on $V$ coincides with the action of ${}^{g}H$ on ${}^{g}V$.
\end{definition}

\begin{theorem}
(Mackey decomposition formula)

$H, K \leq G$. $V$ a $K$-module. Then
\[
\operatorname{Res}_{H}^{G}\left(\operatorname{Ind}_{K}^{G}(V)\right) \simeq \bigoplus_{g \in[H \backslash G / K]} \operatorname{Ind}_{H \cap {}^{g}K}^{H} \operatorname{Res}_{H^{g} \cap K}^{K} V
\]
\end{theorem}

\begin{proof}
$\operatorname{Ind}_{K}^{G} V = \bigoplus_{x \in[G / K]} x \otimes V$.

For a double coset $H g K$, $\bigoplus_{x \in HgK / K} x \otimes V$ is an $H$-invariant subspace.
$x \in H g K$
\end{proof}

\paragraph{Lesson 15}

\begin{theorem}
$\operatorname{char} k = p$, $G$ a $p$-group. Then
\begin{enumerate}
\item the regular module $K G$ is an indecomposable projective $K G$-module that is the projective cover of the trivial module
\item every f.g. projective $k G$-module is free
\item the only idempotent elements in $KG$ are $0$ and $1$.
\end{enumerate}
\end{theorem}
% --- End rewritten content from 抽代3.md ---




\end{document}
